\documentclass[a4paper,11pt,openany]{ltjsbook}
\usepackage{luatexja}
\usepackage{luatexja-fontspec}
\usepackage{amsmath,amssymb,amsfonts,amsthm}
\usepackage{geometry}
\geometry{top=24mm,bottom=24mm,left=24mm,right=24mm,footskip=10mm}
\usepackage{graphicx}
\usepackage{tcolorbox}
\tcbuselibrary{skins,breakable}
\usepackage{hyperref}
\usepackage{enumitem}
\usepackage{booktabs}
\usepackage{fancyhdr}
\setlength{\headheight}{38pt}
\addtolength{\topmargin}{-20pt}

\definecolor{russellblue}{RGB}{18, 48, 92}
\definecolor{softblue}{RGB}{240, 246, 253}
\definecolor{darkslate}{RGB}{30, 42, 56}
\definecolor{deepcrimson}{RGB}{145, 28, 42}
\definecolor{warmgold}{RGB}{175, 120, 20}

\hypersetup{
    colorlinks=true,
    linkcolor=russellblue,
    citecolor=darkslate,
    urlcolor=deepcrimson,
    pdftitle={神秘主義と論理 —— 原著完全直訳・数理論理校閲版},
    pdfauthor={バートランド・ラッセル}
}

\pagestyle{fancy}
\fancyhf{}
\fancyhead[LE,RO]{\thepage}
\fancyhead[RE]{\leftmark}
\fancyhead[LO]{\rightmark}
\renewcommand{\headrulewidth}{0.4pt}

\title{
    {\huge\bfseries\gtfamily 神秘主義と論理、その他の試論}\\[12pt]
    {\Large\gtfamily —— 論理原子論・無限の数学・科学的認識論の全体係 ——}\\[18pt]
    {\large 原題：\textit{Mysticism and Logic and Other Essays} (1917)}\\[6pt]
    {\normalsize Project Gutenberg eBook \#25447}
}
\author{\textbf{バートランド・ラッセル} 著\\[4pt]教育情報分析研究会 訳・数理論理校閲}
\date{2026年10月 完全決定版}

\begin{document}

\maketitle

\frontmatter

\chapter*{訳者まえがき：20世紀数理哲学の金字塔}
\addcontentsline{toc}{chapter}{訳者まえがき：20世紀数理哲学の金字塔}

バートランド・ラッセル（Bertrand Russell, 1872--1970）が著した『神秘主義と論理、その他の試論（原題：\textit{Mysticism and Logic and Other Essays}）』は、ノーベル文学賞を受賞した稀代の知性であり、現代数理論理学・記号論理学・分析哲学の創始者であるラッセルの最も重要かつ美しく格調高い10篇の論文を収めた歴史的傑作集です（Project Gutenberg eBook \#25447）。

本書においてラッセルは、古代ギリシャのヘラクレイトスやプラトン以来の人類の知的探求を駆動してきた「神秘的衝動」と「科学的・論理的衝動」の対立と調和を鮮やかに論じる表題作をはじめ、冷厳な数学的真理の崇高な美を讃えた「数学の研究」、カントールの無限集合論やゼノンのパラドックスを明快に解決した現代数学の到達点を示す「数学と形而上学者たち」、そして現代物理学の時空と知覚経験を厳密な論理構築によって架橋する「センス・データと物理学との関係」「原因の概念について」「直接体験による知識と記述による知識」に至るまで、その思想の全貌を惜しみなく展開しています。

本完全翻訳・数理論理校閲版では、一切の要約や省略を排し、ラッセルの英語原文（序文および全10章）の重厚な文体、精密な論理の歩み、引用文献・脚注に至るまで忠実に日本語へ直訳いたしました。さらに、登場する集合論・実無限・関係論理・物理学的時空概念を厳密に校閲・検証し、ラッセルの知の建築を完全に追体験できる学術的決定版として組版・製本いたしました。

\clearpage
\tableofcontents
\clearpage

\mainmatter


\chapter*{序文（Preface）}
\addcontentsline{toc}{chapter}{序文（Preface）}


これらのエッセイは、様々な時期に執筆され、発表されました。再掲載の許可を得て発表することができたことに、以前の出版者への感謝を表します。

「神秘主義と論理」は1914年7月号の『ヒバート誌』に掲載されました。「自由な人間の崇拝」と「数学の研究」は、かつての集刊（現在は廃版）『哲学的エッセイ』にも収録されており、長蛇堂・グリーン社によって出版されました。これらはともに1902年に執筆され、「自由な人間の崇拝」は1903年版の『独立評論』に最初に掲載され、「数学の研究」は1907年11月号の『新 Quarterly』に掲載されました。理論的倫理学において、「自由な人間の崇拝」で主張されている立場は現在と完全には一致していません：当時よりも、善悪が客観的なものであるという確信が薄くなりました。しかし、そのエッセイで示唆される生活に対する一般的な態度は、神学的信念を持たない人々にとって、内面的な敗北を避けるために、ストレスや困難の時期に採用すべき姿勢だとまだ思っています。

「数学と形而上学者」に関するエッセイは1901年に執筆され、アメリカの雑誌『インターナショナル・ Monthly』に「最近の数学哲学における研究」というタイトルで掲載されました。このエッセイには、後の研究により修正が必要な点があります。これらの点はフ�注で指摘されています。その口調は、編集者が私に記事をできるだけロマンチックにすることを望んだという事実によって一部説明されます。

すべての以上記した論文は完全に一般向けですが、それ以降のものは多少技術的な内容が含まれています。「科学的メソッドに関する哲学」は1914年にオックスフォードで行われたヘンリー・スペナー講演であり、クラレンデンプレスから出版され、本書に収録することを好意的に許可されました。「物質の最終的な構成要素」は1915年初頭にマンチェスター哲学協会に対して行った演説で、同年7月に『モニスト』に発表されました。「感覚データと物理学の関係」は1914年1月に書かれ、同年の『サイエンティア』第4号（国際科学総合誌）に初めて発表されました。この雑誌はM・ユゼッジオ・リガーノ編集で、ロンドンのウィリアムズとノーグェート、ボローニャのニコラ・ザニチェッリ、パリのフェリックス・アルカンによって月刊で発行されました。「因果概念に関する論文」は1912年11月にアリストテレス協会に対して行われた議長演説であり、同年の『 Proceedings 』に出版されました。「認識と説明による知識」もアリストテレス協会に対して読まれた論文で、同協会の1910-11年の『 Proceedings 』に発表されました。

ロンドン
1917年9月

目次

第１章　神秘主義と論理                           1

第２章　自由教育における科学の位置               33

第３章　自由人の崇拝                          46

第４章　数学の研究                             58

第５章　数学と形而上学家                        74

第６章　哲学における科学的メソッド              97

第７章　物質の最終的な構成要素                  125

第８章　感覚データと物理学の関係                145

第９章　因果概念に関する論文                    180

第１０章　認識による知識と説明による知識          209

神秘主義と論理および他の論文


\chapter{第I章：神秘主義と論理 --- {MYSTICISM AND LOGIC}}

Mysticism と論理


形而上学、つまり思考を用いて世界全体を概念化しようとする試みは、神秘主義の衝動と科学の衝動という非常に異なる2つの人間の衝動が結びつき、対立する過程で発展してきました。一部の人々は、これらの衝動のいずれか一方だけを通じて偉大さを達成しましたが、他方では両方を通じて偉大さを達成しました。例えば、ヒumeにおいて科学的衝動は完全に制御されており、ブラックにおいては科学への強い反感と深い神秘主義的洞察が共存しています。しかし、哲学家として最も偉大な人々は、科学と神秘主義の両方を感じた必要がありました：これらの2つの要素を調和しようとする試みが彼らの人生を作り出し、常にその困難さと不確定性にもかかわらず、一部の人々にとって哲学が科学や宗教よりも大きなものであることを可能にしました。

科学的衝動と神秘主義的衝動の明確な定義を試みる前に、これら2つの哲学者の例を用いてそれらを説明します。彼らの偉大さは非常に密接に融合されたことによって達成されました。その2人の哲学者はヘラクレトスとプラトンです。

ヘラクレトスは、誰もが知るところですが、普遍的な流れの信者でした：時間はすべてのものを築き上げ、破壊します。わずかな断片から彼がどのように意見に至ったかを特定することは容易ではありませんが、いくつかの言葉は科学的観察がその源であることを強く示唆しています。

「見えるもの、聞こえるもの、学べるものは最も価値がある」と彼は言います。これは経験主義者の言葉で、彼にとって観察は唯一の真実の保証です。「太陽は毎日新しいものだ」また断片であり、その意見は反転的な性格を持っていますが、明らかに科学的反省によって刺激され、彼にとっては太陽が西から東へ地下を通り抜ける方法を理解する困難性を避けるためのものでした。実際の観察も彼の中心的な教義、つまり火が唯一の恒久的な物質であり、すべての見えるものはその一時的な段階であることを示唆しました。燃焼時に私たちは物事が全く異なる形で変化し、その炎と熱は空気へと上昇して消え去ります。

「すべての人々や神々にとって同じ世界」彼は言います。「誰もが作るのではなく；それは常に存在し、現在であり、永遠に燃える火である。量が燃え上がり、量が消えていく」ということです。

火の変化は、まず海である。そして海の半分は土で、もう半分は旋風である。

この理論は科学が受け入れないものとなったとしても、その精神は科学的である。プラトンが引用する有名な言葉も同様に科学的なものだと言える：「一度同じ川には二度と踏み込めない；新たな水が常にあなたを包むからである。」しかし、残存する断片の中にも別の声明がある：「踏み込みつつ踏み込んでいない。存在しつつ存在していない。」

この神秘的な声明とプラトンによって引用された科学的な声明を比較すると、ヘラク利スの体係において二つの傾向がいかに密接に関連しているかが明らかになる。神秘主義は、宇宙について信じるものを感情的に強化し深めるもので、そのような感情はヘラク利スに奇妙な鋭さを持った人生や世界に関する言葉を生み出す：「時間は子供が将棋を玩んでいるように、王権は子供のものである。」

これは詩的な想像力ではなく科学が、世界の独裁者として時間を見せる。それはまた、対極性の同一性を主張する神秘主義がヘラク利スに導く：「善と悪は一つだ」と彼は言うし、「神にはすべてが美しく正しく善であるが、人間にはそれが正しいとは思えないものもあるし、それらは不正であるとも思う。」

ヘラク利スの倫理学には多くの神秘主義が基盤となっている。科学的な決定論だけでも「人の性格は運命だ」という声明を生み出すことは可能だが、「獣たちは打撃と共に牧草地へと駆り出される」などという言葉は唯に神秘主義者のものである。「心の望みと戦うのは難しい。それは魂の代償で得られるものだ」また「知恵は一つである。すべてのものがその思想によって導かれていることを知ることだ」という言葉も同様である。

これらの例を増やすことはできるが、既に示されたものはヘラク利スの人間性を十分に示している：彼が見た科学的事実は彼の魂の中で炎のように燃え上がり、その光で世界の深淵を見通すことができた。このような性質において、神秘主義者と科学者の真の融合を見ることが可能である。私はそれが思想の世界における最高の地位であると考える。

プラトにおいても、同じ二重の動機が存在します。しかし、神秘的な動機はその二つの中でも明確に強いものであり、衝突が激しいときには最終的に勝利を収めます。彼が洞窟について描写したものは、感覚による知識と現実よりも真実でより本質的な知識と現実に関する信念の古典的な説明です：

「想像してみてほしい。地下の岩穴の中で長く閉じ込められている一群の人々がいる。その岩穴には光が入る入口があり、岩穴全体に沿って延びている。彼らは幼少期から足首と首を縛られており、頭を動かせないために常に前方向を見つめなければなりません。そして想像してみてほしい。火の光が彼らの背後にある上部に燃えているところがあり、その火と囚人との間に高さのある通路がある。その通路上には壁が築かれています。これは Conjurer が観客の前に設ける屏風のようなもので、その上空では奇跡が行われます。

彼は答えた。「そうですね。」

また、壁の向こう側に人々が歩いていて、人形や他の動物の像を運んでいると想像してみてほしい。これらの像は木や石など様々な材料で作られており、壁よりも高い位置にある。そして、当然のように、いくつかの人々は話し合い、他の方々は静かである。

奇妙な場面を描写しているね、奇妙な囚人たちだね。

そうですね、と答えた。「彼らは私たちに似ています。」

今、自然の流れが彼らから鎖を解き放ち、愚かさに対する治療法をとえると仮定してみよう。たとえば、一人が解放され、急に立ち上がり、首を回して目を開けて光に向かって歩くことになる。彼はこれらの行動を痛みの中で行い、眩 blinded 々の輝きにより、以前見ていた影の物体を見分けることができなくなる。もし誰かが彼に昔は愚かな幻影を見ておりましたが、今や少し現実に近づいたり、より現実的なものに向かい、より正確に見えるようになったと告げた場合、彼は何と言ったであろうか？また、彼に各々の物体を指し示し、質問し、答えを要求した場合、どう思うだろうか？

はい、その古い幻影が現在の物体よりもずっと真実であると考えるでしょう。

だからこそ、上位世界の物体を見分けるために習慣が必要となると仮定する。最初に彼は影を見分けることに最も成功するでしょう；次に水中での人物や他のものの反射を識別し、その後現実的なものを見るようになるでしょう。そして最後に、月や星の光に向かって目を上げ、夜には天体や宇宙自体を研究することが比較的容易であると感じるでしょう。しかし日中は太陽やその光を研究するよりも難しいでしょう。

もちろん。

最後に、私は彼が自分の領域で本来の太陽を見ることを観察し、考えることができるだろうと想像します。水の中や異なる土地ではなく、太陽そのものを見ることです。

当然のことです。

彼の次のステップは、太陽が季節と年を生み出し、視覚的世界における万物の保護者であり、また彼自身と同伴者が以前見ていたすべての物事を生成する原因であるという結論を導くことでしょう。

今この仮想的なケース、親愛なるグランコンさん、あなたは以前の声明全体に適用してください。視覚が明らかにする領域を監獄と比較し、その中にある火の光を太陽の力と比較します。そして、上昇と上界の観察によって魂が知性の領域へと昇るという理解があれば、私の推測の傾向に近づくでしょう。あなたはそれを知りたいと言っていますから。しかし、神のみがその正確性を知っているかもしれません。しかし、それはどうであれ、私はこの問題に対する見解は以下の通りです。知識の世界において、真実の善の本質的な形は私たちの探求の限界であり、ほとんど見えない存在です。しかし、それが見えるとすれば、我々は避けて通れない結論を下すでしょう：それはすべての明るさと美しさの源であり、視覚的世界では光とその支配者を生み出し、知性の世界では直接かつ完全な権限で真実と理由を配分します。そして、誰もが賢明に行動するためには、この善の形を眼前に置くべきです。

しかし、この文脈においてプラトンの教えを通じて大半は、善と真正な現実との同一性が哲学的传統に体現され、今日でも大きく影響力を持っています。善に対する立法的機能をとえたことで、プラトンは哲学と科学の間の離婚を生み出し、その結果、両者はそれ以来もまた現在もその影響を受け続けています。科学者にとって何らかの希望があろうとも、彼は自然を研究するときにはそれを捨てなければなりません。そして、真実を得るためには哲学者は同じことをしなければなりません。倫理的な考慮は真実が確定したときにのみ正当に現れるべきです：それは我々の真実に対する態度や、真実に基づいた生活秩序を決定するものとして現れることができるが、それ自体が真実とは何かを指示することはできません。

ソクラテスが善の概念があることを説明した後、パルメニデスは彼に「最も卑しいものすら軽視してはならない」と助言し、この助言は真の科学的精神を示しています。哲学が最大の可能性を実現するためには、神秘主義者の高い現実への洞察と隠された善に対する理解を不偏不党な態度で組み合わせる必要があります。そして、この点での失敗が理想主義的哲学を薄らかで生き生きとしたものではなく、虚無的なものにしたことは多くの人々が知っています。理想的なものは世界との結婚によって果実を結ぶことができますが、それが世界から離れたままだと不稔のままです。しかし、理想的なものは事実から後退したり、世界がその欲求に従うことを事前に要求するものでは、その結婚は達成されません。

パルメニデス自身は、プラトンの思想に特有の興味深い神秘主義の流れを生み出し、この神秘主義は論理上の理論を通じて表現されます。この形式の神秘主義は、西欧においてパルメニデスによって起源したとされ、ヘーゲルやその現代の弟子たちまでの偉大な神秘的形而上学者の議論を支配しています。彼は実在が創造されず、破壊されず、変化せず、分割されないと言います。「巨链の束縛の中で動かず、始まりもなく終わりもない；生成と消滅は遠ざけられ、真なる信仰によって放棄された」と彼は述べます。彼の探求の基本原理は、ヘーゲルの文章にも似た一文で述べられています：「存在しないものを知ることは不可能であり、それを表現することもできない；それは考えられるものと話すことのできるものが同じであるからだ」そしてまた、「考えられかつ言えるものは存在する；それは存在することが可能であり、何もないものが存在することは不可能だからだ」と彼は言います。変化の不可能性はこの原理から導かれます；過去に述べることができるものは現在も依然として存在しているのです。

世界中のどの時代でも、神秘主義哲学は特定の信念で特徴付けられ、我々が考慮した教義によって示されています。

まず、洞察に対する信頼が論理的な分析的知識に対比される信念がある：突然の、深遠な、強制的な智慧への道に対する信頼である。これは、外見を研究する科学に完全に依拠する感覚に基づく遅延且つ不確かな学問と対照されるものだ。内面的興奮に吸収できる者は時として、日常の物事から離れた奇妙な現実感の欠如を感じ、外なる世界の堅固さが失われ、魂は絶望的な孤独の中で自身の深淵から狂った幻想の舞踏を生み出すと感じる。これは神秘主義者の入門の否定的側面であり、共通知識に対する疑念はより高い智慧を受け入れる準備を整える。多くの人々にとってこの否定的な経験が馴染みあるものでも、神秘主義者はそれ以上進まないが、神秘主義者にとってはそれは広大な世界への門戸に過ぎない。

神秘主義者の洞察は、謎が解き明かされ、突然確実性を獲得し、疑いの余地がない隠された智慧があるという感覚から始まる。確信と啓示の感覚は具体的な信念よりも早く現れる。神秘主義者が到達する具体的な信念は、洞察の瞬間に得られた未表現された経験に基づく反芻の結果である。しばしば、この瞬間とは実質的な関連がない信念も中心核に引き寄せられる。したがって、すべての神秘主義者共通の信条に加えて、多くの人々にはより局所的で一時的な信条が見られ、これらは主観的な確信により本質的に神秘的なものと混ざり合う。このような非本質的な付着物を無視し、すべての神秘主義者が共有する信念のみに焦点を当てるべきである。


第一の結果、そして最も直接的なものは、知識の道が存在する可能性に対する信仰です。この知識の道は、感覚、論理、解析と対照的に、盲目の導き手として見なされる幻影の泥沼へと導くものではないと信じられています。この信仰は、現象の世界の背後にある実在を想起させますし、その実在は現象から全く異なるものであると考えられています。この実在はしばしば崇拝に近い敬意で見られ、常にそしてどこでも近くに存在すると感じられます。感覚の演出の薄い被り物の向こう側に隠されていると信じられ、受容的な心にはその輝きが現れる準備ができています。詩人、芸術家、恋人はその輝きを求めています：彼らが追求する haunting 美はその太陽の微かな反映です。しかし神秘主義者はその視覚の完全な光の中で生活しています：他人が曖昧に求めているものは彼には知られています。それ以外のすべての知識とは比べ物にならないほど深い知識で。

神秘主義の第二の特性は、統一性への信条と、どこでも対立や分割を認めない拒否です。ヘラクレイトスは「善悪は一つ」と述べ、「上り道と下り道は同じである」とも言っています。この態度は、矛盾する命題の同時主張にも現れます：「我々は同じ川にかつまたかつまたせず、存在しかつまた存在しない」などです。パルメニデスの主張、「実在は一つで不可分である」というものも、同一性への傾倒から来ています。プラトンではこの傾倒はアイデア論によって制約されますが、彼の論理が許す限り、善の優先性という教義に再現されます。

ほぼ全ての神秘主義的メタ物理学には第三の特徴があります：時間の実在を否定する拒否です。これは分割を否定することから生じます：すべてが一つであるなら、過去と未来の区別は幻覚でなければなりません。パルメニデスではこの教義が目立っていますし、現代においてはスパノジアやヘ겔の体係において基礎となっています。


最後に考慮すべき神秘主義の教義は、すべての悪は単なる表象であり、分析的な知性による分離と対立によって生じた幻覚であるという信念です。神秘主義は、例えば残酷さのようなものは善ではなく、しかし実在ではないことを否定します：それは我々が洞察によって解放されるべき LOWER WORLD OF PHANTOMS の一部に属しています。

しばしば、ヘーゲルにおいては特に、そしてスピノザでは少なくとも言葉レベルで、悪だけでなく善も虚構と見なされます。ただし、現実として認識されるものが持つ感情的な態度は、現実が善であるという信念と自然に結びつくものがあります。すべての場合において、神秘主義の倫理的な特徴は、憤慨や反論の欠如であり、喜びとともに受け入れ、二つの敵対する陣営、善と悪への分裂が最終的な真実ではないという不信です。この態度は神秘体験の性質から直接生じます：統一感があると同時に無限の平和を感じる感情があります。

実際には、平和感は夢における感情と同じように、関連する信念の全体システムを生み出すかもしれません。しかしこれは難しい問題であり、人類が一致見解に達する可能性は低いでしょう。

神秘主義の真偽を考える際に四つの質問が浮上します：

  I. 分析と直感という二つの認識方法があるでしょうか？そして、どちらが他方よりも優れているでしょうか？

  II. 全ての多様性と分離は虚構でしょうか？

  III. 時間は実在ではないでしょうか？

  IV. 善悪にはどのような種類の実在性がありますか？

すべての四つの問いにおいて、完全に発達した神秘主義は私には誤りのように思われますが、それでも十分な自制によって、神秘的な感情から学べる智恵の要素があると信じています。これ以外の方法では得られないようなものでもあります。これが真実であれば、神秘主義は生活に対する態度としては称賛されるべきであり、世界についての信条としてではなくです。私は、その情熱が誤りの結果であると主張しますが、それでもこの情熱は他のすべての思考や感情を色付けし、形づくることによって、人間の最も良い部分の Inspirer となっています。すなわち科学によって真実を慎重で忍耐強く調査することすら、神秘主義のその尊厳の精神によって育てられ、養われることもあります。

I. 理論と直感

神秘の世界についての実在性や非実在性については何も知らない。それを否定する意図も、その洞察が真実の洞察であることを断定する意図もない。しかし、科学的な態度が求められるのは、未検証かつ支持されていない洞察が真実を十分に保証することではないという主張だ——そしてこれが問題の核心である。多くの重要な真実は、その手段によって最初に示唆されるが、未検証で支持されていない洞察だけでは真実性を担保することは不十分だと指摘する。本能と理性との間の対立はしばしば話題になるが、18世紀には理性が優遇された一方で、ルソーとロマン主義運動の影響下では本能が優遇されるようになった。政府や思想の人工的な形式に反発した人々によって最初に、そして伝統的神学に対する純粋な理性的防衛がますます困難になった後には、科学が彼らが生命や世界に対する精神的な見方と結びつける信条に対して脅威となると考えた人々によって、本能が優遇されるようになった。ベルソンは「直観」という名前で、本能を形而上学の真実唯一の仲裁者にまで引き上げた。しかし実際には、本能と理性との間の対立は主に虚構である。本能、直観、または洞察が最初に信念を導き出し、その後理性能が確認したり否認したりするが、確認がある場合でも最終的には他のより本能的な信念との一致によって成立している。理性能は調和と制御の力であり、創造的ではない。最も純粋な論理の領域においても、新しさを最初に到達するのは洞察である。

本能と理性が時として衝突する場面は、本能によって持たれる単一の信念においてであり、その信念は本能的に持たれ、他の信念とはいかなる程度の不一致にもかかわらず捨て去ることなく強固に保持される。本能は人間のすべての能力と同じく誤りを犯す可能性がある。理性が弱い人はしばしば自己に関してこれを受け入れるのを拒否するが、他人に関しては誰もがこれを認めている。本能が最も誤りを犯しにくいのは実用的な事柄においてであり、正しい判断が生存に役立つ場合である：他の人々に対する友情と敵意は、非常に精密な隠蔽を通じて驚くほど正確に感じられることがある。しかし、このような事柄でも、慎重さや奉承によって誤った印象をとえることがあります；そしてより直接的な実用性のない事柄、例えば哲学が扱うようなものにおいて、非常に強い本能的信念はしばしば完全に間違っていることがあり、他の同程度に強い信念との観察された不一致を通じてそれが明らかになることがあります。このような考慮は、信頼性を互いの相対的な整合性でテストし、疑惑のある場合には一方と他方の誤りの可能性的な源を検討する理性による調和的調整が要求します。理性は本能全体に対する反対ではなく、ある興味深い側面への盲目な依存から他のより一般的だがそれだけ信頼できる側面を排すものにのみ反対です。このような偏りこそが、理性が矯正しようとするものです。

これらの多少陳腐な格言は、ベルソンが「直感」を「知性」と対比して主張する場合に例示されることが可能です。彼は、「一つの物事を理解するには二つの根本的に異なる方法がある」と言います。「第一の方法では物体を取り巻く移動が必要である；第二の方法ではその内部に入ることである。第一の方法では視点と表現するためのシンボルに依存しているが、第二の方法では視点にもシンボルにも依拠していない。第一の知識は相対的なものに留まるかもしれませんが、第二の知識はそれが可能であれば絶対的なものを達成します」と彼は述べています。彼が言う二番目のものは直感であり、「個々の物体の中に自分自身を置き、その中で唯一性を持つものを一致するための知的共感」であると彼は言います（p. 6）。例として、彼は自己認識を挙げます：「私たちは少なくとも一つの現実を直感によって内側から捉えることができますが、単なる分析ではなく。それは時間を通じて流れている私たち自身の人格であり、それが存在します—私たち自身が持続するもの」（p. 8）。ベルソンの哲学の残りは、言葉という不完全な媒介を介して直感によって得られた知識を報告し、そして科学や常識から主張されるすべての知識に対する完全な批判です。

この手続きは本能的な信念の衝突に立ち入るため、一方の信念よりも他方の信念の信頼性がより大きいことを証明することで正当化される必要があります。ベルソンはこれを二つの方法で試みます：まず知性は生物学的成功を確保するための純粋な実用的な能力であると説明し、次に動物における驚くべき本能の活用例を挙げ、彼が解釈する知性には直感が理解できるものの、困惑させる世界の特性も指摘します。

ベルソンの理論において、知性は生存競争における実用的な能力であり、真の信念の源ではないという点について、まず言えることは、我々が生存競争や人類の生物学的係譜を認識する手段は知性のみであるということです。もし知性が誤りを含むとすれば、この推測に基づく歴史全体はおそらく偽りとなるでしょう。一方で、進化がダーウィンの考えに従って進行したと同意するなら、知性だけでなくすべての能力は実用的な利便性の下で発達したものであると言えます。

直感は直接的に有用な場面で最もよく現れるため、例えば他人の性格や態度についてです。ベルソンはおそらく、このような知識の能力が生存競争によって説明されるよりも、純粋な数学の能力の方がより明白だと考えています。しかし、偽の友情に欺かれた野蛮人はその間違いを生命の代償で支払うでしょう；一方で、最も文明化された社会においても、数学的な無能さは人を死刑に処すことはありません。

彼が示した直感の動物的例の大半は非常に直接的な生存価値があります。実際、知性と直感は有用であるため発達し、大まかには真実をとえるときには有用であり、偽りをとえるときには害になると言えます。文明化された人間においては、知性は個々人の利便性を超えて発達することがあります；一方で、直感は全体として文明化が進むにつれて減少する傾向があります。通常、子供のほうが大人よりも、教育を受けた者よりも直感力が高いです。おそらく犬には人間よりも高い直感力があるでしょう。

しかし、これらの事実を直感の推奨と見る人々は、木々の中での生活に戻り、ウードで染められ、ヒップやハーブに依存するべきでしょう。

次に、ベルソンが主張するように直感には不思議な確実性があるかどうかを検討しましょう。彼によれば、その最良の例は自分自身に対する認識だとしますが、自己認識は伝説のようにまれで難しいものです。多くの人々にとって、自分の本質の中に卑劣さ、虚栄心、妬みなどがあり、それらは彼ら自身では全く意識していませんが、甚至最も親しい友人でもそれらを容易に察知することができます。直感には論理的な確信がある一方で、それが存在する限りその真実性を疑うことはほとんど不可能です。しかし、直感が検討の結果、少なくとも論理と同程度に不確実であることが明らかになった場合、より強い主観的確信は欠点となり、ますます誘惑的な欺瞞を生むことになります。

自己認識以外で最も目立つ直感の例は、人々が愛人に対して持っている自信を持つ知識です。異なる人格間の壁が透明になり、人々は他者の魂を見通すかのように思えます。しかし、このような状況での欺瞞は常に成功して行われており、故意的な欺瞞がない場合でも、経験は通常、その想定された洞察が虚構であり、長期的には論理の遅い探索的手法の方がより信頼性が高いことを証明します。

ベルソンは、知性は過去に経験されたものに類似しているものを対象とし、しかるべく直感は各新鮮な瞬間に常に存在する独特さと新しさを把握する力を持っていると主張します。各瞬間には確かに何か独特で新しいものが存在することは真実であり、これを見出すためには知性的概念では完全に表現できないこともまた真実です。独特で新しいものを知るためには直接的な経験が必要です。

しかしながら、そのような直接的な経験は感覚によってとえられ、特別な直感の能力を必要としないことが見受けられます。新データを提供するのは知性でも直感でもなく、感覚です；しかし、これらのデータが特に驚くべき方法で新しいものである場合、知性は直感よりもそれらに対処するのに優れています。

鶏が雛たちと共にいるときには、彼女自身が雛たちの中に存在しているような直感があると見なされるかもしれませんが、ただ分析的にそれを認識しているだけです。しかし、雛たちは水に飛び込むと、その全体的な直感は虚構であることが明らかになり、鶏は岸辺で手足を束ねてしまいます。

実際には、直感は本能の側面と発展形であり、他のすべての本能と同じく、特定の環境が動物の習慣を形成した場合に優れていますが、その環境が非慣習的な行動を要求する方法で変化すると、全く無力になります。

世界的な理論的理解は、哲学の目的であるが、その実用的な重要性は動物や原始人、さらには多くの文明人の多くにとって大きなものではない。従って、本能や直感の迅速で粗略な方法がこの分野で有利な土壌を見つけることはほとんど考えにくい。我々と遠い祖先との関連を引き出す古い種類の活動は、直感が最もよく現れる場面である。自己保存や愛のような事柄では、直感は時として（ただし必ずしもそうではない）批判的な知性に驚くべき速さと正確さで作用する。しかし哲学は、過去との親和性を示すような追求ではなく、高度に洗練され、文明化された追求であり、成功のために本能の生活から一定の解放が必要であり、時にはすべての俗世的希望や恐怖からも一定の距離を持つ必要がある。従って、直感が最もよく現れる場所は哲学ではない。むしろ、真の対象とその理解に必要な思考習慣が奇妙で不通常かつ遠くにあるため、ここではより-than どこよりも知性が直感を上回り、素早く分析せずに形成された信念は批判的接受には最も価値がない。

科学的な自制心と均衡を強調し、自信に基づく直感の主張に反対することは、知識の領域においてすべての大宗教によって説かれているような広大な観察、個人的で利己性のない態度、実用的な関心からの解放を推奨することである。従って、我々の結論は、多くの神秘主義者の明示的な信念とは対立するかもしれませんが、その精神に反しているわけではなく、この精神が思考の領域で適用された結果である。


最強力な神秘的洞察の側面の一つは、万物一体性の明白な示唆であり、これにより宗教ではパンテイズム、哲学では単一論が生じます。パルメニデスからヘーゲルとその追随者までを含む精密な論理が徐々に発展し、宇宙は一つの不分割なる全体であることを証明しようと試みています。現象世界の部分として考えられるものさえも、本質的で自己存続的な存在と見做されれば、単なる幻覚に過ぎないとされます。パルメニデスによって西洋哲学には、現象世界とは異なる実在概念が導入されました。少なくとも名目的には神秘的または宗教的理由ではなく、非存在の不可能性に関する論理的議論に基づいています。その後の多くの形而上学体係はこの基本的な考えから派生しています。

神秘主義を支持する論理は論理としては欠陥があり、技術的な批判を受け得る点があります。私はそれらの批判については他處で説明しましたので、ここでは繰り返すことはしません。代わりに、神秘的論理が発生した心持ちを分析しようと試みます。

信念は感覚に現れるものとは全く異なる実在を信じることから、不可避的な力で生じます。この信念が生まれるような気分状態は、多くの神秘主義と形而上学の源となっています。そのような気分状態が支配的であるときには、論理の必要性を感じません。そのため、より徹底的な神秘主義者は論理を使わず、直接その洞察による直感を訴えます。しかし、完全に発達した神秘主義は西洋では稀です。感情的確信の強度が下火になると、論理的に考え慣れた人は自己内に見出された信念に対する論理的な根拠を探求します。しかし、既存の信念があるため、彼は自ずと提示されるどんな根拠でも好んで受け入れます。彼の論理によって示されたパラドックスは実際には神秘主義のパラドックスであり、彼が洞察に合致するように論理が到達すべき目標です。その結果生じる論理は、多くの哲学者を科学や日常生活の世界について説明する能力から無力化させました。もし彼らがそのような説明を欲しければ、論理の誤りを見つけることができたでしょう。しかし、ほとんどの人は科学と日常の世界を理解することよりも、超感覚的な「実在」の世界の虚偽性を証明するためにそれほど熱心ではありませんでした。

このように、神秘主義者である偉大な哲学者たちが論理を追求しました。しかし彼らは通常、神秘的気分による仮定された洞察を前提とし、その論理的教義は一定の乾いた形で提示され、弟子たちはそれらが突然の啓示から生じたものとは全く独立していると信じていました。しかし、その起源は彼らに付き纏い続け、科学や常識の世界に対する「悪意」を保持していました。これが哲学者たちが、最も確立されていることや最も信頼すべき事実と矛盾する教義にも満足して受け入れる理由です。


論理の神秘主義は、当然ながら、何かが悪質であることに内在する欠陥を示します。神秘的な気分が支配的である間には論理への衝動を感じませんが、その気分が薄れると、論理への衝動は再び現れます。しかし、消えゆく洞察を保持しようとしたり、少なくともそれが洞察だったことを証明しようとしたりする望みと共にです。このような論理は完全に無関心や正直ではないもので、日常の世界に対する一種の憎しみから Inspiration を受けています。このような態度は当然ながら最良の結果には至りません。誰もが、作者を単純に反駁するために読むのは彼を理解する方法ではないことを知っています；そして、自然の書物をすべて幻影であるという確信を持って読むのも、理解につながる可能性は低いです。論理が共通の世界を理解可能なものにするためには、敵対的ではなく、通常は形而上学家中では見られない真摯な受け入れに基づいていなければなりません。

III. 時間

時間の虚無性は多くの形而上学的な体係における枢要な教義であり、パルメニデス以来しばしば論理的議論に基づくと名指されるが、その原点は新体係の創始者たちが神秘的な洞察の瞬間に生まれる確信から来ています。ペルシャの sufī の詩人によれば：

「過去と未来は神を私たちの視野から覆うものである。
両方とも火で燃え尽きさせよ！
これらの区切り線に分割されるのはどれほど続くだろうか、稊のように？」

最終的には実在するものは不変であるべきであるという信念は非常に一般的です：これは形而上学的な存在の概念を生み出し、現在でもエネルギーと質量の保存則のような科学的教義から正当化された満足を見出すことができます。

この見解における真実と誤りを分離するのは難しい。
時間は虚しいものであり、感覚の世界は幻であるという主張に対する論拠は、私はその大部分が誤ったものだと考えるべきである。しかし、時間は現実にとって重要な面ではなく、表面的な特性であるという意味では、何らかの意味で正しい感覚があることは認められるべきだ。過去と未来を現在と同じように実在として認めることが必要であり、時間への奴隷から解放されることが哲学的な思考には不可欠である。時間の重要性は主に実用的ではなく理論的であり、私たちの欲望に関連してではなく真実に関連してこそである。より正しく世界を描くためには、永遠なる外部の世界から物事が時間の流れに流入するという視点の方が、時間をすべてを飲み込む暴君として見るよりも良いだろう。思想と感情においても、時間が実在であっても、時間の重要性を理解することは賢明さの門である。

この事実はすぐに自問自答することで確認できる。なぜ過去に対する感情と未来に対する感情が異なるのかを考える。その違いは完全に実用的な理由によるものだ：私たちの願いは未来に影響をとえることができるが、過去には影響をとえることができない。未来は多少なりとも私たちの力によって支配される可能性がある一方で、過去は不変である。しかし、すべての未来はいずれ過去となる：もし現在過去を正しく理解することができれば、その未来はまだ未来だったときも同じように見えていたし、今未来とされているものは、それが過去となったときにどのように見えるかを今知っているように感じられる。したがって、過去と未来の感受的な質的な違いは内在的なものではなく、私たちに対する関係性によるものである：公正な観察からはその違いは消える。そして、公正な観察は知識の領域において無関心性という徳であり、行動の領域では正義や非自発性と現れるものだ。世界を正しく見ること、実用的な欲望から思考で立ち上がることを望む者は、過去と未来に対する態度の違いを克服し、時間の流れ全体を一視野に収める視点を持つことを学ぶべきである。

時間についての理論的哲学思考において、私は時間が介入すべきでないべき方法は、進化というアイデアと関連付けられた哲学を通じて例示される。この哲学は、最低な生命形態から人間までの発展を背景に、宇宙における基本法則として\textit{進行}を見出し、従って\textit{以前}と\textit{以後}の差異を自己反省的な見解の要塞自体に入れ込むことを認める。その世界の過去と未来についての推測的歴史については、私は争うつもりはない。しかし、真の宇宙理解のために必要なものが、成功の迅速な喜びに酔いしれることで忘れ去られてしまったと考える。ギリシャ精神と東洋的な諦めをその急進的な西洋的自己主張と共に組み合わせて、若さの熱意から成熟した人間の睿智へと発展することができるだろう。科学への呼びかけはあっても、真の科学哲学はより困難で、俗世的な希望よりも少ないものに訴えをかけるものであり、成功的な実践には厳しい訓練が必要である。

ダーウィンの『種の起源』は世界に告げたように、動物や植物の異なる種類間の差異は表面的に見られる固定且つ不変なものではないことを示した。自然種の教義は、分類を容易かつ明確にするものであり、アリストテレス伝統において奉守され、正統な神学の必要性によって保護されてきたが、突然生物界から永遠に消え去られた。人間と低等動物との差異は、私たちの人間的な傲慢さによく見られるように巨大なものであるが、中間的な存在が確実に人間家族内または外に位置付けられることのできない段階的な達成であることが示された。ラプラスによって太陽や惑星も既に原始的な未分化なネブulaからおそらく派生したものであることが示され、したがって古い固定された地平線は揺らぎ、すべての鮮明な輪郭はぼんやりとし始めた。物質と種類は境界を失い、誰もがその始まりや終わりを定めることができなくなった。

人間の傲慢が猿との関連性によって一時的に揺らぐことはあったが、そのような状況はすぐに「進化論」の哲学によって克服された。アモーブから人間への進化過程は、哲学者たちには明らかに前進であると見なされた。アモーブ自身がこの意見に賛成するかどうかは不明だが、科学によって示された過去の変遷のサイクルは、宇宙における発展の法則を明らかにし、善への進化またはアイデアの徐々に実現される過程として歓迎された。

しかし、このような見解はスペンサーやヘゲロ派進化論者を満足させるかもしれないが、変革のより真摯な支持者はこれを受け入れることはできない。これらの心には、世界が無限に近づく理想は死と静止し、激励的ではない。理想だけでなく、望みも進化の過程とともに変化し発展するべきであり、固定した目標ではなく、生命という動機によって引き起こされる新しい需要の継続的な形成が必要である。

この哲学における生命は一貫した流れであり、すべての分断は人工的で現実的ではない。個々のもの、開始と終結は便利な架空のものに過ぎず、ただ滑らかな連続性があるだけである。今日の信念が流れる途中を導く限り、その日は真実であり、しかし明日には偽りとなり、新しい状況に対応する新たな信念で置き換える必要がある。私たちのすべての思考は便利な架空のものであり、流れの中止点を想像的に凝らせたものに過ぎず、現実は我々の架空とは逆方向に流れ続け、思慮深く考えることは可能でも、それを思い描くことはできない。未来は予測不能でも、過去や現在よりも良いと確信され、明示的な声明なくその保証が浸透しているように感じる。未来を予見することはできなくても、少なくとも我々にはそれがより良いものであるという保証がある。

論理学、数学、物理学はこの哲学において消え去る。これらは「静止」しすぎているからである。現実とは、進歩の目標に向かう動機と運動であり、その目標は我々が近づくにつれて遠ざかる虹のように後退し、到達した場所が遠くから見えたものとは異なるからである。

哲学のこの思想について技術的な検討に入ることを提案するつもりはありません。私はただ、その動機と利益がいかに実用的で限定的であるか、そして扱う問題がいかに特殊なものであるかを主張したいだけです。それらは真の哲学として考慮に入れられるべき本質的な問いとはほとんど関係がないと考えられます。

進化論の主要な興味は人類の運命、または少なくとも生命の運命に関する問題です。それは道徳性と幸福よりも知識そのものに対する興味に更多ではありません。しかし、多くの他の哲学でも同様であり得ることを認めなければなりませんし、哲学がとえることができるような種類の知識への欲求は非常に稀であることも認めなければなりません。しかし、哲学が真実に到達するためには、まず第一に哲学者たちは真の科学者と同様の無私的な知的好奇心を持つことが必要です。

未来に関する知識—それが人類の運命について知るためには求められるべき種類の知識—は特定の狭い範囲内において可能であり、その範囲が科学の進歩とともにどの程度拡大されるかを予測することは不可能です。しかし明らかなことは、未来に関する任何の主張はその主題によって特定の科学の一部に属し、その科学の方法でしか確認されないことです。哲学は他の科学と同様の結果を得るための短絡的な道ではありません：真実の学問として存在するためには、独自の領域を持ち、他の科学が証明できないことも否認できないような結果を目指す必要があります。

進化論は、変化が悪いものから良いものへと進行することを前提とする「進歩」の概念に基づいており、そのような観点からは、時間という概念が進化論の主人となるよりも奴隷となるべき存在であるように思われ、哲学的な考えや感情の中での公正な見地を失うことがあります。形而上学者はしばしば、時間の実在性を完全に否定してきました。私はそのようなことはやりたいわけではありません；私が望むのは、過去が現在と未来と同じ程度の現実性と重要性を持つという否定から得られた精神的な視点を保存することです。「知性が理由に基づいて何かを考えるとき」という点においては、「それは未来、過去、または現在のものに関するアイデアでも同じように影響を受けます」（スパルザナ）。この「理由に基づいて考えること」が、進化に基づく哲学に欠けていると私は思います。

IV. 善悪

神秘主義はすべての悪を幻視であると主張し、時には善についても同じ見方をしますが、しばしば全ての実在は善であるという見方を持ちます。これらの両方の立場はヘラク利トスにおいて見られます。「善と悪は一つだ」と彼は言いますが、「しかし神にはすべてが美しく、善く、正しくある一方で、人間には何かを誤りとして何かを正しいものとする」（ヘラク利トス）。似たような二つの立場はスパルザナにおいても見られます。彼は「完璧」という言葉を使用して、「単なる人間の善ではなく、実在そのものに属する善」について話すときです。「現実と完璧とは同じことを意味します」と彼は言います。しかし他方で、「私が善を指すものは、私たちが確かに有用であると確信できるものを指す」と定義しています。したがって、完璧性は実在そのものに固有のものですが、善さは私自身や私たちの需要に対して相対的であり、公正な見地からは消えます。このような区別が必要であると考えています。それは現世的な善悪と呼ばれる下位の種類があり、表層世界を対立するように見える部分に分けるものですが、それとは異なる高位の神秘的な善が実在に属し、そのような相補的な悪に対抗することはありません。

この立場を論理的に説明するには、善悪は主観的であることを認めなければならず、善とは特定の感情を持つ方向性であり、悪とは別の種類の感情を持つ方向性であるだけです。行動的な生活において選択を下し、二つの可能な行為の中から一つを選ぶ際には、善と悪、または少なくともより良いものと悪いものの区別が必要です。しかし、この区別は、行動に関連するすべてのことが迷幻的世界に属すると神秘主義は考えるように、時間に関わる性質があるため、実在世界の幻覚に属します。反対に、行動が求められない静観的な生活では、公正な態度を保て、行動が必要とする倫理的二元論を克服できます。我々が単に公正である限りは、行動の善と悪はすべて幻想であると言えるでしょう。しかし、神秘主義の視野を得るために我々はしなければならないように、世界全体を愛し崇拝すべきものとして見ることで、より高い善があることを言えます。そしてこのより高い善は実在の世界に属します。

この二つの態度と神秘主義の表面上の揺れが説明され正当化される方法は、すべて存在するものを普遍的な愛と喜びで満たす可能性が生活の進行と幸せにとって極めて重要であり、その神秘的感動には信仰を構築する上で無価値でもない価値があります。しかし、誤った信念に導かれないようにするためには、神秘的感動が何を明らかにするか正確に理解することが必要です。それは人間性の可能性を明らかにします—より優れた、幸せで自由な生活の可能性です。しかしそれは非人間的なものや宇宙全体の本質について何も明らかにしません。善悪、そして神秘主義が見出す高い善も、我々自身の感情の反映であり、実在としての物事の一部ではありません。したがって、自己に関する一切の執着から解放された公正な静観は、善または悪と判断することはありませんが、普遍的な愛という感覚とともに容易に調和します。そしてそれは神秘主義者が世界全体を善とすると言える感情と組み合わさります。

進化の哲学は、前進という概念を通じて、悪いものと良いものの倫理的二分性と結びついており、その結果、善悪を全く見落としているような調査からは排除され、またすべてが善であるという神秘的な信仰からも排除される。このようにして、善悪の区別は時間のように、この哲学において暴君となり、思考に行動の不穏当な選択性を導入する。

善悪は時間と同様に、思想の世界では一般的または根本的なものではなく、知的階層の中での遅く且つ高度に専門化されたメンバーであると考えられる。

確かに、如実に見れば、神秘主義は善悪が知的に根本的ではないという視点を支持するように解釈できる。しかし、ここでは過去の多くの偉大な哲学者や宗教指導者と口頭での一致から遠ざかる。しかしながら、哲学における倫理的な考慮事項の排除は、科学的には必要であり、（これはパラドックスに思えるかもしれませんが）倫理的な進歩でもあると考えている。

これらの主張を簡潔に防衛する必要がある。

我々が見たように、神秘主義は善悪が知的に根本的ではないという見解と一致させることもできる。しかし、ここでは過去の多くの偉大な哲学者や宗教指導者との口頭での一致から遠ざかる。しかしながら、哲学における倫理的な考慮事項の排除は科学的には必要であり、（これはパラドックスに思えるかもしれませんが）倫理的な進歩でもあると考えている。

これらの主張を簡潔に防衛する必要がある。

我々がより人間的な欲望への満足という希望—世界がこのまたはその他の望ましい倫理的特性を持つことを証明することの希望—は、私が見解できる限り、科学的な哲学が何らかの方法で満たすことができるものではない。良い世界と悪い世界的差異は、これらの世界に存在する特定の物体の特定の特性の差異である：それは十分に抽象的ではなく、哲学の範疇を超えている。愛と憎しみは倫理的な対照であり、しかし哲学にはこれらに対する物体に対する態度は密接に関連しているものである。これらの物体に対する態度の一般的な形式と構造が哲学の問題である一方で、愛と憎しみの差異は形式や構造ではなく、心理学の専門科学に属するものである。したがって、しばしば哲学者を動機付ける倫理的な興味は背景に留まらなければならない：全体的な研究には何らかの種類の倫理的興味が刺激してもよいが、詳細には強制的に現れることなく、特定の結果を求められるものには期待されることはない。

この見解が最初に見せられたときには失望するかもしれません。しかし、他のすべての科学でも類似した変化が必要とされることが発見されてきたことを思い起こすことができます。物理学者や化学者は、そのイオンや原子の倫理的な重要性を証明する必要はありません。生物学者は、解剖した植物や動物の有用性を証明する必要もありません。

前科学的時代にはそうではなかったです。天文学者が研究されたのは、人間の生活に直接かつ重要な影響をとえると考えられていた惑星の動きがあったからです。この信念が衰え、無関心な天文学が始まったとき、多くの人が占星術に熱中していたために天文学が人間にとってあまり興味深いものではないと判断し、研究する価値がないと決めたかもしれません。

プラトンの『ティマウロス』における物理学は倫理的な概念で満ちています：それは地球が称賛に値することを示すことがその目的の一部です。一方、現代の物理学者は、地球が称賛に値することは認めますが、物理学者としてその倫理的属性に関心はありません：彼はただ事実を見つけ出すことにのみ関心があります。善悪や良いものか悪いものかを考える必要はありません。

心理学では、科学的な態度は物理学よりもさらに最近で困難です：人間性が善か悪かであると考えるのは自然であり、善と悪の違いが実践において最も重要であるからこそ理論的にも重要なことだと考えるのは当然です。倫理中立な心理学がこの世紀末に成長したのは最近のことです；そしてここでも、倫理的中立性は科学的成功にとって不可欠でした。

哲学においては、これまで倫理的な中立性を希求することは稀で、実現されることもほとんどなかった。人々は自分の願望を思い出し、哲学をその願望に関連させて評価した。特に科学から追い出された信念は、善悪の概念が世界の理解の鍵を提供するものであるという考えが、最終的な refuge として哲学に避難した。しかし、哲学が楽な夢の集まりにならないよう、この信念さえも排除しなければならない。幸せを得るためには直接それを追求しない方が良いことは常識であり、善についても同様のことが言えるかもしれない。

思想においては、善悪を忘れ、事実だけを求めようと努力する人々の方が、自分の欲望を通じて世界を見ることで善を得るよりもより多くの善を得やすいだろう。

従って、哲学が自らの善悪の概念を強制しようとしない場合、真実を達成しやすく、また進化論や伝統的なシステムのように常に宇宙を評価し、現在の理想の具現を探求するよりも高い倫理的立場であるというパラドックスに帰着される。宗教において、そして世界と人類の運命に対する深い視点においては、人間の力の限界を理解した屈服の要素があるが、これは現代社会では欠如している。その社会は急速な物質的成功を収め、進歩の無尽蔵なる可能性に傲慢な自信を持っているからである。「自分の命を愛する者は失う」ことから、この自信過剰な愛を通じて、生命自体が最高の価値を持つものの中には多くを失う危険がある。宗教が行動において灌输する屈服は、思想において科学が教える屈服と精神的には本質的に同じであり、その勝利を得た倫理的な中立性はこの屈服から生じている。


私たちが記憶すべき善は、私たちの力に属する善であり、自分たちの人生と世界に対する態度にある善です。外部での善実現への信念強調は自己主張の一形態で、その形式は外的善を確保することはできず、内部的な善を損なう可能性があり、事実に対する敬意を破壊し、卑屈さにおける価値と科学的精神における成果性の両方を毀損します。

当然、人間は完全に人間性から超越することはできません；私たちの思考において何か主観的な要素、つまり注意の方向を決定する興味さえも残ります。しかし、科学的哲学は他のいかなる人間の追求よりも客観性に近づいており、したがって、外的現実との最も密接で最も深い関係を達成することができます。原始的な心には、すべてが友好的または敵対的であるという考えがありますが、経験はそれが世界を理解するための概念ではないことを示しています。科学的哲学は、まだ未発達ではありますが、前科学的信念や想像よりも高い形態の思考を代表しており、自我的超越へのすべての接近と共に、範囲と広がりそして理解力の豊かな報酬をもたらします。

進化論は特定の科学的事実に訴えることにより真の科学哲学とはなり得ません；それは時間への従属性、倫理的な関心事、そして私たちの日常的な懸念と運命に対する主要な興味から解放されていないからです。真の科学的哲学はより謙虚で、より段階的で、より困難であり、外的幻影を誇示する光沢が少ないが、運命への関とが少なくて、世界を受け入れる能力が高いです。

フ�注：

 以上引用はBurnetの『初期ギリシャ哲学』（2版、1908年）第146-156ページから引用しました。

 \textit{共和政書} 514番目、 Davies and Vaughanによる翻訳。

 このセクションと、後の一部のページはローレル講義『私たちが外部世界について知る方法』の一環として出版会社Open Court Publishing Companyによって出版されました。しかし、これは元々書かれた文脈のためにここに残しました。


 メタ物理学入門，第1頁。

 ウィンフィールドによる《マサニビ》の翻訳（トゥーベナー，1887年），第34頁。

 規則，IV巻，第LXII条。

 同上，IV部，定義I。

 規則，II部，定義VI。


\chapter{第II章：教養教育における科学の位置 --- {THE PLACE OF SCIENCE IN A LIBERAL EDUCATION}}

II

科学の地位について


I

新聞を読む一般の人々にとって、科学は無線電信や飛行機、放射活性や現代の錬金術のような興奮的な成果の一端を示すものとして表現される。しかし、私が話したいのはこの面ではなく、科学そのものについてである。この面での科学は、新しいものに取って代わられるまで一時的に興味深い最新の断片から成るだけのもので、慎重に構築された知識の体係からはほとんど偶然的に得られた実用的な結果を示していない。自然の力に対する人間の支配力を科学がとえることは確かに科学的研究を奨励する十分な理由であるが、この理由はしばしば強調されやすく、容易に理解されるため、他の重要な理由がしばしば見落とされがちである。以下では、これらの他の理由、特に科学的な思考習慣が私たちの世界観形成における内在的価値について述べることにする。

無線電信の例を用いてこの二つの視点の違いを説明する。この発明の可能性に必要なほとんどすべての真剣な知的労働は、ファラデー、マクスウェル、ヘルツの三人によってなされた。実験と理論の交互的な層から、彼らは現代の電磁気学の理論を築き上げ、光が電磁波であることを証明した。彼らが発見した体係は、深い知的興味をもたらし、無数の表面上に分断された現象を結びつけて統一するものであり、その積み重なる精神的力量は、誰の心にも喜びをとえること必至である。実用的な電信システムのためにこれらの発見を利用するためには残る機械的詳細は確かに大きな巧みな工夫が必要だったが、それらが無関心な観察の対象として内在的な興味を持つべき広大さと普遍性を欠いていた。

从精神训练的角度来看，从文学教育中获得的那種博学而客観的視野，構成了在广泛意義上使用これ个詞所指的文化，似乎普遍認為これ種教育优于基于科学的教育。即使是熱衷于科学的人也往往以牺牲文化来換取実用性的主張為基礎。那些尊重文化的科学家，在と通晓古典学問的人交往時，往往会诚恳地承認自己在某些方面処于劣势，尽管他達可能弁称科学対人類有巨大的贡献，但これ並不改変これ種劣势的真実性。只要これ種态度存在于科学家之中，它就会自我験証：科学中内在有价値的方面往往被実用性的需求所牺牲，而很少有人尝试保持那種能够培养和滋养精细品質的从容不迫、係統化的审視。

即便在目前的事実中，科学教育的价値確実存在某種被認為低于文学教育的劣势，我認為これ並非科学本身的過錯，而是由于以某種精神来教授科学而導致的結果。如果那些教授科学的人能充分実現其全部潜力，我相信它能够産生構成最高精神卓越习惯的能力至少と文学，特别是希腊和拉丁文学相当，甚至更為优越。我説これ些並不是贬低古典教育的价値。我自己並未享受過これ種教育的好処，我対古希腊羅马作家的了解几乎完全来自译本。但我坚信，古希腊人確実値得我達給予他達所有的赞美，並且対他達的著作一無所知是一種非常大的损失。我並非要攻击科学，而是希望通過强調科学中被忽視的优点来展开我的論述。

一方、純粋に古典的な教育には欠点があるように思われる──つまり、過去への過度な偏重である。絶対的に終わったことでかつ再生不能なものだけを学ぶことで、現在と将来に対する批判的態度が養われてしまう。現在において優れている品質は、過去の研究では注目されないものであり、したがって、ギリシャ文明を学んだ学生はこれらの品質を見落としてしまう可能性がある。新しいことや成長するものはしばしば未熟で傲慢、さらには少し俗っぽく、感度の鋭い人には不快な存在となる。そのような人たちは粗末さから身を引いて、整然とした過去の庭園に退去し、忘れてしまうことが多い。彼らが避ける人々は、自分自身の時代と同じくらい荒々しく土臭い人々であることを忘れる。

「昔の人を見ることはなかった；
そして彼らの時代は過ぎ去った。
私は悲しみ──英雄たちを見ることはできないことを思うと」

中国の詩人のように言うが、このような公正さは西で激しい空気の中では稀であり、過去と未来を擁護する戦士たちは永遠に続く戦いではなく、両方の長所を探求するために協力することを選ぶことはめったにない。

この考慮は、古典学の単純な研究だけでなく、静止し、伝統的で学術的なすべての文化的形態にも反対するものであり、結局のところ根本的な問いへと導かれる：教育の真の目的は何であるのか？しかし、この質問に答える前に、我々が「教育」という言葉をどのように使用するかを定義することが適切であろう。そのためには、「教育」の意味を私が用いるものから二つの他の、どちらも完全に正当な解釈、つまりより広範なものとより狭い範囲のものとの区別が必要である。

教育は、狭義では指導を通じた情報の伝達に限定されるが、広義では教養を形成し、生活を通じて人格を育てるすべての経験も含む。この面について非常に重要な点だが、その議論は私たちが扱っている問題とは全く異なるものになるため、ここでは何も述べない。

狭義での教育は、日常生活に有用な様々な分野に関する具体的な情報の伝達に限定されるかもしれない。基本的な教育、すなわち読書、書き込み、算数はほとんどそのような性質である。しかし、必要な指導自体が、私が考える教育の意味においては必ずしも教育を形成するものではない。

私がいう教育とは、指導を通じて特定の心の習慣と生活や世界に対する特定の見方を形成することである。では、どのような心の習慣とどのような見方が、指導によって期待できる結果として得られるだろうか？この問いに答えた後で、我々はどの科学が求められる心の習慣と見方の形成に何を提供できるかを決定することができるだろう。

私たちの全人生は、特定の数少ない基本的な本能や衝動によって構築されています——それらが非常に小さな数ではないのです。これらの本能や衝動と何か関連のあるものは我々には魅力的または重要視されます。理知や徳、あるいはそれを呼ぶどんな言葉でも、私たちの活動的な人生や希望、恐怖をこれらの最初の欲求の制御から外すことはできません。それぞれは女王蜂に例えられるべきで、働き蜂たちが蜂蜜を集めている蜂房のようなものがありますが、女王蜂がいなくなると働き蜂たちは衰弱し、死んでしまい、巣箱には期待された甘さが残りません。文明人における各基本的な衝動も同様です：それ周囲に働く派生的欲望の群れによって囲まれ保護され、そのサービスのために周囲の世界から蜂蜜を貯蔵します。しかし女王衝動が死ぬと、その死亡の影響は習慣により少し遅れて、すべての付属的な衝動を通じて広がり、人生の一帯が不思議なほど色褪せてしまいます。以前は豊かで、明らかにすべき事柄でもなかったのでやる価値があったものが今では退屈で目的がないものになっています：幻滅感を抱きながら人生の意味を問い、結局すべてが虚無であると判断するかもしれません。外側から求めた意味を探し求めても、内面的な反応を得ることは決して叶いません：すべての「意味」は根底から私たちの基本的欲望に関連しているもので、それらが絶たれると世界に再び価値をもたらす奇跡は起こりません。

教育の目的は、未教育者に欠けている何かを創出することではなく、人間性がとえるものの範囲を拡大することである。その目的は、付随する思考の数と種類を増加させることで、最も永久的な満足を見つける場所を示すことであり、それによって未開化の人間に対するカルヴァン主義的恐怖心から生じたこの明白な真実が若者教育においてしばしば誤解されたことは否めない。自然は「最高の自然」を含まないと誤って考えられ、「性質」はすべてが自然の中で最も良いものも除外するとされ、徳を教える努力は、成長した人間ではなく、短く曲がった偽善者を生み出した。このような教育上の誤りから、よりよい心理学や優しい心が現在の世代を守ろうとしているため、教育の目的が自然を妨げたり根絶したりすることであるという理論にはもう一つの言葉も必要ない。

しかし、自然は文明人において、野蛮人が持つような一時的で断片的かつ激しい反応ではなく、それぞれの反応はその性質によって思考と知識そして反省の役割を果たし、衝突する反応の可能性を予見し、一時的な反応は統合反応である知恵によって制御される。このようにして教育は本能の粗野さを破壊し、知識を通じて個々人の外部世界との接触の豊かさと多様性を増加させ、彼を孤立した戦闘ユニットから宇宙の市民に変える、遠くの国々、空間上の遠くの地域、そして過去と未来の大いなる範囲が彼の関心の円の中に含まれる。この欲望の堅持性の緩和とその範囲の拡大は、教育の主要な道徳的な目的である。

この道徳的な目的と密接に関連しているのが、より純粋に知的である教育の目標で、私達を個人的な欲望を通じた歪曲した視点ではなく、可能な限り本来の形での世界を見ることや想像することを試みる努力です。このような客観的な見解の完全な獲得は、確かに理想であり、無限に近づくことができるものの、実際的に完全には達成できないものです。教育が、私達の心の習慣と世界に対する視点形成という過程として考慮されるなら、その成果がこの理想にどれだけ近づいているかで成功を評価すべきです。つまり、それが私達が社会の中でどのような位置にあるのか、全人類社会が非人間環境に対してどのような関係にあるのか、そして非人間の世界が私たちの欲望や利益から独立してどのような性質を持っているのかという真実な見解をとえるかによってです。この基準を認めれば、科学について再考し、そのような目標にどれだけ貢献しているか、また教育において競争相手よりも優れている点があるかどうかを調べることができます。

II

科学に対して文学や芸術と対照的な二つの相反する価値が存在します。一つは先天的に必要ではなく、しかし現在のところ真実である価値で、それは人類の達成に対する将来への希望感であり、特にどの程度知的な学生によって有用な仕事を成し遂げることができるかという点に特に関連しています。この価値とそれを生み出す明るい見解は、科学の別の側面がもたらす可能性のある depressiveness を防ぐことができます。私の考えでは、これは科学における人間の情熱や全体的な主観的装置に関わる不関心性であり、おそらく最大の価値でもあります。つまり、科学の真実性に関しては人間の情熱が関係ないということです。これらの理由のそれぞれについて説明を加える必要があります。まず最初のものから始めましょう。

文学や芸術の研究においては、私たちの注意は常に過去に向けられています：ギリシャ人やルネサンス期の人々が今よりも優れていたこと；前代の勝利は、現在の新たな勝利を促進するどころか、独創性を達成する難易度を増加させます。芸術的な成果は累積的ではなく、文明化によって特定の鮮やかさと素朴さが必要とする衝動と視野に依存しているようにさえ思われます。したがって、過去の文学や芸術作品に養育された人々には、現在に対する不快感と過度な厳格さが生じ、その脱出手段は故意の破壊主義にあり、それは伝統を無視し、独創性を求めることでしかるべきだけの eccentricity を達成することになります。しかし、偉大な芸術から生まれる純朴さと spontaneity には欠けますが、理論はその核心に癌のように存在し、偽りは単なる装った無知の利点を破壊します。

このような破壊主義からは、教育が優れた精神活動を提示しない場合に生じる絶望感が全くありません。科学的方法に関する知識をとえる教育からです。科学方法の発見（純粋数学を除く）は昨日のことです；広く言えば、それはガリレオから始まったと言えます。しかし既にそれが世界を変え、その成功は常に加速的な速度で進行しています。科学では、進歩のために常により大きな天才が現れる必要性がないことが発見されました：科学の後継者は前人を肩に載せることができます。一人の最高の天才が方法を発明した場合でも、千人の下位の人物がそれを適用できます。科学における有用な発見を行うためには超人的な能力は必要ありません；科学の建築物には監督者、石工、労働者が含まれますし、それらと同じくらい重要なのは先任者の肩に立つことができます。芸術では、価値のあることを成し遂げるためには天才が必要ですが、科学では中程度の能力さえも最高の達成に貢献できます。

科学において真の天才は新しい方法を発明する人間である。
重要な発見はしばしばその方法を継承した人物によってなされることがあり、彼らは前の労力を損なうことなく新たな活力を持って方法を適用できる。しかし、その作業に必要な思考の質は第一義的な方法の発明者よりも brillant ではあるが、それほど大きくはない。科学には異なる問題に対する適切な方法が非常に多く存在するが、それらを超えて何か定義しにくいものがあり、それが科学の \textit{方法} と呼ばれるべきものである。かつてこれは帰納法と同一視され、 Bacon の名前とともに語られたことが多かった。しかし、真の帰納法は Bacon によって発見されたものではなく、科学の真の方法は推論も含む帰納法であり、論理や数学も植物学や地質学と同じく含まれる何かである。私は科学的な方法自体を述べることは避けるが、それから生じる思考の態度を示すことに挑戦する。

科学的見解の核は非常に単純で、明らかで、また表面上は trivia に思えるほどであるため、その言及はほとんど嘲りを誘うかもしれない。科学的見解の核は、私たち自身の欲望、趣味、そして利益が世界を理解する鍵ではないという拒否である。これを簡潔に述べれば、これ自体は単なる陳腐な真理に過ぎないように思えるかもしれない。しかし、情熱的な偏見をも掻き立てる問題において常にそれを思い出すことは決して容易でなく、特に証拠が不確定かつ結論的ではない場合尤其である。いくつかの例を挙げることでこの点が明確になるだろう。

アリストテレスは、星が円周運動をとる理由として、円は最も完璧な曲線であると考えていたことを理解しています。反証のない状況において彼は、美術的・道徳的な考慮に基づいて事実に関する質問に答えることを許しました。このような場合、この考慮が正当化されないとすぐに明らかになります。現在では天体の動きを事実として確認する方法を知っていますし、彼らは円周や正確な楕円ではなく、どのような単純に説明できる曲線でも動いていません。これは宇宙の模式的な単純さへの一定の欲求に痛みをとえるかもしれませんが、天文学においてこのような感情は無関係です。この知識が容易であるように思えても、それは科学的方法を最初に発明した人々の勇気と洞察力に感謝しなければなりません。特にガリレオに感謝しなければなりません。

また別の例として、マラトゥスの人口に関する教義があります。彼の実際の教義が現在では大半が誤りであるという事実はこの例をより良いものにしています。彼の結論よりも重要なのは、彼の調査の態度と方法です。誰もが知るところですが、ダーウィンの自然選択説の一部は彼から得たものですし、これはマラトゥスの見識が真実の科学的だったためでした。彼の大切な功績は人間を称賛すべき対象や非難すべき対象としてではなく、自然の一端として、特定の特性を持つ行動から必然的に結果が生じる存在として考えることにあります。もしその行動がマラトゥスが想像したものとは少し異なる場合、または彼が推論した結果が完全でない場合、それは彼の結論を誤りにするかもしれませんが、彼の方法には影響をとえません。彼の教義が新しいときの反対意見は全て科学的態度を持たないものでした；それに対して彼の冷静な人間を自然現象として扱う決意は18世紀の改革者や革命に対し重要な前進を示しています。

darwin主義の影響により、人間に対する科学的な態度は比較的広く普及し、一部の人々には自然なことのように思えるが、多くの人々にとってはまだ困難で人工的な知的な扭曲である。しかし、一つだけ研究分野が科学的精神にほとんど触れられていない－それが哲学の研究である。哲学者や一般の人々は、イオン、胚質、貝類の目などといった参照を伴うページでも科学的精神が満ちていると誤解している。しかし、悪魔も聖書を引用するようにして、哲学者も科学を引用することができる。科学的精神は引用や外部から得た情報の問題ではなく、マナーも儀礼本の問題ではない。科学的な態度とは、知識への欲求の利益のために他のすべての欲望を排除すること－希望と恐怖、愛憎、そして全体的な主観的感情生活を抑制することである。これにより私たちは物質に屈服し、事実を公正に見ることができ、前もっての仮定や偏見なく、何が現実であるかを見るだけの願望を持つことになる。また、それが私たちが望むように決定されるものではなく、または容易に想像できるものとは関連性がないと信じることもない。

哲学においてはこの態度がまだ達成されていない。全体的な宇宙を概念化する試みでは、個人的な自閉症ではなく、人間的な自己中心性がほとんど全ての試みで見られる。心、あるいはその一部－思考や意志や感覚など－が宇宙を概念化する模範とされ、その理由はただそれが奇妙に思われないし、どこでも家にいるような暖かさを感じることができるからである。進歩的な宇宙や衰退的な宇宙を概念化することは、私たちの希望と恐怖に宇宙的な重要性をとえることになる－もちろんこれは正当化される可能性があるが、現在までにはその正当性を証明する理由は存在しない。倫理的に中立的な言葉で考えることを学ばない限り、哲学における科学的精神態度は達成されていない；そのような態度に到達しない限り、哲学は確固たる結果を達成することは難しいだろう。

科学精神には、これまでに述べたように主に否定的な側面があるが、その価値は肯定的な側面から得られる。創造の本能は芸術的創作の主要な動機の一つであり、科学システムの中でより大きな満足を得ることができる。無関心な好奇心は、ほとんど全ての知的努力の源であり、科学が永遠に発見されないか以為「科学が永遠に発見されない可能性がある」秘密を明らかにする驚きと喜びを見出すことができる。より広い人生と幅広い興味、私的な状況から脱出すること、甚至「すなわち人間の生と死という繰り返される人間のサイクルから脱出すること」は、科学の無差別的世界観によって満たされる。これらすべてに加えて、科学者にとって幸せをもたらすためには、偉大な成績への敬意と人類に対する無価値ではない有用性の意識が加わる。因此「したがって」科学に専念する人生は幸せであり、その幸せはこの不安定で情熱的な惑星上で住む人々が利用できる最高の源から得られる。


\chapter{第III章：自由人の祈祷 --- {A FREE MAN'S WORSHIP}}

III

自由人の崇拝


ファウスト博士の書斎でメフィストフェレスは創造の歴史を語り、次のように言った：

「天使たちの無限の賛美がやがて退屈なものになってきた。果たして彼自身にその賛美が必要だったのか？ 彼が彼らに無限の喜びをとえたのであれば、彼が苦しめた存在たちによって無実の賛美を受けることの方がもっと面白いではないか。彼は内心微笑み、偉大なドラマを上演することを決意した。

「無数の時代にわたって熱い星雲は空間を放浪的に旋回していた。やがてそれは形をとり始め、中心部から惑星が放出され、惑星は冷却し、沸騰する海と燃える山が揺れ動き、黒い雲塊から熱した雨のシートが僅かな固体の地殻に降り注いだ。そして今や海洋の深淵で最初の生命の芽が育ち、生殖暖かさの中で急速に巨大な森林の木々、湿った土壌から生える大きな fern、海の怪物たちが繁殖し、戦い、食いちぎられ、消滅した。

「そして怪物たちが物語が展開するにつれて人間が生まれ、思考力を持ち、善と悪を知り、残酷な崇拝への渇望を持つようになった。そして人間は狂気で異常な世界においてすべてが瞬時に死の不変の命令前に奪われる争いに見えていた。そして人間は言った：「隠された目的があるはずです。それを理解できればそれは善であり、何かを崇拝する必要があり、視覚的世界には崇敬すべきものが何もない」と。そして人間は争いから立ち去り、神が混沌から調和を人間の努力によって生み出すことを意図していると決意した。そして彼が祖先からの獣類の本能に伝わったものを追跡するとき、彼はそれが罪だと呼び、神に赦しを求めた。しかし彼は自分自身が正しく赦されるかどうか疑問を持ち続け、神を慰めるための神聖な計画を作り出すまで、その疑念から解放されることはなかった。そして現状が悪いと見ると、それをさらに悪化させることで将来が良くなることを目指した。そして彼は神にとえられた力を感謝し、可能な喜びすら放棄することもできた。そして神は微笑んだ。そして彼が人間が断念と崇拝において完璧になったことを確認すると、もう一つの太陽を空に送り込んだ。それは人間の太陽に衝突して再び星雲に戻った。

「そうだ」と彼はつぶやいた、「それは良い芝居だった；また上演させようと思う」。

SUCH、概略を除いてはなおさら、目的のないもの、意味のないものが科学が我々に信じさせる世界である。そのような世界の中で、あるいはその中にこそ、我々の理想は今後家を構えるべき所がある。人間は予見も無く成就する目的のために存在した原因の産物であり、彼の起源、成長、希望と恐怖、愛と信仰はただ原子の偶然の配置の結果である。火の力、英雄主義、思想感情の強さは個々の人間の生命を死者の太陽係の大死にまで保存することは不可能である。人類の偉業のすべての努力、すべての献身、すべての啓発、すべての正午の天才の輝きは、恒星係崩壊の広大な死の中で滅びる運命にあるし、人間の成就の全廟も必ずや破片の下に埋もれてしまう。これらのことは、争いなく否めないものではないが、それほど確実であるから、これらを拒否する哲学は立脚点を見つけることができない。これらの真実の枠組み内にのみ、そして絶対的な失望の堅固な基礎上にのみ、今後人間の魂の住処は安全に築かれるであろう。

そのような異質で非人間的な世界において、力の弱い存在である人間が自分の抱負を未だ汚染されずに保つことができるのかどうか。自然は無限大でありながら盲目であり、彼女の千年の流れの中で宇宙の深淵を駆け抜ける回転によって最後に子供を生み出した。その子供はまだ彼女の力に従属しているが、彼には視覚があり、善悪の知識があり、無思慮な母のすべての作品を判断する能力がある。死という父なる支配者の印と記章にもかかわらず、人間は短い生涯の中で検討し、批判し、知ること、そして想像力を用いて創造することの自由が残されている。彼一人だけが知る世界においてこの自由があり、これが彼が彼の外在的な生活を支配する力に対して優越性を持つ所以である。

野蛮人は自分自身よりも自然の力に屈服する感情を共有しますが、自分の中でより尊重すべきものは権力だけであるため、彼は自分の神々前に膝を屈め、その神々が自分たちに値するかどうか尋ねることなく、彼らに屈服します。残酷で恐ろしい長い歴史があります。残忍さと拷問、卑下と人間の犠牲が、神々を喜ばせるために受け入れられました。震える信者にとって、最も貴重なものまで自由に捧げられたとき、血の渇望は満たされると考えられますし、それ以上は要求されないでしょう。モロクの宗教というような一般的な信仰は、独立した理想が認められていないため、権力は自由に崇拝され、その恣意的な痛みをもってしても無制限の尊敬を集めることができます。

しかし、道徳が徐々に強くなるにつれて、理想的世界への要求が感じられるようになります。そして、信仰が止まらないなら、野蛮人が作る神々ではなく、別の種類の神々に捧げなければなりません。理想の要求を感じながらも、裸の権力は崇拝に値すると意識的に拒否する人々もいます。その態度は whirlwind から Job に対して神が示す答えに見られます。神の力と知識が晒されますが、神の善さには触れられません。また、私たち自身の時代において生存競争に基づく道徳を基礎とする人々も同じ態度を取っています。生き残った者が必ずしも最適であるという主張です。

しかし、道徳感に反するような答えを受け入れられない人々は、事実の世界と理想的な世界が隠された形で調和しているという立場を採ります。人間は全能かつ全善なる神を作り出し、現実と理想の神秘的な統一性を創造します。

しかし、事実の世界は毕竟善ではなく、その判断に従うことは奴隷性の要素を含み、私たちの思考から除去しなければならない。すべてのことにあって人間の尊厳を高めるためには、可能な限り非人類的な力からの支配から解放することが望ましい。事実上、力を大きく否定し、知識を持った人間は善悪を知るが、無知な世界の中でただの無力な原子に過ぎないという現実を理解したとき、再び我々の前に選択肢が訪れる：力に捧げることを選ぶべきか、還是善に捧げることを選ぶべきか？私たちの神は悪である存在なのか、それとも私た各自的意識によって創造された存在なのか？

この問いに対する答えは非常に重要で、私たちの全体的な道徳に深刻な影響を及ぼす。カリー、ニーチェ、軍国主義の信条によって我々が慣らされた力への崇拝は、自分たちの理想を敵対する宇宙から守る能力を失った結果であり、それは悪に対する屈服そのもので、私たちの最も良い部分をモレッコに捧げた犠牲である。確かに強さを尊重すべきであれば、偽りの「事実への認識」が善悪を知らぬ事実を認めないという誤った認識をもたらすことを拒否する強さを尊重すべきである。我々が生きている世界において多くのことが他に比べて望ましくないと理解し、我々が持続的に追求しなければならない理想は物質の領域で現実化されていないと認めるべきである。真実、美しさ、生活が達成できない理想的な完璧さに対する尊重を失うことなく、これらのすべてが無意識的な宇宙には認められないとしても、それらを尊重すべきである。

力が悪であるように思えるなら、心からそれを拒否すべきである。これが人間の真の自由であり、私たち自身の善への愛によって創造された神にのみ崇拝する決意と、私たちの最高の瞬間に感得される洞察を誘発する天に対する尊重から解放されることである。行動においては、欲望においては、外からの力に永久的に屈服しなければならないが、思考においては希望においては自由であり、他の人間から自由であり、無力な身体が無力に這いずる小さな惑星から自由であり、生きている間は死の支配からも自由である。そして、その視覚を常に前にして行動に降り立つことで、信仰のエネルギーを学び、善のvisionの中で生活することを学ぶべきである。

当実在と理想の対立が完全に明らかになると、自由を主張するために、ファイアリな反乱の精神や神への激しい憎しみが必要となる有時。敵対的な宇宙に対してプロメテウスの恒常性で挑戦し、その悪を常に視野に入れ、常に積極的に憎むこと、力の悪意が思いつくすべての苦痛も拒否しないことが、屈服せずにいる人々の義務であるように思われる。しかし、憤慨はまだ枷であり、それは私たちの思考を悪のある世界に集中させてしまうからだ；反乱の源となる激しい欲望からは、一種の自己主張が生まれ、賢明な者はこれを克服しなければならない。憤慨は思想の屈服であるが、欲求の屈服ではない；賢明さが含まれるスコピックの自由は、欲求の屈服ではなく、思想の屈服から見つかる。欲求の屈服から生まれるのは、諦めの徳であり、思想の自由からは芸術や哲学の世界と美の視覚が生まれ、ついに我々が不情願な世界を半分以上取り戻すビジョンが生まれる。しかし、美のビジョンは、制約のない観察、切迫した希望の重みから解放された思考によってのみ可能であり、そのような自由は、個人的な利益を生活から求める必要がなくなった人々にのみ訪れる。 

虽然放弃的必要性証明了悪的存在，但基督教通過宣揚これ一点，展示了超越普羅米修斯反叛哲学的智慧。必須承認的是，在我達所渴望的事物中，有些尽管最終証明是不可能実現的，但仍是真正的善；而另一些，尽管同様熱切地渴望，也不構成一个完全浄化的理想的一部分。放弃的東西必然是悪的これ一信念虽然有時可能是錯誤的，但比未馴服的熱情假設的錯誤要少得多；宗教信仰通過提供一个理由来証明これ一点从不為假，从而成為通過発現許多厳峻真理而浄化希望的手段。

しかし、屈服にはさらなる善な要素があります：実在の善が手に入らない場合でも、不満げに望むべきではありません。誰しもいずれ、大きな屈服を経験します。若い者は何も手に入らないという状況には遭遇しません；全般的な意志力を総動員して望んだ良いものが不可能であるとすれば、それは彼らにとって信じられないことです。しかし、死、病気、貧困、または義務の声によって、私たちはそれぞれ学びます：世界は私たちのために作られたものではなく、我々が求めるどんな美しいものでも、運命がそれを許さないかもしれません。不幸が来たら、希望の崩壊を嘆きつつ耐える勇気が求められます；虚しい悔いから心をそらすことも求められます。こうした力に屈する程度は、ただ正しく公正なだけでなく、知恵の門でもあります。

しかし、屈服だけでは智慧の全容は語れません：我々自身の理想を崇拝する神殿を築くためには、単なる屈服だけでは足りません。想像の領域、音楽、建築、理由の未揺るがす王国、そして詩歌の金色の夕暮れの魔法の中で、美は遠く離れて輝き光ります；悲しみの手触りから遠く、変化の恐れから遠く、事実の世界での失敗や幻滅から遠く。これらの物事を観察することで、天界の視覚が心に形成され、その結果、周囲の世界を判断するための試金石と、必要なものは神聖な神殿の石として使用できるように形作るためのインスピレーションをとえます。

罪のない生まれた稀な魂以外は、神殿に入る前に闇の洞窟を通らなければなりません。洞窟の門は絶望であり、その床には放棄された希望の墓石で舗装されています。己が死ぬべき場所です；未制御の欲望の切迫感や執着を殺すべき場所です。しかしここで魂は運命の支配から解放されます。しかし洞窟から出てきた屈服の門は、智慧の光に照らされた明るさへと再び導きます。その光によって新しい洞察、新しい喜び、新しい優しさが旅人を喜ばせます。

無力な反乱の渦中でもなく、運命対外部法則的順从と、非人間的世界を崇拝する価値lessnessないことに気づき、学んだとき初めて、無意識の宇宙を変形し、想像の鍛冶炉でその性質を変え、古い泥製の偶像に代わる輝く金の新像が現れることが可能となる。世界の多様な事実——木々や山や雲の視覚的な形状、人間生活の出来事、すら死の無尽蔵なる力さえも——創造的イデアル主義の洞察は、その自身の思考によって初めて生み出された美を発見することができる。この方法で、心は自然の無意識的な力を微細な支配権を持つことになる。扱う素材がより悪ければ、未訓練した欲望に反抗的であれ、屈服させた岩から隠された宝物を得る成就是より大きくなるし、反対する力に屈しない都市を征服することで誇り高き勝利を収めることになる。すべての芸術の中で、悲劇は最も誇り高き、最勝利的なものである。それは敵国の中心、彼の最高山頂にその輝く要塞を築くからである。その不破の望楼から、彼の軍隊と兵器庫、列と砦がすべて明らかになる。その壁内では自由な生活が続けられ、死と苦しみと絶望の連合軍や、暴政運命の奴隷将校たちが、その胆略なる都市の市民に新たな美の見世物を提供する。聖なる土塀は幸福であるし、その全てを見渡す丘に住む者は三度の幸いである。自由の無价な遺産と、不屈の家屋から邪教の侵略者によって汚染されぬことを守り続けた勇敢な戦士たちには栄光を捧げよう。

悲劇の美しさは、生命に常に存在し、明暗を帯びた形で広がっている性質を明らかにするだけである。死の光景、不忍耐な苦痛の耐え忍び、消えた過去の不可逆性から、存在の偉大さ、深淵、無尽蔵の神秘性を感じ取る聖なる感動、圧倒的な畏怖、悲しみに縛られた世界との結びつきが現れる。これらの洞察の瞬間では、一時的な欲望の切迫感は失われ、小さな目的への争いと努力も消え、日常の平凡な生活を構成する trivial なことば一切から目を離すことができる。狭い筏が人間の友情の光で照らされる中、暗い海洋の波に短時間に浮かぶ姿を見ることができる。外側からの寒風が避難所に入り込み、敵対的な力の中での人類の孤独が個々の魂に集中し、その魂は、希望と恐怖を無視する宇宙に対して、手に入るだけの勇気で戦うしかない。この闇との戦いにおける勝利は、英雄たちの豪華な集団への真の洗礼であり、人間存在の圧倒的な美しさへの真の入門である。魂と外側世界の恐るべき遭遇から、放棄、智慧、慈愛が生まれ、それらと共に新たな人生が始まる。不可避の力である死と変化、過去の不可逆性、そして盲目の宇宙の流れに立ち向かう人間の無力さを魂の最深部に取り入れる。これらのことを感じ取り理解すれば、その勝利を得ることができる。

この理由故に過去はその魔法的な力を持つ。静寂で沈黙したその絵画の美しさは、秋の末期の妖精のような純粋さに似ている。葉が一息で落ちるとしても、空に金色の輝きを放ちながらまだ光っているように、過去は変化せず、努力もしない。ドuncanのように人生の不規則な熱狂後には平和に眠りにつく。急せき立てるものや微細で一時的なものは消え去り、美しい永遠のものはその中から星のように輝いている。その美しさは、価値のない魂にとっては耐え難くもあろうが、運命を征服した魂にとっては宗教の鍵となる。

人間の人生を見れば、自然の力と比べて微々たるものである。奴隷は時間や運命や死に敬意を表さなければならない。それらは彼自身で見つけられる何ものよりも大きいからであり、彼の全ての思考がそれらによって呑み込まれるからである。しかし、それらは偉大なものでもある。それらを偉大な思いで考え、その無情な輝きを感じることはさらに偉大である。そしてこのような思考は自由人を作り出す。我々は東洋的な屈服ではなく、それを吸収し、自分自身の一部とする。私的な幸せのために戦うことを放棄し、一時的な欲求を追い出し、永遠のものに対する情熱を持つことこそが解放であり、それが自由人の崇拝である。そしてこの解放は運命自体を征服する思考によって達成される。運命そのものが時間の清め火で浄化されない残滓を去らせる心に服従されることなく、それを吸収し、自分自身の一部とするからである。

自由人は、共通の運命を結ぶ同志たちと強く結ばれ、新たな視点が常に彼の側にあり、愛の光を日々の仕事に注ぐことを発見します。人類の人生は、見えない敵に囲まれた長い夜の旅であり、疲労と苦痛に耐えながら、数少ない者が達成できる目標に向かって進みます。そして誰も長く留まることはできません。一行ずつ進む中で、私たちの同志が見えなくなり、無声の死の命令によって捕らえられます。彼らを助けることができる時間は非常に短く、幸せや不幸が決定されます。私たちは彼らの道に太陽の光を投げかけ、同情の甘露で悲しみを和らげ、絶えることのない愛情による純粋な喜びをとえ、衰えている勇気を強め、希望の時々には信仰を植え付けるべきです。彼らが苦悩し、困難に直面したり、あるいは目指せない場所で失敗した場合でも、私たちは彼らの必要性だけを考えるべきです—悲しみ、困難、おそらくは盲目によって人生の不幸が引き起こされるもの；彼らもまた同じ暗闇の中で苦しんでいる同志であり、私たちと同じ悲劇の舞台に立っている者たちであることを思い出すべきです。そして、彼らの日が終わるとき、彼らの善行と悪事が過去の不滅性によって永遠化したとき、私たちは彼らが苦しみ、失敗した場所で、私たちの行為が原因となったことではないという感覚を持つべきです。しかし、神の火が彼らの心に燃え上がらせた瞬間、私たちは励まし、同情、そして勇敢な言葉で高貴な勇気を輝かす準備ができているはずです。

人間の人生は短く力もありません；彼とその族には無情で暗い運命がゆっくりと降りかかります。善悪に目を閉じ、破壊を恐れず、万能の物質は不動の道を進み続けます。今日、彼は最も愛する者を失うことを告げられ、明日には自身も闇の門を通ることを告げられる人間に対して、残されたのはまだ打撃が降りる前に、その短い人生を栄光に満ちた思想で育てることだけです。運命の怯えを軽蔑し、自分自身の手によって築いた神殿の前で祈りを捧げます。機会の帝国に屈しないで、外なる生活を無駄な暴政から自由な心を守ります。不可避的な力が彼の知識と裁きを一時的に容認する中で、疲れたが屈しないアトラスとして孤独に立ち向かう世界を自分の理想によって形成したものです。

 1903年12月発行の『インディペンデント・レビュー』から再録しました。


\chapter{第IV章：数学の研究 --- {THE STUDY OF MATHEMATICS}}

IV

数学の研究


人間活動に関するすべてについて、時折その目的と理想とは何かという問いを投げかけるべきである。また、それがどのようにして人間の存在の美しさに寄とするかという方法も問うるべきである。生活の機械を提供することで遠くから貢献する追求に関しては、単なる生きる事実よりも、その観察を通じた生活の芸術が望ましいことを思い出されるべきである。なおさら、外部には目的を持たない活動について言えば、それが世界の永久的な財産に実際に加わることでしか正当化されないものについては、その目標を維持し続けることが必要である。創造的想像力を具現化する神殿の明瞭な予想像が見えていなければなりません。

この必要性の充足は、伝統によって若々しい心を鍛えるための科目として決定された学問について特に悲惨に遅れている——その声明すら荒唐無稽であるかのように遠くにある。偉大なる人々は、彼らの生涯を捧げている観察の美しさに十分に敏感であり、他の人々がその喜びを共有することを望むために、次の世代に機械的な知識を伝えることを勧める。この知識なしでは不可能な門を越えるための手段であることを忘れないでほしい。乾いた学究はこの知識を灌漁する権利を得る：彼らはそれが神聖な扉を開く鍵であることを忘れている。これらの聖なる扉に上る階段で一生を過ごしても、彼らは扉自体を見向きもせず、その存在さえ忘れる。熱心な若者はその尖塔と拱廊に導入されるために進むべきなのに、後ろ向きに命じられ、段数を数えるようにさせられる。

数学は、ギリシャやローマの研究よりも、その文明における正当な位置を忘れられた面がより大きいかもしれません。伝統により、少なくとも教育を受けた男性の大部分は数学の基本的な知識を持つことが決められていますが、その伝統が成立した理由は忘れられ、多くの無駄話と微細な事項の下に埋もれてしまいました。数学の目的を尋ねられたとき、一般的な答えは機械の作成や場所への移動、外国との戦争や商取引での勝利を促進することだとされます。しかし、専門家にならない人々に対する単なる基本的な学習がこれらの目標を達成しない場合に、数学は論理的思考能力を訓練するという答えが出されるかもしれません。しかし、この回答を述べる人たちは、既知の誤りであり、素朴な心を持つ聡明な学習者が本能的に拒否する確定的な誤謬を教えることを放棄することにはほとんど抵抗します。論理的思考能力自体は、その開発を主張する人々により、避けるべき罠を避けたり実用生活の指針となるルールを見つける手助けをするための手段と見なされます。これらのすべては数学が持つ明らかに重要な成果ですが、これら全てが数学が自由教育において位置を持つ資格をとえるものではありません。プラトノは、神に値するものとして数学の真理を観察することを認識しました；そしてプラトノは、人間生活における何が天に位置を持つべきかを理解した人物の中で最もその要素を見抜いた人物かもしれません。彼は言います、「数学には何かが必要であり、それを放棄することはできない... そしておそらく神の必要性によるものである。人類の需要について話す多くの人々がこの文脈で用いるような人間の需要については、それほど無意味な解釈はない。クレニアス。そして、これらの知識の需要は何であり、神のものではなく人間のものか？アテネの者よ。人間に真剣に考え、心配することができない何かを知ることなく、世界の神となることや魂となること、英雄となること、人間について真剣に考える能力を持つことはできないものだ。（法律、第818頁）」プラトンはこのような数学の判断を下しましたが、数学教師たちはプラトンを読まず、そのような問題に対する彼の意見を読む人々も数学を学ばず、彼の見解を単なる興味深い誤りだと考えています。

数学は、正しく見れば、真実だけでなく、冷たい厳粛な美しさも持っています——その美しさは彫刻のようなもので、弱い人間の感情を喚起することなく、絵画や音楽の豪華な装飾なしに、純粋で sublime な美と、最高の芸術が示すような厳格な完璧さを持つことができます。最高の喜び、高揚感、人間を超えた感覚は、詩と同じく数学にも見ることができます。最も良い数学は単なる学習以上のものであり、日常的な思考の一部として吸収され、再三再四思い返されるべきです。

現実人生は、ほとんどの人々にとって長い二番目の選択肢であり、理想と可能との間の無数の妥協の中にあるが、純粋な理性の世界では妥協はありませんし、実用的な制限もありません。創造活動は、理想的な追求から生まれる偉大な作品を象徴する壮麗な建築物の中で自由に表現されます。人間の情熱からは遠く離れており、自然の悲惨な事実さえも隔てられている世代が次々と秩序ある宇宙を作り上げ、純粋な思考が本来の故郷である場所で静かに住むことができ、より高い動機の一部が現実世界の厳しい放逐から逃れることができます。

しかし、数学家たちは美を追求するあまりにほとんど何もありません。彼らの作業の中で、多くのものは不可抑制的な本能によって形成され、明示的な信念よりも優れていたが、無意識の味覚によって形づくられました。しかし、多くのものは適切な見解が不適切なものとされ、破壊されました。数学の特性的な優れた点は、論理的に厳格に推理される場所でしか見つけることができません。論理の規則は数学にとって建築構造の規則と同じです。最も美しい作品では、議論の鎖が提示され、その各リンクが独自の価値を持つことから重要であり、全体を通して安らぎと明瞭さがあり、前提がより多くのことが可能であると考えられるよりも多くを達成し、自然で必然的な方法で到達します。

文学は一般的なものを個別の状況の中で表現し、その普遍的な意味が個々の装飾を通じて光り出すようにしていますが、数学は純粋さの中で最も一般的なものを無関係な装飾なしに提示しようと試みます。

数学の教育は、いかにして学習者にできるだけ多くのこの高貴な理想を伝えられるべきか？
ここでは、経験が大いに導きとなるでしょうが、最終的な目的を考慮することでいくつかの格言が得られるかもしれません。

数学を正しく教えれば、その主要な目的の一つは学習者の理由信認を喚起し、証明された真実に対する信頼と、その価値を信じさせる点にある。しかし現在の教育ではこの目的が達成されていない。しかし、それを達成する方法は容易に想像できる。

算数に関しては、生徒や少女には一連のルールをとえられ、これらのルールは真偽ではなく、教師の意志であり、何らかの不明瞭な理由から教師が好むゲームのやり方であると見なされる。このような明確で実用的な学習分野においては、これはおそらく避けられないものかもしれないが、可能な限り早く、これらのルールの理由を子供の心に最も訴える方法で明らかにするべきだ。

幾何学では、ユークリッドの開始部分のように虚偽の証明による冗長な手続きではなく、最初から明白な真理を仮定させ、直感的に驚くべき結論が得られるが、実際に描画によって容易に検証できる定理の証明に教えさせるべきだ。例えば、三つ以上の線が一点で交わることを示すような証明である。

このようにして信頼は生じる。推理が驚くべき結論を導き出すことを見せ、その結果、事実によって検証されることが明らかになると、抽象的または理性的な何かに対する本能的な不信は徐々に克服される。定理が難解であれば、まず幾何学的図形の練習として教え、図形が完全に馴染むまでそれを教授するべきだ。その後、さまざまな線や円の論理的な関係を教えることは喜ばしい進歩となる。

また、定理を説明する図は可能なすべてのケースと形状で描かれるべきであり、幾何学が関わる抽象的な関係が、そのような大きな表面上の多様性の中から類似性として自然に浮き彫りになるようにすべきだ。この方法では、抽象的な証明は授業の一部ではなくなり、具体的な例によって馴染みが得られた後に、視覚的事実の自然な表現として感じられるべきだ。

この初期段階では、証明を学ぶ際には説明的な完全さでとえるべきではない。重ね合わせ法のような確定的に虚偽の方法は最初から厳しく排除され、そのような方法を使わずに証明が非常に困難な場合でも、論理と例によって示されるものと対照的に説明された結果を教えることで受け入れられるようにすべきだ。

代数学の初期では、最も賢明な子供でも通常非常に大きな困難に直面します。文字の使用は、目的が不明瞭で、ただ混乱を増すだけのように思われます。最初は、教師がその文字が表す具体的な数値であることを明らかにしなければならないと信じてしまうことがほとんどです。実際、代数学では、心は一般的な真理を考えるよう教えられます。これらの真理は特定のものだけではなく、一連の多くのものの一つでも通用するものです。このような真理を理解し発見する能力こそが、知性が現実と可能性の世界全体に支配力を行使する根拠であり、一般に対する対処力は数学教育によって最もよくとえられる贈り物の一つです。しかし、通常、代数学から算術への隔たりを説明できる教師はどれほど少ないことか、また学習者は理解のための探求的な努力にどれほど助けられることか。通常、算術で採用された方法が継続されます：ルールが述べられますが、その根拠に対する適切な解釈はありません。生徒は盲目的にルールを使用し学び、不久のうちに教師が望む答えを出すことができると感じ、科目の難しさを克服したと感じるでしょう。しかし、使用された過程の内なる理解を得たことはほとんどないかもしれません。

代数が学ばれれば、すべては順調に進むが、無限の概念を用いた研究、微分積分学や高等数学全般に達すると、問題が生じる。過去に数学的な無限に関する困難性が取り囲んでいたのはおそらく、私たち自身の時代が誇れる最大の成績である。ギリシャ思想の初期からこれらの困難は知られており、毎年、最高の頭脳もゼノンのエレアコスによって提唱された不可解な質問に答えようとして失敗した。ついにゲオルク・カントルがその答えを見出し、無限という新たな広大な領域を知性に征服し、混沌と古い夜から解放した。

無限に関する疑問が存在する集合から何かを取り除いた場合、残されたものの数は元のものよりも少ないという仮定は自明なものとして受け入れられていた。しかし、カントルとデデキントによって反証されると、この仮定は有限な集合に対してのみ成り立つことが明らかになった。無限に関連する場合においてその否定が拒否されたことで、これまで人間の理性を困惑させたすべての困難が解消し、無限に関する正確な科学の創設が可能となった。

この驚くべき事実は、数学教育における革命をもたらすべきである。それは科目自体の教育価値を著しく向上させ、最近まで誤謬と不明瞭さに包まれていた多くの研究を論理的な精度で扱う手段をとえた。旧来の方法で教育を受けた人々には、新しい作業は恐ろしく困難で抽象的であり不明瞭であると見なされるが、発見者はしばしばその光によって霧が晴れるまでまだ曖昧に包まれている。

しかし、無限に関する新理論は、すべての誠実で探求心のある人々にとって、高等数学の掌握を容易にするものである。これまで、論理的に厳格な要求を満たさないと思われる誤った議論を最初から理解し、信認するためには長い巧妙化の過程が必要だったが、現在はそのような議論を一挙に学ぶことができる。

数学教育の過去2世紀間で、無作為に受け入れるべきことが真実であると信じられていた多くの事柄について、理性的な調査によって拒否されるべきものもまた受け入れるべきだと教わるという方法は、恐縮し妥協的な精神や、神秘主義を信じる祭司的な信念が生まれさせた。これらすべては今や取り除かなければならない；数学の奥秘に penetretare 欲する者は、論理的純粋性と関連性そのものによって確立された順序で真の理論を一挙に学ぶべきである。

数学をその目的自体として考え、エンジニアリングの技術的な訓練ではなく、それ自体の美を追求する場合、論理の純粋さと厳密性を保全することは非常に望ましいです。したがって、より簡単な部分に十分な馴染みを持つ者は、自己明了として受け入れた命題からさらに根本的な原理へと逆行し、それ以前には前提と考えられていたものが導き出せるような原理へと進むべきです。彼らは未熟な心で自明と感じられた多くの命題が、精密な検討の結果虚偽であることが示されることを教えられるべきです。この方法によって、最初の原理への疑問的な探求に導かれ、論理の全体構造が築かれる基礎を再調査するようになるでしょう、またはより適切な比喩として、広がる枝から根元が生えている大きな幹を調査するようになるでしょう。この段階では、数学の基本的な部分を再度学ぶことが好ましく、とえられた命題は真偽だけでなく、論理の中心原理からどのように発展したかについても問い直すべきです。このような質問は以前には全く不可能な精度と確実性で今や答えられるようになります；そしてその解答のために必要となる論理の連鎖の中で、ついにすべての数学的研究の統一性が明らかになるでしょう。

数学的教科書の大部分においては、手法に統一性がなく、中心的なテーマの係統的な展開が欠けている。種々異なる種類の命題を、最も容易に理解できる方法で証明し、多くのスペースは主な議論には一切寄としない単なる奇異なものに費やされている。しかし、最高の作品では、ドラマの展開のように統一性と不変性が感じられる；前提からある主題が考慮のために提案され、その後のすべてのステップでその本質的理解へ一定の進歩がなされる。システムへの愛着、相互関連性への愛着は、おそらく知性的な衝動の最も内なる本質であり、数学においては他のどこよりも自由に発揮できる。学習者がこの衝動を感じるなら、意味のない例題の羅列や興味深い奇異なものによる誘引から repellir されず、中心的な原理に留まることを奨励され、提示される様々な主題の構造に馴染み、重要な推論のステップを容易に旅することを奨励されるべきである。このようにして良い精神状態が育てられ、選択的注意は重要で根本的なものに偏って留まることを学ぶことができる。

数学が分野別に分けられたそれぞれの研究を、論理的な全体として見直し、その原理となる命題からなる自然な発展過程として捉えたとき、学習者は統一と体係化された全演繹的推論の基礎科学を理解できるようになる。これが象徴論理である—アリストテレスによって起源を得たものの、広範な発展においては十九世紀にほぼ完全に産出された研究であり、現在でも非常に速いペースで成長している。

象徴論理の真の探求方法、そしておそらく他の数学分野を学んだ学習者にとってこの研究を導入する最良の方法は、実際の演繹的推論の例を分析し、その中で用いられている原理を見出すことである。これらの原理の大半は、私たちの推理本能に深く埋め込まれており、意識的に用いられ、多くの忍耐強い努力によってのみ明らかになる。

しかし、ついにこれらが見つかったときには、数は少ないことが判明し、純粋数学におけるすべての唯一の源泉であることがわかる。全ての数学が少数の基本法則から必然的に導かれるという発見は、全体の知的な美しさを著しく向上させるものであり、大部分の既存の推論鎖が断片的かつ不完全であることに苛まれていた人々には、その発見は啓示のように強烈な力を持つ。旅行者がイタリアの丘を登りながら秋の霧から現れる宮殿のように、数学的な建築物の壮麗な階層が適切な順序と比例で現れ、全体の各部分に新しい完璧さが生まれるようになる。

Until 符号論理が現在の発展に達するまで、数学に依存している原理は常に哲学的であり、これまで哲学者たちが使用してきた不確実で進歩しない方法によってしか発見できないものとされてきた。この考えが存在した限りでは、数学は自立しておらず、自身の方法とは異なる他の研究分野に依存していたように見えた。さらに、算術、解析、幾何学から導き出される公理の性質は伝統的な論理討論の奥深さに包まれており、そのような疑問視できる基礎に建立された建築物は空想上の城のように見られ始めた。この点において、真の原理が数学の一部であり、その結果と同様に重要な部分であるという発見は、学習者の満足感を著しく向上させた。このような満足感は享受できる能力を持つ学習者には拒否されるべきものではなく、人間の能力に対する尊敬と抽象世界の美しさに関する知識を増す種類のものです。

哲学者は一般的に、数学の基礎となる論理法則が思考の法則であり、私たちの心の運作を規制する法則であると信じてきました。この見解は理性の真の尊厳を非常に大きく下げるものです：理性は実在や可能性すべての核心と不変の本質に及ぶ究極的な調査となりますが、代わりに何かより或いは少々人間的で私たちの制限に従属するものについての調査となります。非人間的なものの観察、私たちは創造したわけではない素材と対処できる能力を発見し、特に内面と同様に外側にも美しさがあるという実感は、異邦力による支配下での無力感、弱さ、敵対勢力との放逐の恐ろしい感覚から解放される主要な手段となります。そのような認識を示すことで我々を運命の支配下へ屈服させず、文学的な表現としてこれらの力を人間化した「運命」として描く悲劇の役割は、それ自体が過酷なものです。しかし数学はさらに人間から離れて、絶対的必然性の領域に私を連れて行き、実在する世界だけでなく、どんな可能性の世界も従属しなければならない世界へと導きます；ここでは私たちの理想が完全に満たされ、私たちの最高の希望が妨げられることなく、住まいが築かれます。理性がこの世界を見つけるとき、自分たちの完全な独立性を理解するときのみ、その美しさの深い重要性を十分に認識できるのです。

数学は私たちや私たちの思考とは独立しており、また別の意味で、私たちは存在する万物も数学とは独立している。この純粋に理想的な性格を理解することは、数学が芸術の一環として正しい位置を把握するために不可欠である。かつては純粋な理性が実在の世界について、ある程度の点でその性質を決定できると考えられていた：少なくとも幾何学は私たちが生息する空間に関することだと考えられてきた。しかし現在では、純粋数学は実存に関する問いには決して言及できないことが分かっている：理性の世界は実在の世界を制御するが、実在を創造することはなく、その結果を時間と空間における現実に適用するとき、その確固たる正確性は近似や作業仮説によって失われる。数学で考慮される対象は過去には主に現象から提示されたものであったが、抽象的な想像力はこのような制約から完全に解放されるべきである。したがって相互の自由を許すことが必要となる：理性は実在の世界に対して命令することはできないが、実在の世界は理性の権利を制限することができない。美に対する愛着によって考慮価値があると感じられるどんな対象に対しても、理性は扱うことができる。ここでも他の場面でも、我々自身の理想は世界から見つける断片を組み立てていく：最終的にはそれが創造か発見であるかは明確に言い切れない。

指示に従い、以下の英文を忠実に日本語へ翻訳します：

数学教育においては、重要な定理の正確性だけでなく、その証明方法自体が最も美しく、すべての可能な方法の中で最も美しいと学生に説得させることが非常に望ましいです。証明の真の興味は、従来の解釈模式が示すように、結果に集中しているわけではありません；これが起こった場合、それは可能な限り補正されるべき欠点であり、証明の手順をより一般的なものにすることで各ステップが自身で重要であるようにする必要があります。結論を立証するだけの議論は、ある道徳的な教訓を教え込むために下位にある物語のように扱われます：全体のどの部分も手段であるべきではありません。実用的な精神や新しい領域を征服するために急速に進展することへの欲求が、数学教育において結果に対する過度な強調を生み出しています。より良い方法は、考慮すべきテーマを提案することです—幾何学では重要な性質を持つ図形、解析では洞察をとえる関数などです。証明が研究対象を定義するためのマークの一部に依存している場合、これらのマークは自身で調査されるべきです。なぜなら、結論には不要な前提を使用することは議論における欠点であり、数学では新しい公理が導入されると結果の定理の一般性が低下しますので、最大限の一般性を求めることは最も重要です。ユークリッドの功績は、並行線の公理を使わずにできる限り進むことができたことです—これはこの公理が本来的に問題があるからではなく、新しい公理が導入されると定理の一般性が低下するためです。最大限の一般性を求めることが最も重要であることを示しています。

数学の影響は、その自己の範囲を超えた分野においても、自身の理想的な範囲について書かれたものよりも多くが記述されている。過去には哲学に対する影響が最も目立っていたが、多種多様なものでもあった；17世紀には理想主義と理性主義、18世紀には実在論と感覚主義が数学の子孫のように見えた。将来に於ける数学による影響については、多くを言うのは非常に冒険的であるが、一つの点では良い結果が期待できるようだ。疑問主義者に対する数学の反論は、その道が困難で目標が確実には達成できないため理想追求を放棄するような疑問主義に対し、自範囲において完全な答えとなる。しばしば言及されるように絶対的な真実は存在しないとされ、意見や私見のみがある、それぞれの視点は個々の特徴、好み、偏見によって形成されている、確実な真理の外なる王国は忍耐と規律によって最終的に訪れるものではなく、ただ私ごとの中だけにある真実があるというような考えが示される。このような疑問主義は人間努力の主要な目的を否定し、勇敢な現実への認知である誠実性が我々の道徳的な視野から消え去る。しかし数学はその真実の建築物が、疑問主義的 cynicism のすべての武器に対して揺るがないと破壊不能であることを常に戒める。

数学が実際の生活にとえる影響は、私たちの研究の動機とすべきではないものの、孤立した学習者が常に抱く疑問に対する答えとなることができる。悪や苦しみで満ちたこの世界において、瞑想の隠れ家に退屈し、いかなるものでも少数しか享受できない尊い喜びを楽しむことによって引き篭もるのではなく、他人が運命的な出来事によって負担されている苦しみを助けることに拒否する自意識過剰な行為には、必ずしもそうであるとは思えない。私たちは、私たち自身が困難で厳格な生活を送りながら、現在の悪から逃避して、同僚の人々に援助をとえることなく生きる権利があるのかと問うことができるだろうか？

これらの疑問が生じたとき、真実の答えは、聖なる火を息づかせる者や、常にそのような努力の果てに影を落としている haunting vision を維持する者が一定数必要であるということであろう。しかし、時折起こる場合があるように、この答えが冷たく感じられるとき、私たちが助けることができない苦しみの光景を見守り狂うような状況が生じたときには、数学はその他の実用的な活動よりも多くの人々に幸福をもたらすことがしばしばあることを考えることができる。科学の歴史は十分に証明しているように、抽象的な命題の集合体（双曲線の場合でも二千年間日常生活に影響をとえないとしても）がいつかは市民の日常的な思考や活動を根本的に変革する手段となる可能性がある。

蒸気や電気の利用など、著しい例として挙げられるものは、数学なしでは不可能である。抽象的な考えの結果は、世界が所有している資産であり、その使用によって一般的な生活を豊かにする効果はこれまでに発見された限界がない。また、どの部分の数学が有用となるかについては経験からは決定する方法がないため、有用性は失望の時における慰めであるだけで、学習を導く指針にはならない。

数学の学問において、道徳的生活の健康を保ち、時代や民族の品位を高めるためには、厳格な德性が奇妙な力を持っています。その力は、思考によって啓蒙され浄化された者だけにとえられるものを超えています。これらの厳格な德性の中で、真実への愛が最も重要であり、数学においては他のどの分野よりも、真実への愛は衰えた信仰を慰める機会を見つけることができます。偉大な学問は、自分自身の目的だけでなく、高い精神習慣を創造し維持する手段でもあります。そしてこの目的は、数学の教えと学びを通じて常に目指すべきものとして保たれなければなりません。


\chapter{第V章：数学と形而上学者たち --- {MATHEMATICS AND THE METAPHYSICIANS}}

数学と形而上学家たち


十九世紀は蒸気の発明や進化を自慢としたが、純粋数学の発見からもっと正当な名誉を得られたかもしれない。この科学は他の多くの学問と同じく生まれる前に洗礼を受けた；そして十九世紀以前の著者は、彼らが呼んだ純粋数学について言及していることがわかる。しかし、もし彼らにその科目は何であるかと尋ねられたなら、 Arithmetic（算術）、Algebra（代数）、Geometry（幾何学）などといった内容しか説明できなかっただろう。これらの研究が共有するものは何か、そしてそれらが応用数学とはどのように異なるのかについて、祖先たちは完全に見当違いだった。

Booleによって発見された純粋数学は、彼自身が『思考の法則』（1854年）と呼んだ著作の中で明らかになった。この著作には、それが数学ではないという断言が溢れており、実際にはBooleは自分の本が初めての数学に関するものであるとは思っていなかったことを示している。また、彼は自分が思考の法則に取り組んでいると考えていたが、人々が実際にどのように考えているかという問題は彼にとって関連性がなく、もし本当にそのような法則が含まれていたなら、なぜ以前からそんな風に考えていなかったのか不思議だった。実際には彼の本は形式論理に関心を持っていたものであり、これは数学と同じことである。

純粋数学は、何かが特定の命題で真である場合、その何かについて別の命題も真であるという主張のみを含む完全に由来するものである。最初の命題が本当に真かどうかを議論したり、それが何かについて述べられているかどうかを指摘することは絶対的に必要ではない。これらの点は応用数学に関連するものである。我々は純粋数学において、推論規則から始め、その規則により一つの命題が真である場合に別の命題も真であることを導くことができる。これらの推論規則は形式論理の主要な原理を構成している。

その後、何かが amusement として見受けられる任意の仮定から始め、その結果を導き出す。もし仮定が何かについてではなく、特定の一つまたは複数のものについてである場合でも、我々の推論は数学となる。したがって、数学は我々が何について話しているのか、また我々が述べていることが真偽如何に関わらず知る必要がないという意味で定義できる。数学の始まりに困惑した人々にとって、この定義は慰めとなり、おそらくその正確性を認めるだろう。

現代数学の主要な勝利の一つは、数学そのものとは何かを発見したことにあり、この話題について少し追加の言及を行うことも無駄ではない。通常は、幾何学などのどの数学の分野でも、定義不能とされる原始的な概念や証明不能とされる原始的命題や公理から始める。しかし実際には、応用数学においても必ずしも全てが定義不能または証明不能であるわけではないが、純粋数学においてはそれらは一般論理に属するもの以外存在しない。論理学は広い意味で、その命題が何でも適用できる形で表されることが特徴とされる。すべての純粋数学、すなわち算術、解析学、幾何学は論理の原始的な概念による組み合わせによって構築され、その命題は論理の一般公理、つまり三段論法や推論の他の規則から導き出される。そしてこれは夢や希望ではなく現実となっている。むしろ数学のより大きな難易度の高い領域においては既に達成されているし、残された少数の例でも特別な困難はなく、現在急速に達成されている。哲学者は長い間、このような推論が可能かどうかについて議論して来たが、数学者はその推論を実行した。そして今では哲学者の役割は礼賛するだけとなった。

形式論理の科目は、ついに数学と同一であることを示し、その起源はアリストテレスにまで遡る。アリストテレスによって創始され、中世では神学以外の主要な研究分野となった形式論理だが、アリストテレスは結論論理までしか進展せず、学校哲学者たちもアリストテレスの範囲内で止まっていた。中世医師たちに対する私たちの優越性を証明する必要があれば、この点を見ればよいだろう。中世を通じて、ほとんど全ての最上級の知性が形式論理に捧げられていた一方で、19世紀には世界全体の思考の中でもその科目に入るものは僅少だった。

しかし、1850年以来、毎十年ごとにアリストテレスからレーニッツまでの全期間よりも多くの進歩が見られた。人々は論理を代数のように符号化し、数学的な規則によって推論を行う方法を見出した。さらに結論論理以外にも多くのルールが発見され、旧来の論理では対処不可能だった主題に対応する新しい論理分野、関係論理が創始された。これは数学の中心的な内容を完全に超えるテーマに対処するために作られたものであり、その多くは旧来の論理では扱うことができなかった。

symbolismは、数学の基礎を論じる際に重要であるが、一般の人々にはその重要性を理解するのは容易でない。そして、説明は奇妙なパラドックスのように思われるかもしれない。実際、 symbolismは有用であり、それは物事を難しくするからである。（これは高度な数学の部分に関しては真ではないが、基礎的な部分に関してのみ真である）我々が知りたいことは、何から何を導き出すことができるかである。しかし、基礎的な部分では、すべてが自明である；そして一つの自明な命題が別のものから導き出されるかどうかを見分けるのは非常に難しい。自明さは常に正確性の敵である。したがって、新しいかつ困難な symbolismを発明し、何も自明でないようにする。その後、符号に対する特定の操作ルールを設定し、全体として機械的に処理できるようになる。この方法で、前提条件は何であり、何が導き出されるかまたは定義されるかを見つけることができる。例えば、算術と代数学は三つの未定義概念と五つの未証明命題を必要とすることが示されている。しかし、 symbolismがない場合、これを見つけ出すのは非常に困難だった。二加二が四であることの自明さは如此して、自分自身に十分な疑念を持てずにそれを証明できるかどうかを疑うのが難しい。そして他の自明なことを証明する場合も同様である。

しかし、自明な命題の証明は、未経験者には少し軽薄な仕事のように思われることがあります。これに対して答えるには、しばしば一つの自明な命題が別の自明な命題から導かれることが自明でないことがあり、従って、明らかであることを明らかにすることを証明する方法が明らかでないという手法によって新しい真実を見出しているということです。しかし、もっと興味深い反論は、人々が明らかであることを証明しようと試みた結果、多くの命題が偽であることが判明したことです。自明さはしばしば単なる霧の中の灯火で、それを指南者としすぎると間違いに導かれることがあります。例えば、全体には部分よりも多くの項があることや、数は一つを加えることで増加することなどは明らかです。しかしこれらの命題は現在では通常偽であることが判明しています。ほとんどの数は無限であり、無限の数に対してはいくらでも一つを加えても全く影響を受けません。証明の一つの長所は、その結果に対する一定の疑念を植え付けることです。明らかであることを一部では証明できても他ではできない場合、それらの場合でそれが偽である可能性が生じます。

arty事理の偉大な専門家は、私たちの時代の人々の中でもイタリア人教授ペアノです。トーリン大学に在籍しています。彼は数学の大部分（そして時間とともに全体も）を厳格な符号表現にまで簡素化しました。通常の数学書では、言葉は読者にとって必要以上に少ないかもしれません。それでも、「従って」「仮定してみましょう」「考慮する」「したがって」といった短いフレーズが現れます。しかしこれら全ては妥協であり、ペアノ教授によって取り除かれます。例えば、算術、代数学、微分積分学、そして通常の純粋数学（幾何学を除く）の全般的な理解を学ぶためには、三つの単語で構成される辞書から始めなければなりません。一つの符号は「零」を表し、もう一つは「数」を表し、さらに一つは「次に来るもの」を表します。これらの概念が何を意味するかを知る必要があるのは、数学家として始めるためです。しかし、これら三つの概念の符号が作成されると、その後の全ての符号はこれら三つによって説明されます。さらには、これら三つも「関係」や「クラス」という概念で説明できるかもしれませんが、これは関係論理が必要とし、ペアノ教授はこれを取り上げていません。最初に数学家として知るべきことは多くはありません。全ての純粋数学（幾何学を含む）における概念は至多十二つの基本概念から成り立っています。ペアノ教授は非常に能力のある若きイタリア人の弟子たちの支援を得て、これらがどのように組み合わさるかを示しました。彼が発明した方法はさらに進められる余地がありますが、先駆者的な功績は彼に属するべきです。

二百年前、レーニッツはペアノが完成させた科学を予見し、それを創造しようと試みました。彼はアリストテレスの権威に対する敬意から成功できなかったが、彼が望んでいたものは現在存在しています。その「普遍的特性」について彼は言いました：「争議が起こった場合、二つの哲学者の間でも二つの会計士の間でも、論争をする必要はありません。計算機を取り出し、机に座って互いに（友人が証人であれば）『計算しましょう』と言えばよいのです。」

この楽観主義は現在では少し過度なものと見なされるようになりました。まだ解決が困難な問題や、計算によって決着をつけられない争議もあります。しかし、以前は論争の中心だった広大な領域において、レーニッツの夢は現実となりました。数学哲学全体についても、かつては他のどの部分よりも疑問に満ちていた哲学の一部分でも、混乱と迷いが秩序と確実性に取って代わりました。当然、哲学者たちはこの事実をまだ発見しておらず、旧来の方法でこのような主題について書くことを続けます。しかし、少なくともイタリアの数学者は、ペアノやゲオルク・カントルの著作を通じて数学の原理を正確かつ巧みに扱う力を持ち、これにより数学の確実性は数学哲学にも及びました。そのため、以前は大きな謎とされた多くのトピック—例えば無限、連続性、空間、時間、運動の本質など—現在ではどの程度でも疑問や議論の対象ではありません。これらの事物の性質を知りたい人は、ペアノやゲオルク・カントルのような人物の著作を読むだけで、これら昔の謎が正確で疑いの余地のない説明で提供されています。

この世は不思議なもので、もっと不思議なものは後世の評価だ。その最も顕著な例として、エレティアのゼノンを挙げることができる。この人間は無限の哲学の創始者と見做され、プラトンの『パルメニデス』においてソクラテスの教師として優遇されている。彼は４つの議論を考案し、これらはすべて無比の精巧さと深さを持ち、運動が不可能であることを証明し、アッキレスが亀を追いつくことはあり得ないことを示し、飛行中の矢は実際には静止していることを主張した。アリストテレスやその後の哲学者たちによって駁難され、その議論は再評価され、無限小を使用しない数学の再興の基礎となったドイツの教授によって復活した。ウィーラートルスは、無限小を排除して数学から infinitesimals の使用を厳格に禁じることで、我々が不変の世界に生きていることを示し、飛行中の矢は実際には静止していることを証明した。ゼノンの唯一の誤りは（もし彼が推論したとすれば）、変化の状態が存在しないという事実から、世界はどの時間でも同じ状態であるという結論を導いたことだった。これは必ずしも成り立つわけではなく、この点ではドイツの数学者の方が建設的だと言える。ウィーラートルスは彼の見解を数学に具体化することで、真実に対する親しみが庶民的な偏見を排除し、ゼノンのパラドックスに尊敬すべき平凡さをとえられた。そしてその結果は、理を愛する者には少々魅力的ではないかもしれないが、少なくとも学術界の大衆に対して平和をもたらすことに成功したと言えるだろう。

ゼノは実際、運動によって提示された3つの問題に取り組んでいましたが、それぞれが運動よりも抽象的であり、純粋に算術的な扱いが可能な問題でした。これらの問題は無限小、無限、連続性の問題です。これらの困難を明確に述べることは、哲学者の任務の中でもおそらく最も難しい部分を達成することでした。ゼノによってそれがなされました。ゼノから現代まで、各世代の最高の知性が順番にこれらの問題に取り組みましたが、概して何も達成できませんでした。しかし、現代において、3人の人物—ワイルシュタイン、デデキント、カントール—はこれらの3つの問題を進展させただけでなく、完全に解決しました。数学に通じている人にとっては、その解は明確で、今後も最小限の疑問や困難が残らないほどです。この達成はおそらく当時の時代が誇れる最大のものであり、私は他の時代（ギリシャの黄金時代を除く）よりも、その偉大な人物たちの超人的な才能をより説得力のある証拠として提示できる時代を見つけることができません。この3つの問題の中で、無限小に関する問題はワイルシュタインによって解決されました；他の2つの問題の解法はデデキントによって始まり、カントールによって決定的に完成しました。


古代において無限小は数学において重要な役割を果たしていました。ギリシャ人は円を非常に多くの非常に小さな等しい辺を持つ多角形から微細に異なるものとして捉えていました。無限小は徐々に重要性を増し、ライプニッツが無限小解析を発明したときには、より高度な数学の基礎となるべき概念と考えられました。キャリーは『フリードリヒ大王』の中で、ライプニッツが普魯士女王ソフィア・シャルロッテに無限小について語ったことと、彼女がそのような話題については教育を必要としないと言ったことを語っています。しかし、哲学者や数学者は、宮廷の知識よりも多くそれを議論しましたが、進展は見られませんでした。解析学には連続性が必要であり、連続性には無限小が要求されると考えられていましたが、誰も無限小とは何かを発見することができませんでした。明らかにゼロではなかったことは確認されました。十分な数の無限小を加えると有限な全体が形成されることを見ました。しかし、誰も非零で有限ではない分数を指摘することはできませんでした。したがって、停滞状態が生じました。しかし、最後にはワイエルシュトラスが無限小は必要ないことを発見し、それなしでもすべてのことが可能であることを明らかにしました。したがって、そのようなものは存在する必要がなくなりました。現在では数学者はライプニッツより礼儀正しく、無限小についてではなく無限大について話しています。しかし、王侯にとって適切なこの主題は彼らにはあまり興味を引かないようです。

The 無限小の放逐には様々な奇妙な結果があり、これらは徐々に慣れなければなりません。例えば、次の瞬間というものは存在しません。一つの瞬間と次の瞬間との間に存在する区間は無限小であるべきです。なぜなら、有限の区間に二つの瞬間を選んだとしても、その間に他の瞬間が常に存在します。したがって、無限小が存在しない場合、どの二つの瞬間も完全に連続してはおらず、どのような二つを選んでも中には他の瞬間が存在します。そのため、任意の二つの瞬間間に無数の瞬間があるべきです。なぜなら、有限の数だけの瞬間があったら、その中の最も近い瞬間が次の瞬間となるからです。

これについては難題と感じるかもしれません。しかし実際には、無限に関する哲学がここに登場し、すべてを解決します。

空間においても同じことが起こります。任意の物体を二つに切って、さらに各部分を半分に切り裂き続けると、その断片は小さくなり続け、理論上は何度でも小さくすることができます。どの程度小さいとしても、それらはまだ更に小さくできるという可能性があります。しかし、どんなに小さな物体であっても、必ず有限のサイズを持っています。無限小に達することはなく、有限回だけ分割しても点には到達しません。それでも点は存在しますが、これらは連続的な分割によって到達するものではありません。またここでも、無限に関する哲学がこの可能性を示し、なぜ点が無限小の長さではないかを説明します。

運動や変化に関しては、類似の奇妙な結果が得られます。人々は昔から、何かが変化すると必ずその過程において変化していると考え、何かが動くと必ずその過程において動いていると信じていました。これは現在では誤りであることが判明しました。

物体が動くときには、ただその物体が一定の時間に一定の場所にあり、別の時間には別の場所にあるというだけが言えるのです。次の瞬間には隣接する場所にいるとは言えません。なぜなら、次の瞬間は存在しないからです。

哲学者たちはしばしば、物体が動いているときにはその瞬間に位置が変化すると主張します。しかし、ゼノンは古くからこの見解に対して致命的な反論をしました。「物体は常にその場所に存在している」という点で、哲学者の主張は否定されました。

しかし、哲学者たちにとって単純かつ短い反論は重視されず、彼らは今日まで同じ言葉を繰り返し続けました。ゼノンの破壊的な熱意を掻き立てるために。

最近になって初めて、ゼノンの簡潔な教訓に従って運動を詳細に説明する方法が見つかり、哲学者のパラドックスと対立しました。今や物体が動いているときも静止しているときと同じようにその位置にあるという安心した信念を持つことが可能になりました。

運動はただ物体が時々一定の場所にあり、時々別の場所にいるという事実に過ぎません。また、中間の時間には中間の場所にいることもあります。この単純で明快な常識を理解するためには、この主題に関する哲学的な推測をくぐり抜ける必要があります。そして、古来的な偏見から解放される喜びは、そのような簡潔な教訓の中に含まれているのです。


 infinitesimal の哲学については、我々が見たように、主に否定的なものです。人々は一度その存在を信じていたが、現在ではその誤りを見つけることに成功しました。一方で、無限の哲学は完全に肯定的です。かつて無限の数や数学的な無限が自己矛盾であると信じられていました。しかし、数には無限があることが明らかだったため、無限の自相矛盾は避けられないと考えられ、哲学は「行き止まり」に陥る可能性がありました。この困難は Kant の対立原理へと導き、それから多少間接的に Hegel の弁証法方法へとつながりました。現在のほとんどすべての哲学は、非常に少ない哲学者たちを除いてまだ認識していないが、無限に関する古い且つ信頼性のある矛盾は一挙に解決されたという事実によって揺さぶられています。この問題を解決した手法は最も興味深く、教益的なものです。まず第一に、ギリシャの思考が始まった以来、人々は流暢に無限について話していましたが、誰も一度も「無限とは何か？」と尋ねたことがありませんでした。もし哲学者が無限の定義を求められた場合、彼はおそらく意味不明な冗談を述べるかもしれませんが、全く意味のある定義をとえることは不可能でした。20年前ごろから、Dedekind と Cantor がこの質問を提起し、さらに驚くべきことにその答えを見つけました。つまり、無限の数や無限の集合について完全に正確な定義を見つけることができたのです。これが最初でおそらく最も重要なステップでした。次に、この概念における仮想的な矛盾を検討する必要がありました。ここで Cantor は唯一適切な方法で進みました。彼は対立する命題のペアを取り、通常両方の側が証明されるものと見做されていましたが、その仮説的証明を厳格に検討しました。彼は無限に対する反対の証明が特定の原則を含んでいたことを発見しました。この原則は最初に見えていたly真実であるように思われましたが、その結果としてほとんどすべての数学が破壊されるという悪影響がありました。一方で、無限を支持する証明にはこのような悪影響をとえる原理はありませんでした。したがって、常識が説得力のある最大公約数に惑わされたことが明らかになり、この原則を拒否すると全てがうまくいくことがわかりました。

問題の格言は、一つの集合が別の集合の一部である場合、その部分集合の方が項数が少ないという点です。この格言は有限な数に対して真です。例えば、イギリス人はヨーロッパ人の一部であり、イギリス人がヨーロッパ人よりも少ない数があります。しかし無限の数に到達すると、これはもう成り立ちません。この格言の崩壊が無限の正確な定義をとえます。項からなる集合が無限であるとは、その集合の中に、項数が同じ他の集合が含まれていることを指します。もし集合からいくつかの項を取り除いても項数は減らない場合、その集合には無限の項があります。例えば、偶数の数と全体の数は同じ数であり、なぜならすべての数を2倍にすれば偶数になります。これは以下のように奇数と偶数を一列に並べて確認できます：--

          1, 2, 3, 4, 5, \textit{ad infinitum}.
          2, 4, 6, 8, 10, \textit{ad infinitum}.

下の行にある項の数は上の行と同じであることが明らかで、それぞれに対応する関係があります。この性質は以前には矛盾と見做されていましたが、現在では無限の定義として無害化され、上記の場合では有限な数の数が無限であることを示しています。

だが、未経験者は無限の数を一つずつ数えることが可能なのか疑問に思うかもしれない。すべての数を一つずつ数えることは不可能である。なぜなら、数えた数だけ数があると感じるときでも、さらに続く数が常に存在するからだ。実際、数えることは、集合の中に何個の項があるかを知る非常に俗野で基本的な方法に過ぎない。そして結局のところ、数えることは数学的に「序数」の数を私達にとえ、それは項を順序や係列に並べ替えることによって得られ、その結果はこの並び方からどのような係列が生じたかを教えてくれる。つまり、何かを数えるときには最初に数えるものとその後に数えるものがあるため、数えることは常に順序に関連している。有限の項しかない場合、それらを好きな順序で数えることができるが、無限の項しかない場合、その操作の方法によって得られる結果は全く異なるものになる。したがって、一般的な意味での数え方から生じる「序数」の数は、項の数だけではなく（項の数が無限の場合）、項がどのように並べられているかにも依存する。

基本的な無限数は順序数ではなく、 CARDINAL と呼ばれるものである。これらは項を順に並べて数えることによって得られるものではないが、異なる方法で得られるものであり、その方法により最初から二つの集合が同じ数の項を持っているかどうか、もしくはどちらの集合の方が大きいかを知ることができる。この方法は数を数えるようにして項の総数を教えてくれるわけではないが、ある集合の項の総数を特定の集合の項の総数として定義する場合、この方法により他の参照された集合がより多くまたは少ない項を持っているかどうかを発見できる。

例を挙げてこの方法がどのように機能するかを見ると、その国で人口調査が不可能でも、男性は全て妻を持ち、女性は全て夫を持つということが分かっている場合（多妻制が国家の制度ではないと仮定すると）、無数に数えることなく、その国の男性と女性の総数は同じであることを知ることができる。この方法は一般的に適用できる。ある関係がある二つの集合のそれぞれの要素を他方の集合の一つの要素と一対一に対応付けるもので、逆もまた同様である場合、その二つの集合は同じ数の項を持っている。

例えば、偶数と自然数の総数が等しいことを示す方法として、全ての数は倍になり、全ての偶数は半分になるという関係がある。それぞれの過程により一つの数に対応する唯一の数が得られる。このようにして無限の集合の何れもが有限の数と同じ数の項を持っているような複数の集合を見つけることができる。

ある集合の各項を数に結びつけ、全ての有限の数が一度だけ使用される場合、その集合は有限の数と同じ数の項を持っていることがわかる。この一般的方法により無限集合の数が定義される。

但すべての無限数が等しいとは思われてはなりません。むしろ、無限数よりも有限数の方が極めて少ないのです。異なる種類の数列に有限数を配置する方法の数は、有限数よりも多くあります。空間中の点や時間内の瞬間の数も、おそらく有限数よりも多いでしょう。整数と分数の数と同様に、分数の数も無限にありますが、任意の2つの整数間には無数の分数があります。しかし、整数や分数よりも無理数の方が多く存在します。空間中の点の数はおそらく無理数と同じであり、百万分の一インチ長の線にある点の数も無限空間全体にある点の数と同様です。最も大きな無限数が存在し、それがすべてのもの全て、種類や形を問わず、ある数であることは明白です。これ以上の数は存在しないでしょう。なぜなら、すべてを取り尽くした後には追加できるものが何も残っていないからです。カントールは最大の数が存在しないことを証明しており、この証明が有効であれば、無限に関する矛盾が sublime な形で再出現することになります。しかし、この一点において、彼は非常に巧妙な誤りを犯している可能性があります。私は将来的にその点について説明する予定です。

我達現在理解了芝諾為何認為阿基里斯無法追上烏亀，同時也明白了実際上他可以追上的原因。我達将看到，所有と芝諾観点相左的人其実都没有理由これ様做，因為他達都接受了導致芝諾結論的假設前提。

論理は以下の通りです：アキレスとウッキーが同時に道路を進むことにします。公平のため、ウッキーには差し引きをとえます。アキレスがウッキーの二倍速で走る、あるいは十倍速、百倍速で走ることにします。すると、彼は決してウッキーを追いつかないでしょう。なぜなら、どの瞬間でもウッキーはどこかにいて、アキレスもどこかにいるからです；そしてどちらも競技が進行中では二度と同じ場所には存在しません。したがって、ウッキーはアキレスと同じ数の場所に行きます。なぜなら、それぞれは一つの瞬間に一つの場所にあり、他の瞬間には別の場所にいるからです。しかし、もしアキレスがウッキーを追いつくことができた場合、ウッキーがいた場所はアキレスがいた場所の一部しかありません。ここで、芝諾は全体よりも部分の方がより多くの要素を持つという原理に訴えました。つまり、もしアキレスがウッキーを追いつけるなら、彼はウッキーよりも多くの場所に存在したはずです。しかし、私たちはどの期間でも彼とウッキーは同じ数の場所にいたことを確認しました。従って、彼は決してウッキーを追いつかないという結論を導きます。

この議論は、全体が部分よりもより多くの要素を持つという公理を許す場合に厳密に正しいです。結論が荒唐無稽であるため、その公理を否定し、それによって全てがうまくいくはずです。しかし、過去二千年以上にわたる哲学者たちには良い言葉はなく、彼らはこの公理を認めつつ結論を否定してきました。

この公理を保持すると絶対的な矛盾が生じますが、それを否認すると奇異な結果しか出てきません。これらの奇異のうちいくつかは非常に不思議なものでした。その一つ、私はトリストラム・シャンディのパラドックスと呼ぶものがあり、アッキレスの逆数であり、その結果トートリオが時間をとえられればアッキレスと同じ距離を進むことを示しています。

トリストラム・シャンディについては、彼は最初の二日間について二年を費やし、このペースでいくと記録すべき資料が彼が処理できるよりも速く増えるため、年月が経つにつれて歴史の終わりに更に遠ざかると嘆いていました。しかし私は、彼が永遠に生きていて、その仕事から疲れることなく続けた場合、彼の生涯は最初と同じように進行したとしても、彼の伝記のどの部分も未書かれることはありません。なぜなら、百日目は百年目に、千日目は千年目に記述されるためです。どんな日にでも到達不可能だと考えられるものを選んでも、その日は対応する年に記述されます。したがって、どのような日を指しても、いずれかの時点で書かれることになりますから、伝記のどの部分も永久に未書かれることはありません。このパラドックスでありながら完全に真実な命題は、すべての時間における日の数が年数と等しいという事実に基づいています。

無限に関する結論は最初見たときには不思議そうですが、その理由は多くの哲学者が無限に固有の矛盾があると信じていたからです。しかし少しの練習でカントールの教義の真の原理を理解し、真偽について新しいそしてより良い直感を持つことができます。するとこれらの奇異は、頭を逆さまにして立つことを好まない反対極の人々よりもさらに不思議ではありません。

無限に関する問題の解決により、カントールは連続性に関する問題も解くことができました。彼はこの概念について完全に正確な定義をとえ、その概念において矛盾がないことを示しましたが、この主題は技術的なもので、ここではその概要を提供することは不可能です。


連続性の概念は順序に依存しており、連続性は単なる一種の順序であるからです。現代数学では、順序がますます重要な位置を占めています。かつては（そして今でも哲学者たちはしばしばそう信じているが）、量が数学の基本的な概念だと考えられていましたが、現在では量は幾何学の一部を除いて完全に排 displaces され、順序が支配的になっています。数列の種類とその関係性に関する調査は現代数学において非常に大きな部分を占めていますし、この調査は量に関連することなく、大半の場合では数に関しても関わらずに行われることが分かりました。すべての数列は形式的な定義が可能であり、それらの性質は符号論の原理に基づいて relatives の代数によって導き出すことができます。より高度な数学の大半において基礎となる極限の概念も、かつては量を用いて定義されていましたが、近づくことができる任意に近い項という意味で定義されていました。しかし現在では極限は全く異なる方法で定義され、その制限する数列は必ずしもそれを近づいているとは限りません。この改善もカントールの功績であり、数学を革命的に変えるものでした。極限に関連するのは順序のみです。例えば無限整数の最小値は有限整数の数列の極限ですが、すべての有限整数はその最小値から無限遠くにあります。異なる種類の数列に関する研究は一般的な主題であり、順序数（先述）の研究はその特殊で非常に興味深い分野です。しかしながらこの主題の避けられない技術的な側面により、それらを非専門家には説明することは不可能です。

幾何学は、算術と同様に、最近では順序の一般的な研究の下に収斂されている。かつては幾何学が我々が住む空間の性質を研究するものだと考えられていたが、実在のみが経験的にしか認識できないとする人々にとっては、幾何学は実用数学の一部として扱われるべきであると主張された。しかし、非ユークリッド係の増加により、幾何学が空間の性質に比肩するほど光を投げかけないことが徐々に明らかになった。算術が米国の人口に比肩するほど明確な洞察をとえるわけではなく同様である。

幾何学は、対応する公理の集合に基づく一連の推論的科学の集まりである。一つの公理の集合はユークリッドのものであり、他の同等に良い公理の集合も存在し、それぞれ異なる結果を導き出す。ユークリッドの公理が真であるかどうかという問いは、純粋数学の研究者にとっては関心事ではなく、実質的に確実な肯定的回答を得ることは理論上不可能であることもまた同様である。非常に精密な測定によってユークリッドの公理が偽であることを示すことは可能かもしれないが、観測上の誤差によりどの程度正確であるかを絶対に確認することはできない。

したがって、幾何学者は科学者に任せるべきで、できるだけ適切に最も近い真実性を持つ公理を選択し決定させる。幾何学者は興味深い任意の一組の公理から推論を行い、その結果を導き出す。この意味での幾何学が定義されるのは、公理が1次元を超える係列を生むことである。そして、これにより幾何学は順序の研究における一分野となる。

幾何学においては、数学の他の部分分野と同様に、ペアノとその弟子たちは原則に関する最高の価値を持つ業績を残しました。かつては哲学者や数学者ともに、幾何学の証明が図形に依存すると考えられていましたが、現在ではこれは誤りであることが知られています。最良の教科書には図形は一切含まれていません。論理は最初から定められた公理に基づいて形式的論理の厳格な規則によって進行します。図形を使用する場合でも、あらゆる種類のことが直感的に明らかに思えるものがあり、これらは明示的な公理から形式的な論理で証明することはできませんし、実際には直感的に受け入れられるだけです。図形を排除することで、必要なすべての公理を見つけることが可能になります。これにより、それらが以前は検出されなかった可能性も明らかになります。

正確性に関して大きな進歩として、点を必要に応じて導入することによって、全体の空間を前提とする昔の方法から脱却しました。この手法はペアノだけでなく、イタリア人のファノによっても発展させられました。未経験者にとっては少し強情な教条主義的であるように感じられます。このような方法では、以下の公理から始めます：（1）存在する一類のエンティティを「点」と呼びます。（2）少なくとも一つの点が存在します。（3）点 $a$ がある場合、$a$ 以外の少なくとも一つの点が存在します。次に、二つの点を結ぶ直線を導入し、再度（4）について開始します：直線 $ab$ 上には、$a$ と $b$ 以外の少なくとも一つの点が存在します。（5）直線 $ab$ 上の点でない少なくとも一つの点が存在します。そしてこのように続けます。必要な数だけの点を得る手段を持ちます。しかし、ペアノは冗談交じりに言及しているように、「空間」という言葉は幾何学には全く必要のないものであると言います。

現代の幾何学者が用いる堅実な手法は、ユークリッドを正確性の頂点から引き下ろしました。直到最近，人達還認為，如亨利·薩維尔爵士在1621年所指出的那様，欧几里得只有両个瑕疵：平行線理論和比例理論。現在我達知道，在これ些之外，几乎没有任何其他地方存在瑕疵。他的前八个命題中充満了無数錯誤。也就是説，並非僅僅関于公理的真実性存疑——これ相対来説是較為次要的問題——而是肯定地説，他的命題並不从他所提出的公理推導出来。証明其命題所需的公理数量遠遠超過欧几里得無意識地使用的数量。即使在他的第一个命題中，即在給定底辺的基礎上構造等辺三角形時，他也使用了両个假設相交的円。但没有任何明確的公理保証これ一点，而在某些類型的空間中它達並不総是相交。我達的空間是否属于これ些類型之一尚存疑問。因此，在欧几里得的第一个命題中他完全未能証明其論点。由于他肯定不是一位易于理解的作者，並且極其冗長，他已経不再具有実際的教学价値。在これ種情况下，仍然在英国的男孩達中間教授《几何原本》簡直是一種耻辱。一本書応該具備清晰性或正確性；両者兼備是不可能的，但缺乏其中之一則使它不配占据欧几里得在教育中所占的地位。

現代の数学における最も注目すべき結果は、象徴論理と厳格な形式主義の重要性である。ウィーナストラスの影響下にあり、現代の数学者たちは正確性への配慮と、曖昧な論理を避ける傾向を示し、かつてギリシャ時代以来見られなかったものであった。17世紀の偉大な発明、解析幾何学と無限小計算法は新たな結果を豊富に生み出し、数学者たちはそれらの基礎を検討する余裕や意欲がなかった。哲学者はこの役割を担うべきだったが、数学的な能力が足りず、現在必要とされる新しい数学分野を発明することができなかった。そのため、数学者たちはウィーナストラスやその追随者によって多くの最も愛された命題が一般に偽であることが示され、初めて「神学的睡眠」から覚めざるを得なかった。マカロイは数学の確実性と哲学の不確実性を対比させ、誰がテイラーの定理に対する反動を聞いたかと尋ねる。もし彼が現代に生きていたなら、自身もそのような反動を耳にするかもしれない、なぜならこれは現代の調査によって覆された定理だからである。このような数学信仰への粗野な衝撃は、象徴主義への愛着を引き出し、それが motive が不明瞭な人々には単なる過度な学問的礼儀作法のように見られるようにした。

すべての純粋数学、包括的な幾何学も含め、それが形式論理のただの一種であるという証明は、カントの哲学に致命的な打撃をとえます。カントは、ユークリッドの命題がユークリッドの公理から導き出されるためには図形が必要であると正しく認識し、その事実を説明するための知識論の理論を作りました。この理論は非常に成功し、ユークリッドの欠点が幾何学的推論の性質の結果ではないことが示されたとき、カントの理論も棄却せざるを得ません。純粋数学の可能性を解釈するためにカントによって説明された先験的直観の全理論は、現在の数学の形態において適用できません。アリストテレスの学問の教義は現代の数学が inspires する教義に精神的に近いですが、学問は形式論理が非常に欠陥であり、三段論法に基づく哲学的論理が対応する狭さを示すことが妨げでした。現在求められているのは、最大限の可能性で数学的論理を発展させ、関係の重要性を十分に考慮に入れ、その上でこの安全な基礎上に新しい哲学的論理を建立することです。これにより、純粋哲学において近い将来が数学の原理における直近の過去ほど大きな時代となる理由があります。偉大な勝利は偉大な希望をもたらし、我々の世代には純粋思考が、その数学的な基礎と同等の正確さと確実性を得ることで、ギリシャ最盛期に匹敵するような結果を達成することが期待されます。

FOOTNOTES:

 この主題は主にMr. C.S. Peirceによって作られました。

 私がこの記事を書いたときFregeの著作はまだ知らなかったので、彼を追加すべきでした。[1917年に追加された注]

 ベルリン大学の数学教授。1897年に亡くなりました。

 [1917年に追加された注] ある2つの無限数のうちの一つが必ず大きいとは証明できません。いくつかの無限数は他のものよりも大きいかもしれませんが、任意の2つの無限数についても一方が必ず大きいとは証明できません。

 Cantorはこの点で誤りを犯していません。彼の無限数が最大のものではないことを示す証明は有効です。このパズルの解決策は複雑であり、型論理に基づいて説明されています。\textit{Principia Mathematica} 第1巻（Camb. Univ. Press, 1910）で解説されています。[1917年に追加された注]

 この記述はゼノンが実際に考えていたことを歴史的に正確に述べたものではない。これは彼の結論を支持する新しい議論であり、影響をとえた議論とは異なるものである。この点については、C.D.ブラッド氏の「アキレスと亀に関する補足」（Mind, N.S., Vol. XXII, pp. 318-19）を参照されたい。

 この文書が執筆された後、彼は教科書として使用されることをやめることになった。しかし、現在使用されている多くの本が非常に悪いものであるため、この変更があまり改善ではないという懸念がある。

 ペロポネソス戦争によってギリシャの黄金時代は終わりを迎えた。


\chapter{第VI章：哲学における科学的方法について --- {ON SCIENTIFIC METHOD IN PHILOSOPHY}}

VI

哲学における科学的方法


我々が哲学的な問題への調査の動機を確認しようとするとき、広く言えば、これらは二つのグループに分けられ、しばしば対立し、非常に異なる体係へと導きます。これらの二つの動機群は、一方で宗教や倫理から派生したものであり、他方で科学から派生したものです。プラトノ、スパニオラ、ヘゲロールは、主に宗教や倫理に関心を持つ哲学者の典型例として挙げられ、レービンツ、ロック、フームは科学的な翼を代表する者として考えられます。アリストテレス、デカルト、ベルキル、カントでは、両方の動機が強く存在します。

今日、我々がハーバート・スペンサーに敬意を表して集まったように、彼は自然哲学者の列に分類されるべきでしょう：彼は主に科学からデータ、問題の定式化、方法論の概念を得ました。しかし、彼の強い宗教感覚は多くの著作で明らかであり、倫理的な関心が進化的な概念を重視させ、それが全体の世代において科学と道徳が結婚する果実豊かな不可分の婚姻となるという信念を形成しています。

私はこれらの倫理的・宗教的な動機が、尽管它達産生了壮麗而富有想象力的体係，総体上来説反而阻碍了哲学的進步，並且希望発現哲学真理的人応該有意識地将其擱置一旁。科学最初也是被類似的動机所糾纏，从而妨碍了其進步。我認為，哲学応当从科学中汲取灵感，而不是从倫理和宗教中。

しかし、哲学が科学に基づくためには二つの異なる方法がある。一つは最も一般的な「結果」を強調し、これらの結果にさらなる一般性と統一性をとえることを目指すことができる。もう一つは科学の「手法」を研究し、これらの手法を必要に応じて適応させて自身の特殊な領域に適用することである。科学に基づく多くの哲学が、一時的に達成されたと見做された「結果」への関心から迷い込んだ。しかし、利益を得るためには、特別な科学の範囲から哲学の範囲へ移行できるのは「手法」であり、「結果」ではない。私がお知らせしたいのは、科学的な問題に成功的に適用される特定の広範な方法論原則を哲学的問題に適用する可能性と重要性である。

科学的手法に基づく哲学と宗教的・倫理的思想が支配する哲学との対立は、哲学者の著作で非常に一般的な二つの概念を通じて例示できる。すなわち、「宇宙」の概念と「善悪」の概念である。哲学者は、全体的な宇宙の性質について何かを我々に語り、楽観論や悲観論の根拠をとえることを期待される。これらの期待は私には誤りだと考えられる。私は「宇宙」という概念を、「宇宙」という言葉の由来が示すように、コペルニクス以前の天文学の遺物であると信じる。そして私は楽観論や悲観論の問題を、哲学者はその範囲外と考えるべきであり、おそらくそれが解決不能であることを主張する程度に限られるべきだと信じる。

インコネリウス以前の時代において、「宇宙」の概念は科学的な根拠で正当化されていた：太陽係の天体が昼夜の回転によって結びつけられ、地球が中心である一つのシステムの一部として統合されていた。この表面上の科学的事実を中心に、多くの人間的欲望が集まり：人間が万物の秩序において重要な存在であることへの願望、全体に対する包括的な理解を求める理論的欲求、自然の流れが私たちの希望と sympathize する何かによって導かれるという期待。このような方法で、倫理的に動機付けられた形而上学の体係が発展し、天文学的な地心説に基づく人間中心主義は明らかに正当化されたように見えた。コペルニクスがこの思考体係の天文的基礎を取り除いたとき、それは非常に馴染み深くなり、人々の希望と密接に関連していたため、その力はほとんど減少することなく生き延びた——「コペルニクス革命」すらもその力を弱めることはなかった。そして今でも、多くの形而上学的な体係の無意識的な前提となっている。

世界の一体性は、ほとんどの形而上学においてほぼ不問の公理である。実在は単に一元的で自己矛盾しないだけでなく、「相互決定的な部分から成るシステム」である——このような声明はほとんど無視され、単なる真理として通過するだろう。しかし私は、その中には「コペルニクス革命」の完全な実施が欠けていると信じており、世界の一見一体性は単に単一の観察者によって見たものや単一の心で apprehended されたものの一貫性であると理解されるべきだと考える。批判的哲学は、多くの世界の表面上の特性において主体的な要素を強調する意図を持っていましたが、その本質について無知であると見做すことで、注意は主観的な表現に集中しすぎてしまい、その主観性はすぐに忘れられてしまった。分類概念を心の働きとして認識した後で、自身の認識によってパラライズされ、主観的虚偽化の工作を取り戻そうとする試みを諦めてしまった。おそらく一部においては、その絶望は正当なものだったかもしれませんが、絶対的な或る意味での最終的なものとは言えなかった。ましてや喜びの理由ともなり得ず、それを正当に交換できる形而上学的独断主義と誤って考えることもできない。


現時点での問題、すなわち世界の統一性に関する問題については、ウィリアム・ジェイムズが正しい方法を示唆していると考えます。「では、今度は世界の一体性について説明する無能な或いは理解不能な方法から背を向け、一体性が原則であるか否かではなく、その「一体性」が単なる「存在」のような名前で、経験的流れの一部に見られる特定の且つ検証可能な関連性を描写するものではないかと問い合せてみましょう……私たちは互いに関連性のないものを思い描くことができます。これらは異なる時間と空間に存在し、現在でも異なる人々の夢のように独立しています。それらは全く異なって且つ不可比較であり、互いに対立せず、干渉することもありません。現在でも実際には我々が知る世界とは全く異なる多くの宇宙があり、その存在を我々は認識する手段がないかもしれません。しかし、その多様性を考えることによって、これら全ては論理学において「議論の対象となる宇宙」となります。この例から明らかであるように、「宇宙」を形成することはさらなる種類の関連性を要求しません。特定の一元主義者の著者が何らかの混沌が名前で呼べばただちに一元的な世界になるという事実に重大な重要性を付とすることについては、私には理解不能です」したがって、経験的世界における二種類の統一性に残されるものがあります；一つは我々が体験するものの共通性であり、これは単に私の経験的世界が存在全体から選ばれた一つの経験によるものであるという事実に基づくエピステミオロジー的な統一性です；もう一つは科学法則によって示される暫定的で部分的な統一性です。これらのいずれかに基づいて一般化することは誤りです。私たが体験するものの共通点は、我々がそれを経験しているという事実から導き出す重要な何かを推論することから明らかに導くことができない：つまり、何でも私たが体験しているという事実から、従って全てが私たが体験しているという結論を導くことは誤りです。科学法則に基づく二番目の統一性の一般化も同様に誤りですが、その誤りは少し基本的ではありません。それを説明するために、法則支配の意味を考えましょう。人々はしばしば物理世界が不変の法則によって支配されているという事実を驚くべきこととして述べます。しかし実際には、このような世界が一般法則に従わない理由を見つけるのは難しいです。空間内の任意の点群に対して、これらの点に対応する時間の関数があります、すなわち、これらの点を通過する粒子の運動を表現します：この関数はそのような粒子の行動が従属する一般法則として見做すことができます。宇宙全体のすべての粒子に対するこのような関数を考えると、理論的にはそれら全てを包み込む一つの公式があり、これは空間時間世界の単一で最高の法則と見做すことができます。したがって物理学における驚くべきことは一般的な法則の存在ではなく、その極めて単純さです。自然の均一性自体に驚くべきことではない、十分な解析的巧妙さがあれば、どんな自然の進行も均一性を示すことが示されるからです。驚くべきはその均一性が我々が発見できるほど単純であることです。しかし、これまで発見された法則のこの特徴的な単純性を一般化することは誤りで、それは明らかにそのような単純性が発見の一因であり、したがって他の未発見の法則も同様に単純であるという仮定の根拠をとえることはできないからです。

これらの二つの種類の単一性に生じた論理的誤謬は、哲学において科学が達成した一般化された結果を使用する際の注意を示唆します。まず第一に、過去の経験を超えてこれらの結果を一般化する際に、その結果が経験されたすべてに対して有効であるよりも普遍的に有効である可能性が高い理由がないか十分に慎重に検討する必要があります。人類が経験した総和は存在全体から選択されたものであり、この選択によって示される一般的な特性は、選択の方法ではなく、経験が選択するものの一般的な特性によるものである可能性があります。第二に、科学の最も一般的な結果は最も不確実で、その後の研究により容易に覆される可能性が高いです。これらの結果を哲学の基礎として利用する際には、科学法の最も貴重かつ特出した特性を犠牲にする必要があります。すなわち、科学においてほとんどすべてが最終的に修正が必要となることがありますが、その修正は通常、既に誤りと判明した前提から導き出された結果の大半を根本的に変更せずに残ります。慎重な科学者には、現在の科学的信念を使用する適切な方法についての一定の直感が育ちますが、これはおそらくその後の発見によって導入される修正により完全かつ絶対的な否認から危険を被ることなく利用できるものです。不幸にして、広範囲な科学的推論を利用した哲学の基礎は、科学的慎重さの直感が避けるべき種類の使用であり、通常、そのような推論に基づく一般化が修正の必要性がない限り真実の結果を導き出すことはありません。

これらの一般的な考慮事項を二つの例によって説明することができます。すなわち、エネルギー保存則と進化原理です。

（1）まずエネルギー保存則について始めましょう。ハーバート・スペンサーはこれを「力の持続」と呼んだことがあります。彼は言います:

進化の理的な解釈に初めての一歩を踏み出す前に、物質は破壊不能であり運動は連続的であるという事実だけでなく、力も持続するという事実を認識することが必要である。全体と詳細に対する変貌が原因となる機関によって引き起こされるなら、その機関が存在し始めるか消失するかというような試みは進化の原因を割り当てるものとして明白に無意味である。その場合、現象の順序は全く恣意的であり、推論科学は不可能となる。

この段落は、哲学者が経験的総括を絶対性と必然性に装う傾向を示しています。科学の未調査領域での近似的な真実だけが、未補助の科学的手法によって保証される限界内においてです。しばしば言及されるように、何かが持続していることはすべての科学的調査にとって必要な前提とされ、その前提は物理学が定める一定量を例示すると考えられています。私にはここに三つの異なる誤りがあると見えます。

まず第一に、詳細な科学的な自然調査は、その結果が確認されるような一般的法則を前提としていません。特定の観察以外を除き、科学は何も前提する必要はありませんし、論理の一般原理のみを前提します。これらの原理は自然の法則ではなく、仮説的であり、実際の世界だけでなく何が可能であるかに対しても適用されます。

第二の誤りは一定量と持続的な存在の同一視です。エネルギーは物理係の特定の関数ですが、物理係の変化を通じて持続的に存在するものではありません。質量も同様で、物質の量としてしばしば定義されるにもかかわらずです。

量の概念自体が、その定義に多くの規約に基づく数値測定を含むことから、数学的な利便性の実現形態よりも遥かに人工的であり、一般的に物理哲学を行う人々はそれほど信じていないことが明らかです。したがって、何かの存在の持続性が科学の必要前提であるとしても（私はそのような見解には一時的にも同意できません）、これからどの物理量の恒常性を推論したり、そのような恒常性の先験的必然性を経験的に発見した場合の先驗的必然性を推論することは誤りです。

第三に、物理学の進展とともに、エネルギー保存則や質量保存則のような大きな総括は、決して確実ではなく、おそらく近似であることが明らかになっています。質量は昔は最も疑いの余地のない物理量とされていましたが、現在では速度に応じて変化すると考えられ、実際には方向によって異なるベクトル量であると信じられています。物理学で研究していたような運動に対する質量の恒常性から導かれた詳細な結論は非常に正確であり続けますので、古い調査分野における古い結果の修正はほとんど必要ありません。しかし、質量またはエネルギーの保存則を普遍的な先驗的法則として建立すると、絶対的な正確さのわずかな欠如も致命的となり、その基礎に築かれた哲学的な構造は必ずしも破綻します。したがって、慎重な哲学者は、方法によって得られた最も一般的な結果が一時的にも最適であると見做されるものに基づいて何かを立てるのには非常に慎重になります。

(2) 進化論の哲学は、私たちの二番目の例となり、同様に急いでした一般的化の傾向を示します。また、倫理的な概念への過度な関心という別の種類の一般的化も示します。進化論には二つの種類があります。ヘーゲルとスペンサーは古い未発達な種類を代表し、プラグマティズムやベルゲンソンはより現代的で革命的な種類を代表します。しかし、これらの両方の種類の進化論には共通点があり、それは「進歩」への強調点です。すなわち、悪いものから良いものへと継続的に変化するか、単純なものから複雑なものへと変化することです。ヘーゲルに科学的な動機や基礎を atribu することは不公正ですが、他のすべての進化論者、特にヘーゲルの現代的弟子も含め、彼らは生物発展の歴史から得た推動力が非常に大きいです。この歴史から普遍的な進歩の法則を得た哲学に対して二つの反論があります。まず第一に、その歴史自体は空間と時間の無限小の一部に限定された極めて狭い事実を選択しており、科学的に見てもおそらく世界全体の事件の平均的なサンプルではない可能性が高いです。私たちは知っていますが、成長だけでなく衰退も世界で正常な現象です。外宇宙の哲学者が一人の青年を21歳まで観察し、他の人間と出会ったことがない場合、その人は人類は無限に完璧に向かって無限に高くなり、賢くなる性質であるという結論を下すかもしれません。そしてこの一般的化は、進化論者がこの惑星の過去の歴史に基づいて行う一般的化と同じ根拠を持っているでしょう。しかし科学的な反論以外にも、進化論が進歩という概念自体に過度な倫理的概念を混ぜ込むことから派生するもう一つの反論があります。生命はプロトゾイアから哲学者へと徐々に発展し、これは明らかに前進であると確信されています。不幸にも、それが哲学者ではなくプロトゾイアがこの確認をとえているのです。私たちは誰もが公平な第三者が哲学者の自慢げな仮定に同意することを保証することができません。この点は、哲学者・荘子によって以下のような説得力のある逸話で示されています：

大 augur は祭服を着てslaughter場に近づき、豚たちに対して次のように述べた。「汝等が死を拒否するのは如何なる理由があるか。汝等を三ヶ月で肥育する。自分自身を十日間律し、さらに三日間断食する。細かい草を敷き、汝等を彫刻された祭壇の上に置く。これで不満はないか？」

その後、豚たちの視点から述べた。「しかし、結局は、糠のみを食べながらslaughter場から逃れる方が良いのかもしれない...」

「しかし、それでは」と彼自身の視点から続けた。「生きて栄光を得るためには、戦場の盾や行刑者の籠で死ぬのが当然である。」

そこで彼は豚たちの視点を棄て、自分の視点に立った。すると彼は豚とは何か異なるのか？

進化論者もしばしば大 augur と豚と同じような態度を取っていることを恐れる。


多くの有名な哲学体係において重要な役割を果たしている倫理的要素は、私の見解では、科学的方法が哲学的な問題調査における勝利を得るための最も深刻な障害の一つである。人類の倫理的概念は、 Chuang Tz[)u] が指摘したように、本質的に人間中心であり、これらが形而上学で使用されるときには、現在の人々の欲求に基づいて宇宙を立法する試み、無論それが隠されているとしても、そのような試みは世界に対する科学的態度の基本である事実性への受容性に干渉する。世界の理解のために倫理的概念を鍵と見なすことは、コペルニクス以前の考え方である。それは、現在の瞬間にある希望や理想を持つ人間を宇宙の中心とし、そのような目標と目的の解釈者とする試みである。倫理的形而上学は、自らの願望に立法的な力をとえる試みであり、それが隠されているとしても、当然疑問視されるべきものかもしれないが、私は倫理的概念がどのように生じるかを考慮する際にそのような主張が正当化されることを確認している。

倫理的概念は、群れ本能、つまり自らのグループと他者のグループとの協力の本能から発生すると考えられる。自らのグループに属する者は善であり、敵対的なグループに属する者は悪である。自らのグループが追求する目的は望ましいものであり、敵対的なグループが追求する目的は不道徳なものである。この状況の主観性は群れ本能を持つ動物には明らかではなく、その動物は自己の群れの正義の一般原則が自分たちの側にいると考える。倫理的概念を形而上学的に考案したとき、それは自己の群れの正義の信念に基づいて倫理的概念を考案する。そこで、Grand Augur は倫理的概念を Augurs と豚との争いにおける正当性の証拠として引用する。しかし、倫理的概念には自殺的な概念が含まれていないという反論があるかもしれない。しかし、これは誤りである。群れ本能を持つ動物が生存競争において成功するのは、自己の群れ内の協力によるものであり、その協力は個々人の利益を一部犠牲にすることを要求する。そのため、自衛と自己の群れの保存という双方の本能が生物的に個々人にとって目的であるため、欲求と本能の衝突が生じる。倫理的概念は本来、自分自身との協力を推荐するために他人に対して必要な犠牲を推奨する芸術である。したがって、反射的に社会的正義によって自己に対する犠牲も推奨されるが、いかなる程度にでも主観的な倫理は残る。すなわち、ベジタリアンたちは発熱中の人間の命を救うために多くの微生物の命を破壊しても躊躇しない。倫理的概念に基づく哲学が世界を見ることで得られる視点は常に不偏なものではなく、したがって完全に科学的ではない。科学と比較して、それは人間が望むような世界理解のために必要な自己からの想像的な解放を達成することができず、その哲学は常に多少の地方性を持ち、時代や場所の偏見に汚染されている。

私は倫理的な概念に基づく哲学の重要性や価値を否定するものではありません。そのような哲学の中で、スピノザの倫理学は私にとって極めて重要な意味を持つように思われます。しかし、その価値は、それが生じ得る世界観に関する任意の形而上学的な理論、あるいは論証によって証明できるか否かに関わらず何かではないのです。価値があるのは、生活や世界に対する新たな感受性の示唆、つまり私たち自身の存在が深く望む特性を持つものになる可能性のある感受性です。このような工作的な重要性は無限大であり、実践的な側面と関連しています。その理論的重要性は人間性に対してのみあり、宇宙全体に対しては存在しません。したがって、世界を理解することだけを目指す科学的哲学は、事実に屈服するという科学的精神の本質から逸れてしまうかもしれません。

II

万有論や善悪の概念が科学的哲学から排除されるとすれば、哲学者と科学者との間で残る具体的な問題は何となるのでしょうか？この問いに対する正確な答えを出すのは難しいかもしれませんが、哲学の領域と特殊科学の領域を区別する特徴はいくつか指摘できます。

まず、哲学的命題は一般性を持つべきである。地上の事物や太陽係、あるいは空間と時間の他の部分について特別に扱うものでない。この一般性が必要な理由から、哲学が宇宙全体を扱うものだと信じられている。私はその信念が正当化されないと考えているが、一方で哲学的命題は存在するすべての事物や将来存在する可能性のあるすべての事物に対して適用されるべきであると信じている。

この承認が私が否定したい見解とほとんど区別できないと考えるかもしれない。しかし、これは誤りであり重要な誤りだ。伝統的な見解では、宇宙そのものがさまざまな述語の対象となり、これらの特異な述語を宇宙に付とすることは哲学の仕事だとされる。一方で私は、「宇宙」が主題となる命題は存在しないと主張する。つまり、「宇宙」というものは存在しないと主張する。

私が推奨したい哲学は、論理原子主義や絶対多様主義と呼ばれるべきである。なぜなら、多くの事物があることを認める一方で、それらの事物からなる全体が存在しないと主張するからだ。したがって、哲学的命題はすべての事物を総合的に扱うのではなく、個々の事物についての性質を分布的に扱うべきであることが明らかになる。そして、これらは偶然的な事物の性質に依存せず、どのような可能性のある世界でも真実であり、感覚によってしか発見できない事実から独立したすべての事物の性質についても言及すべきである。

この文章は、哲学的命題の二つの特性を導き出すもので、まずそれらが\textit{a priori}であるべきであることを示しています。哲学的命題は経験的な証拠によって証明されたり否定されたりすることのないようなものでなければなりません。しばしば哲学書には歴史の流れや脳の構造、貝類の目などといった特別な事実に基づく議論を見かけますが、これらは哲学とは無関係であり、哲学はどのような実在の世界であっても真となるような主張のみを述べるべきです。

これらの二つの特性を要約すれば、哲学は可能性の科学であると言えます。しかし、この声明が解説されないままでは誤解を招く恐れがあります。可能性とは一般性と異なるものと考えられがちですが、実際には両者は区別できません。

もし先に述べたことが正しいとするならば、哲学は現在の意味での論理学から区別せずに済むことになります。論理学の研究は大まかに言えば二つのあまり明確に区別されない部分で構成されています。一方では、何か一つのものを言及することなくすべてのものについて述べられるような一般的な主張に関心を持ちます。例えば「$x$がアルファクラスのメンバーであり、アルファクラスのすべてのものがベータクラスのメンバーであるなら、$x$はベータクラスのメンバーである。$x$, アルファ，ベータは何であれ同じです」といったような主張です。他方では論理的な形態の分析と列挙に関心を持ちます、すなわちどのような命題が発生するか、どのような事実があるか、そして事実の構成要素をどのように分類するかについてです。このようにして論理学は可能性の目録を提供し、抽象的に支持可能な仮説のレパートリーを作り出します。

SUCH 研究は、あまりに曖昧で一般的であるため、非常に重要な意味を持つとは思われないかもしれません。また、その問題が特定の点で十分に明確になると、それらは一部の特殊な科学の問題と合流するだろうと考えられるかもしれません。しかし、実際にはそうではありません。例えば空間と時間の分析、感覚の性質、判断理論などのいくつかの問題では、事実の論理的な形式を発見することが作業の中で最も困難でかつこれまで最も不足していた部分です。このような問題は、適切な論理的仮説が欠如しているために、これまで不十分な方法で扱われ、哲学者たちの中には理由を否定する者たちによって常に喜ばれてきた矛盾や反律法に至る原因となっています。

論理的な形式の調査に焦点を当てることにより、哲学は問題を個別的に取り組むことが可能になります。そして、科学が得られるような部分的でおそらく完全ではない結果を得ることができます。その後の研究によって補完され改善されるものとして利用することができます。これまでの大多数の哲学は一塊で構築されており、それらが完全に正しくない場合でも完全に誤っているため、さらに詳しい調査の基礎には使用できません。この事実により、哲学は科学とは異なり、これまで進歩的ではなく、各原始的な哲学者は最初から仕事を再開しなければならず、前人からの何らかの確定したものを受け入れることができませんでした。私は推奨したい科学的な哲学は、他の科学と同様に部分的で仮説的なものになります。特に、不完全であっても修正が施された後でも依然として生産的な仮説を発明する能力があります。真実への逐次近似の可能性は、科学の勝利の源よりも何よりも重要であり、これを哲学に転送することで、方法論における進歩がどれほど重要であるかを過大評価することもほとんど不可能になります。

哲学の本質は分析であり、合成ではない。世界のシステムを構築するような作業、ヘイネのドイツの教授が生活の断片を組み合わせて理解できるシステムを作り出すような作業は、私には philosopher's stone の発見よりも実現不可能であると考えられる。しかし、一般的な形式の理解と、伝統的な問題を個別のそしてより解明しやすい質問に分割することにより解決する可能性がある問題の分類が可能だ。成功のための格言「 Divide and conquer 」はここでも他の場面でも当てはまる。

これらの格言を哲学の空間についての適用例で説明することで、その意味や重要性を理解することができます。仮にカントの『超越的美学』において空間に関する問題が提示され、この問題の要素は何であるか、それらの解決の希望があるのか探求したいとしましょう。すぐに明らかになるのは、カントが取り組んでいると思われる単一の問題は、論理、物理学、認識論という異なる分野に属する三つの完全に異なる問題を混乱して組み合わせているということです。論理に関する問題、物理学に関する問題、そして認識論に関する問題があります。これらの三つの中で、論理に関する問題は正確かつ完全に解決できる；物理学に関する問題は実証的な領域において希望される程度の確実性と正確さでおそらく解決できる可能性がある；しかし認識論に関する問題は非常に曖昧で扱いが難しいまま残ります。これらの三つの問題がどのように生じるかを見てみましょう。


非ユークリッド幾何学の提案を通じて論理的な問題が生じました。幾何学的命題の体をとえられた場合、その命題から導き出すための最小限の公理を見つけることは難しくありません。また、これらの公理を一部省略したり変更することで、より一般的なまたは異なる幾何学を得ることも容易です。純粋数学からは同様の論理性と尊重される権利を持つ Euclid の幾何学と同じように、論理抽象から見れば無数の算術係が存在し、それぞれはEuclidean空間と等しく、不変かつ正当な名を有します。したがって、空間は論理的または数学的な研究対象として唯一性を失います；種類は多く、各種類には無限の例があり、物理学の空間がそのいずれかの例であるかどうかは困難で、物理の空間が確実にそのいずれかの例であるかどうかは不可能です。論理的システムの幾何学の一例として、実際の空間における「間接」関係と類似する一定の形式的な点に関する三項関係すべてを考慮することができます。空間は、このような三項関係によって定義されます。空間内の点は、何かに対してこの関係を持つすべての項であり、その空間での順序はこの関係によって決定されます。一つの空間の点は他の空間の点でもあり、それらと同じ点を持つ他の三項関係が存在するからです。実際のところ、空間はその点のクラスではなく、その定義をなす三項関係によって定められます。このような関係の十分な抽象的な論理的性質を数え上げて、例えばEuclid幾何学のような種類の幾何学を決定した後では、純粋な幾何学者はそのような関係のいずれか一つではなく、それら全体を対象とします。したがって、とえられた種類の幾何学を研究する際には、特定の抽象的な論理的性質によって定義される一クラスの関係を研究していることになります。幾何学的な推論は従って純粋に帰納的かつ論理的です；幾何学における特殊な認識論的特異性が見つかるなら、それは推論ではなく、特定の空間における公理に関する知識の中にあるはずです。


物理的な空間に関する問題は、論理的な問題よりももっと興味深く、難しくなります。この物理的な問題は次のように述べることができます：物理的世界において、あるいは物理的材料から、論理的な幾何学で列挙されたいずれかの種類の空間を見つけるか、または構築するという問題です。この問題が困難である理由は、実在の世界の粗さと曖昧さに適応しようと試みるためであり、これは純粋な数学において邏輯清晰性と正確性を持つシステムを持っています。もし私が三人の人間A, B, Cが並んで座っているのを見ると、BはAとCの間にいるという事実を認識し、それよりもAはBとCの間にいる、またはCはAとBの間にいるという事実よりも認識します。この「間にいる」という関係性は、我々が見た三項関係が幾何学を生む抽象的な論理的性質を持つ部分を持っていますが、その性質は正確ではなく、経験的にとえられたものとして幾何学が目指すような扱いには適応しません。抽象的な幾何学では点、直線、平面と対象を扱いますが、私は三つの人間A, B, Cを見ているとき、これらは正確な点ではなく、並んでいる列も正確な直線ではありません。しかし物理学は形式的に点、直線、平面を含む空間を仮定しますが、実感的世界に適用可能な結果を生み出します。したがって、物理的なデータやそれに似た仮定的な追加を用いて物理的点、直線、平面の解釈を見つけることが可能であると推測できます。すべてのデータは一定の大きさを持ち、多少輪郭が曖昧であるため、もし点という概念が実感材料に適用されるなら、点はデータや仮定的な追加ではなく、データとその仮定的な追加を用いて構築されたものでなければなりません。任意の仮定的なデータの補完がより確実で不満足な場合、それはデータに類似した追加よりも根本的に異なる種類の場合です。例えば、私があなたの目をそむけた後でも見える物体が、私が見ていたときと比較的類似しているという仮定は、数学的な点からなる無限数の物体よりも激しい仮定ではありません。したがって物理空間の幾何学的研究においては、論理的な幾何学で最初に仮定されるように点を仮定してはなりませんが、データとその仮定的な類似体から構成されたシステムとして解釈する必要があります。これにより物理世界の最終的成分である物理的な対象全体の特定のクラスを物理的な点として定義することができます。これはすべての対象が点を含んでいるという自然な言葉で説明できます。この定義をとえるために、物理的な物体が点で構成されていることを事前に仮定せずに得られる結果は、数学的論理学における快適な問題となります。この問題の解決とその重要性の認識は私の友人ドクター・ホワイトヘッドによるものです。点を物理的存在のクラスとして見なす奇特性は慣れることで薄れていき、実際には幾何学的な仮想的存在であるという主張をする人々にとっては感じられません。多くの男性がこのような状況において「仮想的存在」という言葉を使いこなすにもかかわらず、その必要性を説明することなく、仮想的存在が実在の世界の研究で有用になる方法を理解することは不思議です。我々の定義によれば、点を物理的存在のクラスとして見なすことにより、点を使用することで重要な物理的結果につながる方法と、点が物理的世界の存在であるという前提を避けることができる方法が説明されます。

多くの抽象論理空間に固有の数学的な便利な性質は、物理学の空間に属するかどうかを知ることはできず、またそうではないことも知ることはできません。これらはすべて連続性に関連した性質です。実際の空間がこれらの性質を持つことを知るには、感覚的認識が無限に正確であることが必要です。もし実際の空間が連続的なものであるとすると、それでもなお多くの可能性のある非連続的な空間があり、これらは実験的に実際の空間から区別することはできません；逆もまた真で、実際の空間が非連続的であっても、それは実験的に可能性のある連続的な空間から区別することはできないかもしれません。したがって、連続性は算術の\textit{a priori}領域において得られるものですが、物理世界の空間や時間において確実に得られるものではありません：これらのものが連続的であるかどうかは、答えられていないだけでなく、永遠に答えられないような問題であるように思われます。しかし、哲学的な視点から見れば、質問が回答不能であることが発見されたことは、何らかの回答よりも完全な回答と言えるでしょう。そして物理学において、区別するための実験的方法がない場合、最も単純な仮定、すなわち連続性に対する実験的な反対は存在しないと言えます。

物理的空間論の主題は非常に広範で、これまであまり探索されていません。これは時間に関する類似の理論と関連しており、両者は相対性理論に関する議論によって哲学的に志向した物理学者たちの注意を引きつけました。


 Kantが『超越的美学』において取り組んでいる問題は、主に認識論的な問題です：「我々はどのようにして先験的に幾何学的知識を得るのか？」この問題の意味と範囲は、几何学の論理的問題と物理的問題との区別によって大きく変わります。純粋な幾何学的知識は先験的ですが、完全に論理的なものです。物理学的幾何学的知識は合成的ですが、先験的ではありません。純粋な幾何学的知識は仮定的であり、平行公理が物理的世界で真であると主張することはできません。物理学的幾何学的知識は、この公理が大体正確であることを主張することができますが、観測の避けられない不正確さにより、それが完全に正確であるとは主張することができません。したがって、純粋な幾何学と物理学的幾何学との区別によって、カントの問題は崩れ去ります。「合成的な先験的知識は可能か？」という問いに対する答えは、少なくとも幾何学に関して言えば、「不可能」です。もし「合成的」という言葉が「論理だけで導き出されるものではない」という意味であるなら、幾何学的知識は論理と感覚から部分的に得られるものです。カントの時代に幾何学が占めていた特別な位置づけは幻影でした。真実を追求する一部の哲学者は、平行公理が実際的な空間で真であることを知る知識は、経験的説明ではなく、カントが主張したように先験的直観から得られるべきだと主張します。この立場は論理的には反証不能ですが、物理空間の概念がいかに複雑かつ派生的なものであるかを理解すると、その説得力は失われます。我々が見たように、幾何学を物理学的世界に適用する際には、実際的な存在の中に点や直線がある必要はありません。したがって、経済性の原則から、点と直線の存在を仮定することは避けるべきです。しかし、点と直線は物理的要素のクラスによる複雑な構築であるという視点を受け入れた場合、直線が無限に延長されたときにどのようなことが起こるかを先験的に知っているという仮説は非常に強引で不自然になります。このような仮説は、物理空間の性質を理解した哲学者には思いつかないでしょう。カントはニュートンの影響を受け、絶対的な空間の仮説を採用しましたが、この仮説はオッカムの剃刀によって削除され、絶対的な空間は物理的世界の解釈において不要な存在です。したがって、我々はカントの先験的直観理論を反証することはできませんが、問題の分析を通じてその根拠を一つずつ取り除くことができます。したがって、多くの哲学的問いと同様に、解析法は論理的に結論を導くことはできませんが、ある理論の肯定的な根拠すべてが虚偽であることを示し、より自然な理論が事実を説明できる可能性があることを示すことができます。

分析法の能力を示すもう一つの問題は現実主義に関する問題です。現実主義を主張する者と反対する者がいるが、両者は私が彼らが議論している問題の性質について十分に明確でないように思われます。もし私たちは「我々の感覚的対象は\textit{現実的}であり、それらは\textit{感知者}に依存していないか？」と問うとするなら、「現実的」と「依存しない」という言葉に何か意味を付とする必要がありますが、現実主義論争のどちらの側もこれらの二つの言葉を定義しようと求められると、その答えには論理的な分析によって明らかになる混乱が含まれていることがほとんどです。

Let us begin with the word "実在"。確かに感知対象があることは確かで、したがって、「これらの対象は実在であるか否か」が重要な問いとなるなら、世界には二種類の対象が必要であり、すなわち実在と非実在があり、しかし非実在とは存在しないものとして本質的に定義されている。ある対象が実在であるためにはどのような性質を持つべきかという問いは、しばしば適切な答えを生み出すことが稀である。もちろんヘゲロールの答えがあるが、それは全体のみが自己矛盾しないと主張し、しかし現在の議論とは異なるレベルで進行しており、物理学的に扱われる他の対象と感知対象の地位についての問いに関心を持っている。感知対象は実在と虚在との区別よりも、一方では心理状態と他方では物質とを対比する傾向がある。したがって我々が考慮すべき問いは、世界の一部の存在に対してのみ「実在」という言葉を適用することの意味は何であるかという問いである。この意味での実在性には二つの要素が含まれていると考える。ある対象が感知されていないときでも時々持続する場合、またはまたある対象が他の対象と経験によって予想される関連性を持つ場合、その対象は実在であると言える。これらのいずれの意味での実在性も必ずしも必要ではなく、実際にはどちらの意味でも実在でない世界全体が存在する可能性があることが明らかになる。したがって感知対象が一方または両方の観点から実在でない場合でも、それが外部世界の一部であることを否定的に推論することはできない。同様な議論は「独立」という言葉にも適用される。この言葉に関連する多くの関連性は現在維持不可能な因果に関するアイデアと密接に結びついており。AがBから独立しているとは、BがAの原因の一端でなければ意味がない。しかし因果関係は単なる関連性であり、同時性の関連性だけでなく順序性の関連性もあることを認識したとき、特定の事象の因果的前件列には一貫した独特性がなく、同時に異なる前件列から新たな因果的前件列を得るための手段があることが明らかになる。因果法則を考慮する基準を明確に指定することが必要となる。先日私はある相手から手紙を受け取り、彼はいくつかの哲学的な問いについて困惑していたと述べていた。その問いを数えると彼は言っている。「これらの問いがボンからストラスブールへ私を導いた。そこでシメル教授を見つけた。」現在、これらの問いが彼の体をボンからストラスブールへ移動させたという事実を否定することは異常であるが、それでも純粋に機械的な前件列も存在し、これが物質の場所からの移動を説明できると考えられる。因果的係列の多様性により、「その」原因という概念は不確定になり、独立性に関する問いもそれに応じて曖昧になる。したがって、「AはBから独立しているか否か」という問いではなく、「BからAへ至る特定の因果法則に基づく前件列があるか否か」という問いを尋ねるべきである。この点は、特に感知対象に関する問いと関連して重要である。実際には我々が感知しないときでも全く同じような存在が存在することはないかもしれない。その場合、感知対象の存在は感知されることに依存する因果法則があることになる。しかしそれでも、純粋な物理的因果法則によって、他の感知されない対象を通じて感知される対象が存在することがあるかもしれない。その場合、そのような因果法則において感知対象は感知されることから独立していると言える。したがって、感知対象が感知されることから独立しているかどうかという問いは、その問い自体が確定的でないため、方法によっては肯定的に、また否定的に答えられる。この混乱は、この問題に関する議論を長期化させた重要な役割を果たしてきたと考える。私は推奨したい見解は、感知対象は感知されていないときにも変化しないわけではないが、おそらくそのような時間には類似した存在が存在するかもしれない；感知対象は物理的な物理学の実質的対象の一部であり、それ自体は物理的に呼ばれるべきである；純粋な物理法則が存在し、それが感知されることから独立して存在と性質、持続時間を決定しており、そのような法則の確立において物理学の命題は心理学の命題や心の存在を前提にしない；という見解である。実在論者がこのような見解を実在主義として認識するかどうかは知らないが、私はこの見解が以前の議論で提起された実在主義と理想主義に遭遇した困難を避けることができると主張する。さらにこの立場を支持し、詳細化するための防御は、現在時間がないため不足しているが、私の『外部世界に関する知識』という本の中で示されている。

科学的方法を哲学に導入した場合、私は誤りのない限り、多くの伝統的な哲学で野心的かつ人間的に興味深い問題を解決するという希望を放棄せざるを得ません。これらのうち的一些は特別な学問分野に relegatesされ、成功的な解決が期待されていません。他には、私たちの能力が根本的に解決不能であると示されるものがあります。しかし、哲学で認識されている多くの問題については、推奨される方法は独立して個別の質問への分割、仮説的かつ部分的な進歩、そして原則への訴えなどの利点を提供します。気性に関係なく、すべての適格な学生が同意するべきものです。これまでの哲学の失敗は主に急いでいて野心的であるため、忍耐と謙虚さは他の科学同様、堅固で永続的な進歩への道を開くでしょう。

FOOTNOTES:

 ボーソンケット, 『論理』, 第2巻, p. 211.

 『哲学の問題』, p. 124.

 『最初の原理』（1862年）, 第2部, 第8章の冒頭.

 オープン・コート社, 1914年.


\chapter{第VII章：物質の究極的構成要素 --- {THE ULTIMATE CONSTITUENTS OF MATTER}}

VII

物質の最終的な構成要素


この記事では、「matterとは何か？」という古くからの形而上学的問題について議論したいと思います。「matterとは何か？」という問いは、哲学的な観点から言えば、私はその答えが原則的に完全な解答と同等であると考えられる程度まで可能になったと考えています；つまり、問題を根本的に解決可能な部分と、根本的に解決不可能な部分に分けることができ、少なくとも主な輪郭についてはどのように解決するかを見ることはできるようになりました。これらの輪郭について、本稿で提案したいと思います。私の主張は現実主義的ですが、私は教授アレクサンドラの著作により多くの利益を得たことを願い、また信頼しています。彼女の立場は私の立場と密接に一致しているとも言えます。また、ドクター・ヌンの立場も同様です。

一般的な感覚は世界を心と物質で分けることに慣れており、哲学を学んだことのない人々は、心と物質の区別が完全に明確かつ容易であると考えています。二つはどの点でも重複しないし、特定の存在が精神的か物理的なかについて疑問を持つのは愚か者か哲学者だけだと考えられています。この単純な信仰はデカルトにおいて生き残り、スパノザでは多少変形した形で生き延びていますが、レービンツからは失われ始め、その時代から現在までほとんど全ての有名な哲学家たちは一般的な感覚の二元論を批判し、否認してきました。この記事では、この二元論を擁護することに意図しますが、それを擁護する前に、その拒否に至る理由について数分間議論しなければなりません。

私たちの物質世界に関する知識は、視覚や触覚などの感覚を通じて得られます。最初には、事物がそのように見えると仮定されることがありますが、二つの対立する巧妙な説明によってこの素朴な信念はすぐに崩れ去ります。一方では、物理学者は物質を分子、原子、粒子に切り分けており、将来的な必要性により更なるような分割を仮定することがあります。彼らが到達した単位は日常の視覚的・触覚的な物体とは非常に異なります。物質の単位は徐々に、空間全体を満たす電磁気場のようなものになりつつあり、その強度は小さな領域で最も高いです。このような要素からなる物質は日常の生活から遠く離れており、どのような形の metaphysical 理論よりも遥かに隔絶しています。この理論が metaphysician の理論と異なるのは、実用的な効果がその中には一定量の真実を含んでいることを証明し、商人たちがそれを基にして金銭を投資することだけです；しかし、金市場との関連性があるからといって、それゆえに metaphysical 理論とはならないのです。


感覚の世界がさらされた第二種類の sophistication は、心理学と生理学から来ています。生理学者は、私たちが見るものは眼によるものであり、聞くものは耳によるものであり、すべての感覚は脳に影響される何か（アルコールやハシなど）によっても影響を受け得ることを指摘します。心理学者は、私たちが見ると信じている中にはどれだけの部分が関連性や無意識的な推論によって補完されているか、どれだけが心的解釈であり、どれだけ残されたものが原始的なデータとして扱えるかどうかは疑問であることを指摘します。

これらの事実から心理学者は、知覚が心に静観的に受け取られるという概念は幻覚であると主張し、生理学者は、経験の分析によって純粋な感覚データが得られたとしても、それは普通の常識が仮定するように外的世の中のものではなく、私たちの神経や感覚器官の性質によって完全に制約され、変化し、それが観察される対象物と関連付けることが不可能であると考えられる方法で変化することから、そのデータは存在しないと主張します。

この生理学者の主張に対して、より説得力があるように見えて実際には堅固ではない反論があります。すなわち、感覚器官や神経の存在を知るための知識は、その生理学者が信用を失墜しようとしている過程によって得られるものであり、感覚器官や神経の存在は感覚自体の証拠を通じてしか認識されないという点です。この主張は、生理学的結果の再解釈が必要になることを示す可能性があるかもしれませんが、それが直感的な現実主義を否定するための単なる論理矛盾（reductio ad absurdum）である限り、生理学者の主張自体を覆すものではありません。

これらの Various lines of argument は、私は考えると、常識的な信念の一部を放棄せざるを得ないことを証明していると思います。それらの議論は、これらの信念全体を取り扱うと、自体矛盾する結論に至ることを余儀なくされるかもしれませんが、そのような議論自身では、どの部分が修正が必要であるかを決定することはできません。

常識的な考えには、私たちが見ることは物理的であり、心外に存在し、目を閉じたり視線を変えたりしても引き続き存在すると信じられています。私は、私たちが見ることは物理的で（いくつかの可能な意味での）心外であるという点では常識は正しいと考えますが、それがあなたが見ていない間でも引き続き存在すると考える点ではおそらく誤っていると思います。私には、物質に関する議論全体が二つの相互支持する誤りによって曖昧さを生じているように思われます。その一つは、私たちが見たり、他の感覚を通じて感知したりすること自体が主観的であるという誤りです。もう一つは、物理的なものは恒久的であると信じる信念です。

物理学が物質の最終的な構成要素として考慮する何らかのものであれ、それらは常にこれらの構成要素を不滅であると仮定します。直感的なデータが不滅ではなく永久的な変動状態にあることから、これらのデータ自体が物質の最終的な構成要素には含まれないと主張されます。私はこれに全く誤りがあると考えます。

数学物理学における恒久粒子は、非常に複雑な事実の簡潔な表現を可能にする論理的構築やシンボリックな架空と見做すべきであり、一方で、感覚上の実際のデータ、視覚や触覚や聴覚の直接対象であるものは心外にあり、純粋に物理的であり、物質の最終的な構成要素の一つと考えるべきです。

物理的な存在の非持続性についての私の意味は、ベルソンが好んで用いるシネマトグラフの例によってより明確になるかもしれません。私は最初にベルソンの主張を読んだとき、シネマトグラフを見たことがありませんでした。その主張を検証するために初めてシネマトグラフを見に行きました。そして、実際にその主張は私にとって完全に真実であることを確認しました。

映画館で人間が坂を滑り落ちるのを見て、警察から逃げるのを見るとき、河に落ちるのを見るとき、あるいはその他の場所で人々が好むような他の行為を見たとき、我々は一つの人間が動いているわけではないことを知っています。実際には一連の映画があり、それぞれに異なる瞬時の人物が描かれています。持続性の錯覚は、瞬時の人物の一連の係列における連続性に近づくことで生じます。

ここで私が提案したいのは、この点ではシネマトグラフが一般的な感覚や物理学、哲学よりもより優れたメタフィジシャンであるということです。私は警察がその同一性を証明しようと主張するいかんに関わらず、実際の人間も一連の瞬時の人物であり、それぞれが互いに異なるもので、数値的な同一性ではなく、連続性和特定の内在的な因果法則によって結びついています。そしてこれは人間だけでなく、テーブルや椅子、太陽、月、星にも同様に適用されます。

これらの各々は一つの持続的な存在としてではなく、時間と共に続く一連の存在として見なされるべきです。それぞれが非常に短い期間だけしか持続しないものの、おそらく純粋な数学的瞬間ではない期間です。この主張を述べることで、私は空間における分断と同様の時間における分断を強調しているだけです。

一立方フィートを満たす物体は多くの小さな物体から成り立つと認められます。それぞれが非常に小さな体積を占めています。同様に、一時間持続するものはより短い期間を持つ多くの存在から成り立つと見なされるべきです。物質の真実的な理論には、空間コロプルだけでなく時間コロプルへの分断も必要です。

世界は一定の模式に沿って配置された多数の存在から成り立っていると考えることができます。これらの配置をなす存在を「個体」と呼びます。個体間の関係が結果として模式や配置を作り出します。ある性質によって一括して扱うことが便利であるため、個体を集めた集合や係列を「論理的構築」または「シンボリックな架空」と呼びます。個体は、建物のブロックに例えるよりも、交響曲の音符に例える方が適切です。交響曲の最終的な構成要素（関係を除く）は音符であり、それぞれは非常に短い時間だけ持続します。一つの楽器で演奏されたすべての音符を集めて考えることができます。これらは一般的な感覚が一つの「物」の連続的な状態として考える個々の個体と類似することができます。しかし、「物」という概念は、トランペット役と同じくらい「実在的」または「本質的」であるべきではありません。このような観点で「物」を捉えると、直感的な対象が物理的なものとして扱いにくいという困難が大幅に軽減されます。

人々が「感覚の対象は精神的か物理的なものであるか？」と尋ねるとき、しばしば「精神的」や「物理的」という言葉の意味について明確な概念を持っていないし、どのような基準でとえられた存在がどちらのクラスに属するべきかを決定すべきかについても明確な考えを持っていません。私は「精神的」という言葉に対する鋭い定義をとえる方法は知らないですが、何かは精神的な出来事の一覧を挙げることで解決できるかもしれません：信じる、疑う、望む、意志を持つ、喜びを感じる、苦しみを感じるなど、これらは明らかに精神的な出来事です；また、視覚、聴覚、嗅覚、一般的な感知も同様です。しかし、これは見られたもの、聞こえたもの、嗅がれたもの、感じられたものがすべて精神的であることを意味するわけではありません。私は雷の閃光を見たとき、私の見ることは精神的なものですが、私が見ることそのものは、他の人が同じ瞬間に見ているものとは全く異なるものであり、物理学者が説明するような閃光とは非常に異なるように思えても、それは精神的ではないと言います。実際には、私は雷の閃光を見たときに現在見ている物体が存在したとしても、私の意識が消滅した場合でも、その物体は存在し続けるでしょう。ただしもちろん、私がそれを見ることはなく、私の見ること自体が精神的なものであるためです。

人々がこの観点を拒否する主な理由は二つあり、私はそれらを次のように考えています：まず第一に、私の見ることと見ているものが混同されていることが十分に区別されていなかったこと；第二に、私が見ているものは私の体の依存性によって、それが私「外部」にあるとは思われないという誤解が生じたことです。これらの理由のうち最初のものは私たちに留まらなくても良いでしょう。混乱は指摘すれば解決できますからです。しかし二つ目の理由については議論が必要で、それは一方では空間の本質、他方では因果的依存性の意味に対する現在の誤解を除去する必要があります。

人々が色やその他の二次的質量は心の中にいるのか外側にあるのかと尋ねるとき、彼らはその意味が明確であるべきであり、はいまたはいいえと答えることが可能であるべきだと考えているように見える。実際には、「内側」や「外側」という言葉は非常に曖昧である。この質問で何かが「心の中にいる」とは何を指しているのか？心は袋やパイのようなものではない；空間の特定の領域を占領するものではない、もしそういう意味で占領しているなら、その領域にあるものはおそらく脳の一部であり、それは心の中にあるとは言わないだろう。人々が感覚的質量が心の中にいると言った場合、彼らは「空間的に含まれている」という意味ではなく、パイの中で黒鳥が存在したように解釈しないでほしい。私たちは心を特定の個体の集まりと考えることができる、つまり、「心状態」と呼ばれるものであり、それらは特定の共通の特性によって一緒に属することになる。すべての心状態の共通の特性は「精神的」であるという特徴であり、それに加えて各個人の心状態には他の人々と異なる共通の特性があると仮定する必要がある。この後者の点を無視し、言葉「精神的」というもので指定される性質が観察上実際に色や音声などの感覚的対象に属しているかどうか考えてみよう。私はどんなに「精神的」の意味を知るのも難しいとしても、それは信念や希望や意志のような内在的な特性を持つものではなく物理的世界にはないものであると見なすことは明らかだと答えるべきだろうと考える。

ベーコンはこの問題について説得力のある議論を展開しているが、その根拠は「痛み」という言葉の曖昧さにあるように思われる。彼は現実主義者が火に近づくにつれて感じる熱は心の中にあるものだと考えているが、さらに近づくと熱感覚は不自覚的に痛みに変わり、痛みを心の中にあるものと考えることは不可能であると主張している。この議論に対する反論として最初に注意すべきことは、我々が直接意識する熱は火ではなく自分の体内にあることである。熱の原因を体の中で感じる熱源と判断するのは推論によるものである。次に（これはより重要な点だが）、痛みについて話すときには二つの意味を持つことができる：感覚や経験の対象である痛みそのもの、またはその性質である痛み自体。人間が大趾に痛みがあると言った場合、彼は大趾に関連する感覚がありそれが痛みという特性を持っていることを意味している。感覚自体はすべての感覚と同様に、経験的な対象を経験することであり、その経験には「痛み」という特性があるが、それは思考や欲望などと同じく精神的現象だけが持つ特性である。しかし一般的な言葉では、痛みを感じる感覚的対象を痛みと呼ぶことがあり、これがベーコンの議論の説得力に依存する混乱の原因となっている。何か非精神的なものに「痛み」の性質を atribu 付けることは荒唐無稽であるため、私たちが指す足の付け根の痛みは必ずしも精神的であるとは考えられるようになる。しかし実際には、そのような場合の感覚的対象自体が痛みではなく、それは経験である、つまり感覚的対象を経験することである。火から感じる熱が増えるにつれて、経験は徐々に快適から痛みへと変化するが、その快楽や痛みは経験に対して対象の特性ではなく、したがってこの対象が精神的であるという理由でそれを説明することは誤りである。

もし何かが心の中に存在するとすれば、それが思考や欲望のような特定の認識可能な内在的特性を持つものであるとするなら、直感に基づいて主張せねばならないのは、感覚的対象はどの心にも存在しないということだ。

しかし、「心の中に存在する」という言葉の別の意味は、感覚的対象が心の中に存在すると考える人々によって提示された議論から推測できる。これらの議論の大半は、感覚的対象が観察者に依存している因果的な関係を証明することを試みるものだ。しかし、因果的な依存という概念は非常に曖昧で難解であり、実際には哲学者たちが認識するよりも遥かに複雑である。すぐにこの点に戻って議論したいが、現時点ではこの因果的な依存の概念について批判せずに受け入れることを望む。そのような依存は私たちの体に対してではなく、私たちの心に対して存在すると主張したい。

例えば、一つの目を閉じたり、斜めに見たり、眩い物を見た後に物体を見るといった行為によって物体の視覚的姿が変化するが、これらはすべて身体的な行動であり、その結果生じる変化は生理学と光学によって説明されるべきものだ。つまり、これらの変化は眼鏡や顕微鏡による変化と同じ種類である。したがって、物理的世界の理論に属し、私たちが見ることが心に対して因果的に依存しているかどうかという問題には関係がない。

これらは実際には、私たが見たものが私たちの体に対して因果的に依存していることを証明するものであり、粗野な常識が仮定するよりも、眼や神経、脳が存在しなければ同じように存在しない何かではないという点を示すものである。つまり、顕微鏡を取り除いたとしても、顕微鏡を通じて見た物体の視覚的姿は変わらないということだ。

物理的世界が安定したとされる恒常的な構成要素で構成されていると仮定する限り、私たちが見ることが体の変化によって変化することから、その見えるものが最終的な物質の構成要素ではないという理由があるように思われる。しかし、物質の最終的な構成要素は空間的範囲だけでなく時間的にも制限されていると認識されるなら、この全ての困難は消え失う。

空間に関するもう一つの難題が残されている。太陽を見ると、その本質的な性質について九十億三千万マイル離れた太陽自体を知りたいと思うが、見ているものは眼によって影響されるものであり、九十億三千万マイル離れた場所で何かが起こっていることを眼が影響するとは思えにくい。物理学は、特定の電磁波が太陽から発生し、約八分後に眼に達することを教えてくれる。それらは棒と椎を刺激し、その結果視神経、そして脳で混乱を引き起こす。この純粋な物理的な係列の終わりには、何か不思議な奇跡によって「太陽を見ている」という体験が生まれ、これが光学神経、棒と椎、九十億三千万マイル、電磁波、そして太陽そのものについて我々が信じる理由全体を形成する。物理学と認識論によって示される因果関係の方向性と証拠の順序との奇妙な対照は、物理的実在の本質に関する最も深刻な困惑を引き起こす原因となっている。私たちの視覚が知識源として無効化されれば、物理学や生理学全体も同様に無効化される。しかし、日常的な観察から眼の信頼性を受け入れた物理学は、その因果係列を段階的に構築し、最終的には眼が見ている対象自体が九十億三千万マイル離れた最初の一因とみなされないことを示す。

この困難さをいかに強調しても過言ではない。私はこれを根本的な分析と再構成によってしか解決できないと考えるからである。

空間、時間、物質、因果関係はこれらの概念の主なものである。まず因果関係という概念から始めよう。

因果関係は、先ほど述べたように、その表面的な意味で受け入れる危険性の高い概念です。どのような事象に対しても、その事象を引き起こす何かが存在すると考えられています——特定の明確な出来事があり、その出来事がなければ事象は不可能であり、その出来事が存在すれば事象は必要不可欠となります。事象は、その原因によって一定の方法で依存すると考えられていますが、この依存性は他のものに依存していません。したがって、人々は心は脳に依存するか、あるいは逆に脳は心に依存するかを主張します。十分な知識があれば、人間の心の状態からその人の脳の状態を推測したり、その人の脳の状態からその人の心の状態を推測したりできるかもしれません。通常の因果関係の概念が保持される限り、この状況は実在論者が脳の状態が思考を引き起こすことを主張し、理想論者が思考が脳の状態を引き起こすことを主張するのに利用されます。どちらの主張も等しく正当か等しく不正確かです。事実上は、多くの因果関係のような相関関係があり、例えば物理的事象や心理的事象は、十分な数の物理的前駆体からあるいは十分な数の心理的事象から理論的には予測できると思われます。何か事象の原因を言及するのは誤解を招くでしょう。事象を理論的に推論するための適切な前駆体のセットであれば、その事象の原因と呼び得ますが、何か事象の原因を言及するのは唯一性を示すものであり、それが存在しないことを意味します。

この文章は、私たが「太陽を見ること」という経験と関連していることが明らかです。私たちが見ることを説明するための鎖の一連の前因があり、それによって視覚は眼、神経、脳に依存することがあるとしても、物理的事物としての眼、神経、脳を無視した別の一連の前因があることを示す証拠にはならないのです。私たちは、太陽を見ることを述べる際に生じたジレンマから逃れるためには、少なくとも理論的には、物理世界における因果法則を述べる方法を見つけなければなりません。その方法では、単位は眼や神経、脳などの物質的なものではなく、私たが太陽を見ているときの瞬時の視覚的対象と同じような瞬時の個体となります。太陽自身と眼、神経、脳は瞬時の個体の集合として見なされるべきです。

物理的に受け入れる自然な態度から出発し、物質が物理世界における「本当に実在する」ものであると考え、感覚の直接的な対象はただの幻影であると仮定することが一般的であるとしても、私たちは物質を論理的構築として見なすべきであり、その成分は観察者が存在するときには感覚データとなるような瞬時の個体が含まれるでしょう。物理学が八分前までの太陽とみなすものは、異なる時間に存在し、光速で中心から広がり、現在太陽を見ている人々によって見られるすべての視覚的データを含む一連の全体となります。したがって、八分前の太陽は一連のクラスであり、私は現在太陽を見たときに見ることのできるものはそのクラスの一員です。このクラスを構成するさまざまな個体は、中心から外へと進むにつれて一定の継続性と一定の内在的な変化法則を持つことと、他の非該当の個体と関連した一部の変更が存在します。これらの外部的変更が以前の説明で眼や神経が太陽の見せ方をどのように変えるかという影響として現れたような事実を表すものです。

The \textit{prima facie} 難題は、空間の過度な形式主義的な理論から主に生じています。最初の見方では、世界をこれほど詰め込む必要があるかのように思われることがあります。私と太陽との間に存在するすべての場所において、特定のものは数分前と同じ太陽の一部であるべきです。もちろん、その場所からどの惑星や固定星が見えても、それらのいずれも同じ場所に一部となる特定のものが必要です。私はいる場所では、現在私が感知している「ものの」それぞれの一部となる特定のものが存在します。したがって、世界中で、あらゆる場所において、同じ場所に非常に多くの特定のものが共存することになります。しかし、これらの問題は、学校教師が私たちに慣れさせた単なる三次元空間に甘んじることから生じます。実際の世界の空間は六次元の空間であり、これに気づくと、必要なすべての特定のものに十分なスペースがあることがわかります。これを理解するためには、物理学の滑らかな空間から、直感的な経験の粗末で散漫な空間に戻る必要があります。一人の人間の感覚的対象の空間は三次元の空間です。二人が同時に同じ感覚的対象を見たり聞いたりすることはおそらくないでしょう；彼らが同じものを見たり同じ音を聞いたとされる場合でも、常に微細な違いがあります。これが真実であり、空間における位置は一般的に相対的なものであるという仮定が正しいとするなら、一人の人間の感覚的対象の空間と別の人の感覚的対象の空間には共通の場所はありません。つまり、異なる人間が感知する感覚的空間の部分間の直感的な関係は、同じ人間が感知する感覚的空間の部分間の直感的な関係とは異なります。したがって、世界には多数の三次元空間があります：それらは観測者が感知するものであり、おそらく観測者が適切な位置にいないため感知されないものも含まれます。

これらの空間は互いに同じ種類の空間的関係を持つとは限らないが、それらを三次元的な順序で配置することは可能である。これは、一つの物理的事物のメンバー（または側面）として見なされる相関する個々の要素によって達成される。例えば、「同じ物体を見ている」と言われる人々について言えば、近い位置にある人々はその物体をより大きな視野領域で見ることになる一方、遠くにいる人々はそれに匹敵する部分を見ることができない。このような考慮に基づいて、すべての異なる空間を三次元的な係列として配置することは可能である。各空間が自身の三次元的空間を持つことから、個々の具体的な存在全体は六次元空間の中で配置される。つまり、任意の特定の具体的な存在の位置を完全に指定するには、その存在自体の空間における位置を3つと、他の空間との相対的位置を3つ必要とする。

具体的な存在を分類する方法は二つある：一つは「視覚」に関連するすべてのもの、または常識的に言えば同じ「物」の異なる「側面」とされるすべてのものを一緒に取ることである。例えば、「私は太陽を見ている」と言われる場合、私が見ているものは次の2つの集合に属する：（1）現在の感覚対象全体の集合、これが私にとっての「視覚」である；（2）8分前までの太陽の異なる具体的な側面を指す集合、これは私が定義する「8分前の太陽」となる。したがって、「視覚」と「物」は具体的な存在を分類する二つの異なる方法に過ぎない。注意すべきは、具体的な存在がこの双方向の分類を受け入れる必要性はあらかじめ決定されているわけではないということである。ある種の「野生的」な具体的な存在がある可能性があり、通常の関係によって分類されるものとは異なる関係を持つ具体的な存在が含まれているかもしれない。たとえば、夢や幻覚はこの意味での「野生的」な具体的な存在から成るかもしれません。


視覚とは何を指すのかの正確な定義は容易ではありません。見えている物体や触れる物体について限定すれば、「とえられた個体と単純な空間関係を持つすべての個体」を「視覚」と定義できます。私は今見ている色の斑点同士には直接的な空間関係がありますが、異なる人によって見られる色の斑点同士には直接的な空間関係はなく、物理空間（これは視覚全体から構成される空間と同じ）における「物体」配置を通じて間接的に空間関係が構築されます。とえられた個体と直接的な空間関係を持つ個体は同じ視覚に属します。しかし、例えば私は聞く音が私が見ている色の斑点と同一視覚に属する場合、直接的な空間関係を持たない個体もその視覚に属することがあります。一つの観察者による一時的なデータすべてを視覚として定義することはできません。なぜなら、視覚が誰にも感知されていない可能性があることを考慮に入れる必要があります。したがって、視覚を定義する際には心理学や空間から導き出される原理ではなく、ある原理が必要となります。

そのような原理は時間の観点から得られます。包摂的な時間と包摂的な空間同様に、包摂的な時間も構築されたものです。私の視覚と他者の視覚に属する個体間には直接的な時間関係はありません。一方で、私が認識している任意の二つの個体は同時であるか連続的であるかであり、その simultaneity または succession が私にとってデータとなることもあります。したがって、「とえられた個体と同一視覚に属するすべての個体」と定義できます。「simultaneous」は直接的な単純な関係を指し、物理学における間接的な構築された関係とは異なります。相対性原理によって提案された「ローカル時間」の導入は、科学的理由により我々が先ほど主張したような時間の増加を実現しました。

すべての個別の事象がとえられた特定の事象と同時にまたはそれ以前またはそれ以后に存在するもの全ての総和を、その特定の事象が属する「生譜」と定義できる。観察される必要はなく、同様に生譜も誰かによって実際に生きる必要はない。誰にも生きられない生譜は「公式」と呼ばれる。

「物」の定義は連続性と一定の微分独立を持つ関係を通じて達成される。つまり、一つの視点においてとえられた特定の事象に対して、隣接する視点では通常非常に類似した特定の事象があり、これはとえられた特定の事象から最初に小さな量で異なるものであり、その差は視点空間における二つの視点の位置差のみに関連し、宇宙中の他の「物」には関係ない。この連続性と変化を移行するにつれて現れる法則の微分独立が、「一つの物」と呼ぶべき特定の事象のクラスを定義する。

一般的に言えるのは、物理学者は個別の事象を「物」として分類することが便利であり、心理学者はそれらを「視点」および「生譜」に分類することが便利である。なぜなら一つの視点は一瞬の観察者のデータを構成し得る一方で、一つの生譜は一生を通じて一観察者の全データを構成する可能性があるからである。

私たちの議論を要約することができる。我々の目的は、可能な限り物理世界の最終的な構成要素の性質を探求することだった。私が「物理世界」という言葉を使用した場合、最初に物理学で扱われる世界を指すものだと理解してほしい。物理学は明らかに経験的科学であり、感覚を通じて得られる証拠に基づいて一定量の知識をとえてくれる。しかし、物理学自身の発展や生理学、心理学、或いは形而上学から導かれた議論によって、物理世界の最終的な構成要素が直感データそのものではなく、「精神的」、「心の中に存在する」、「主観的な」というような何かに属していると考えられるようになった。この見解の根拠が物理学に依存する部分については、象徴論理に基づく複雑な構築によって適切に対処されるべきであることが明白だ。感覚から提供される材料を使って物質にとえられた性質を持つクラスと係列を構築することが可能であることを示すことができる。この議論は困難で技術的なものであり、本文ではそれに取り組むことはしなかった。しかし、直感データが「精神的」であるという見解が生理学、心理学、或いは形而上学に依存する部分については、恒常性の偏見と混乱が存在することを示そうとした。物質の最終的な構成要素における恒常性の偏見、空間に関する単純な概念からの混乱、感覚器官との因果関係から生じる感覚データの混乱、そして直感データと感覚との区別欠如がその根拠となっている。これらの主題について述べたものが正当であるならば、直感データの存在は論理的には心の存在に依存せず、観察者の体に対して因果的に依存していることになる。観察者の体に対する因果的な依存は表面上よりも複雑であり、すべての因果的依存と同様に誤った信念を生む可能性がある。もし我々の主張が正しければ、直感データは物理世界の最終的な構成要素のうち偶然に直接認識しているものだけである。それ自体は純粋に物理的であり、それに関連するすべての精神的なものは我々の認識であり、それはその性質や物理学における位置に関係なく無関心である。

不適切に単純化された空間に関する概念は、実在論者たちにとって大きな障害となっています。二人の男が同じテーブルを見ている場合、一人が見ることと、もう一人が見ることは同じ場所にあると仮定されます。しかし、二人の視覚は形状や色において全く同じではないため、これは困難を引き起こします。この困難は、それぞれが見るものは単なる「主観的」であると断言し、それによって素早く解決されたり、むしろ隠蔽される傾向があります。しかし、「主観的」という言葉を使う人々がその意味を何であるか説明できるかどうかは疑問です。実際のところ、空間—そして時間も同様に—は物理学の完成された構造から見るとより複雑であり、一つの包摂的な三次元空間は六次元空間からの関連性によって得られる論理的構築であることが明らかになります。この六次元空間に存在する個々のものは、一種の方法で分類されると「物」を形成し、別の方法で分類されると「視覚」や「経歴」を形成します。これらの個々のものがそれぞれ瞬間的な経験の感覚データや全体的な経験の感覚データとなるためには、適切な観察者が存在する必要があります。物理的「物」が係列として解体され、我々が行ったように特定のクラスの個々のものに分けられるまで、物理学と心理学の視点との間の衝突は克服できません。この衝突は、述べた内容が誤りでない限り、異なる分類方法から生じるものです。その源を発見するとすぐに消失します。

理論を簡潔に提示した私の主張については、絶対的な真実であるとは主張していません。誤りの可能性や仮説的要素が含まれている部分もあります。しかし、この理論は他のどの理論よりも真実である可能性が高いと主張します。特にライプニッツの類似理論を除いては、これほど言えるものはありません。現実主義に伴う困難、物理学に対する哲学的な説明を妨げる混乱、感覚データを否定する結果とするジレンマが残る外世界に関する知識の唯一の源泉——これらすべては、私は推奨する理論によって避けることができます。この理論が真実であることを証明するものではありません。おそらく他の多くの理論も同様の長所を持つかもしれませんからです。しかし、この理論は現在の競争相手よりも真実である可能性が高いという点で優れていることが証明され、そして何が確実に知られるかについては、その理論を出発点とし、無関係、不要、または根拠のない仮定から徐々に解放することで、それらが明らかになる可能性があることを示唆します。これらの理由から、私はこの理論を仮説として注意を集め、さらなる研究の基礎とするために推奨します。しかし、その自身が問題解決の完全なまたは十分な解答ではないという点は認識しておくべきです。

脚注：

23. 1915年2月にマンチェスター哲学会に対して発表された講演を再録したものです。

24. サミュエル・アレクサンダーコーの「現実主義の基礎」（ブリティッシュ・アカデミー、第6巻）を参照してください。

25. 「感覚的性質は感知に依存するか？」（アリストテレス協会議事録、1909-10年、pp. 191-218）

26. ヒラスとフィロノウスの最初の対話（フレーザー編、1901年版、第I巻、p.384）

27. アメリカの現実主義者はこの点をよく指摘しています。

28. T.P. ヌンの「感覚的性質は感知に依存するか？」（アリストテレス協会議事録、1909-10年）を参照してください。


\chapter{第VIII章：センス・データと物理学との関係 --- {THE RELATION OF SENSE-DATA TO PHYSICS}}

VIII

感覚データと物理学の関係


I. 問題の提起

物理学は観察と実験に基づく経験的な科学だとされる。

それは予め計算可能であり、その後観察と実験によって確認される結果を提供することが可能であると仮定されている。

観察と実験から何を学べるか？

物理学に関しては、何も学べない。それ以外のものとして、直感的な感覚データ：特定の色の斑点、音、味、臭いなど、一定の空間時間関係を持つものだけである。

物理世界の仮定される内容は最初に見えて非常に異なる：分子には色がなく、原子は音を立てず、電子には味がないし、粒子すら匂わない。

このような物体が検証されることは、感覚データとの関係性を通じてのみである。それらは感覚データと何か種類の関連を持つべきであり、その関連性だけで検証されるべきである。

しかし、その関連性自体をどのように確認するか？関連性は常に共に見つかった被関連物体によって経験的に確認される。しかし、私たちの場合、関連性の一方の項のみが常に見つかることになる：すなわち感覚的な項である。もう一方の項は本質的に見つかることができないかのように思われる。したがって、物理を感覚データとの関係で検証するためには、その関連性自体は永遠に不確実であり、絶対に確認できないことになる。

この結果を避ける二つの方法がある。

（1）私たちは直観的に原理を知っていると主張し、経験的な検証の必要がない場合がある。例えば、私たちの感覚データには自己以外の原因があり、その原因についての知識は効果から推論することによって得られるべきである。この方法はしばしば哲学で採用されてきた。これが必要な程度であれば採用してもよいが、物理学が経験と観察のみに基づくものではなくなるのであれば、可能な限り避けるべきである。

（2）物理の対象を感覚データの関数として実際的に定義することができる場合がある。物理学が予測を導き出す程度であれば、これは必ず可能である。なぜなら、私たちは経験できるものだけを期待できるからである。また、物理状況が感覚データから推論される限り、それは感覚データの関数として表現可能であるべきである。この表現を行う問題は興味深い論理数学的な作業につながる。

物理学においては、通常のように知られているところでは、感覚データは物理的物体の関数として現れる：ある種の波が目に入ると、そのような色を見ることになるなどと主張される。しかし実際には、波は色から推測され、逆ではない。物理学が経験的数据に基づくものとして正当化されるのは、波を色その他の感覚データの関数として表現した後である。

従って、物理学が検証可能であるとするなら、我々に以下の問題が残る：物理学は感覚データを物理的物体の関数として示すが、検証は物理的物体を感覚データの関数として示せる場合のみ可能である。したがって、感覚データを物理的物体の関数として表す方程式を解く必要があり、それによって代わりに感覚データを物理的物体の関数として表すようにする。

II. 感覚データの特性

「感覚データ」という用語を使用する際には、私が指しているのは一時的な感覚でとえられるものの全体ではなく、その一部であることを意味する。特定の部分の色、特定の音など、注意によって選り分けることができるような部分を指す。一つの感覚データとして何が考慮されるべきかを決定することは難しい：しばしば注意は以前には存在しなかった分割を生み出す。例えば、この赤い領域がその青い領域の左にあるという観察された複雑な事実も、現在の立場から見ればデータとみなすことができる：認識論的には、単純な感覚データと比較してそのとえる知識に対する機能的な違いは大きくない。しかし、その論理的な構造は感覚とは非常に異なるものである：感覚は個体との親しみを提供し、これは二項関係であり、対象は名前でしか指すことができず、真偽の根拠がない一方で、複雑な事実の観察、適切には認識と呼ぶべきもの、は二項関係ではなく、主張形式を含み、真理についての知識をとえ、単なる個体との親しみとは異なる。この論理的な違いは重要であるが、現在の問題に関連しないので、本稿のために感覚データとして認識データを含むことに便利である。特に注意すべきは、認識データの構成要素は常に厳密な意味での感覚データであるという点である。

Concerning センスデータ，我々はそれらがデータであるときには存在することを知っています。これはすべての外部的個別の知識の認識論的基礎です。（もちろん、「外部」という言葉の意味は後で扱う問題を引き起こします。）我々は、感覚データが一度データであった物体が感覚データではなくなる時間に引き続き存在するかどうかについては、より或る程度不安定な推論によってしか知ることができません。感覚データがデータであるときのそれらはすべて、我々が直接原始的に外部世界を認識するものであり；したがって、認識論において、それらが \textit{データである} という事実は全て重要です。しかし、それらが我々が直接認識する唯一のものであるという事実は、もちろん、それらが存在するすべてのものであるとは証明しません。私たちは個人的な形而上学を構築できる場合、知識と無知の偶然性に依存しない独立した形而上学があれば、実際のデータの優遇された位置はおそらく消失し、それらはおそらくより類似する多くの物体からなる一連の中から乱雑な選択として現れるかもしれません。このようなことを述べるとき、私はただ認識されていない特定の個体が存在すると推定しているだけです。したがって、感覚データの特別な重要性は認識論に関係していますが、形而上学には関係ありません。この点で物理学者は形而上学に数えられるべきであり、それは個人的ではなく、感覚データに特別な注意を払っていないと名指します。しかしそれらの重要性は、物理学がどのように \textit{認識されるか} という問いが出るときに出直ります。

以下が翻訳です：

感覚的対象を「sensibilia」と呼び、感覚データと同じ形而上学的・物理学的な地位を持つもののことを指しますが、必ずしもどの心に感知される必要はありません。したがって、「sensibile」に対する関係は人間に対する夫の関係に似ています：人間は結婚によって夫となることができ、同様に「sensibile」は認識の関係によって感覚データとなります。両方の用語を有することが重要です。なぜなら、ある時感覚データであり、別の時感覚データでない物体が依然として存在できるかどうかについて議論したいからです。感覚データが感知されずに存在できるかという質問は「夫が結婚されていない状態で存在できるか」と同じようなものなので、「sensibile」が感知されずに存在できるかという質問と、特定の「sensibile」がある時感覚データであり、別の時ではそうではないかという質問が必要です。これらの用語がないと、このような質問は軽微な論理的なパズルに陥る可能性があります。

すべての感覚データは「sensibilia」となります。「sensibilia」がすべての感覚データであるかどうかは形而上学的問題であり、特定の「sensibile」を推測するための方法が感覚データから得られるものだけであるかどうかは認識論的問題です。

進行するにつれて補足的ないくつかの前説は、「sensibilia」という用語を使用する目的を明確にするのに役立ちます。

感覚データは精神的ではなく、実際の物理学の対象の一部であると私は考えます。感覚データが主観的なものであるという議論は短時間に検討される予定ですが、これらの議論は私には、つまり生理学的な主観性、すなわち感覚器官、神経、脳に対する因果的依存性を証明するだけだと考えられます。何かが私たちに対して示す現象は、遮断物や霧や色ガラスなどの中間のものと同様に、その感覚器官、神経、脳による因果的依存性があります。両方の依存性は、ある場所から物体を観察したときに示される現象が、物体だけでなく中間の媒体にも依存する関数であるという声明の中に含まれています。（この声明で使用される用語「物質」、「特定の場所からの視点」、「現象」、「中間の媒体」は本稿を通じて定義されます。）脳や神経、感覚器官を取り巻かない場所から何かがどのように見えるかを確認する手段はありませんが、連続性により、そのような場所では何かしらの現象が存在すると仮定するのは不合理ではありません。そのような現象はすべて「感覚的もの」に含まれます。

もし不可能なことでもありますが、脳や神経、感覚器官を取り巻かない完全な人間の体があった場合、その体に対して存在するすべての「感覚的もの」が存在し、体内に意識がある場合の感覚データとなります。実際には、意識を加えることはただの認識であり、それ以外はすべて物理的なものや生理学的なものです。

IV. 感覚データは物理学的である

この質問について議論する前に、「精神的」および「物理学的」という用語の意味を定義することが適切です。すべての前段階の議論において、「物理学的」という言葉は「物理学によって扱われるもの」を意味することとします。物理学は明らかに、実際の世界の一部の構成要素について何かを教えてくれます；これらの構成要素が何であるかは不明ですが、それが物理的なものであり、その性質がどのようなものであれ、それらを物理的と呼びます。

以下の原著英文を忠実に日本語へ翻訳しました：

「精神的な」という用語の定義はより難しく、多くの困難な論争が討議され決定されるまで適切にはとえられません。そのため現在の目的のために私は、これらの論争に対する断定的な答えを仮定せざるを得ません。私は何かに意識している特定のものを「精神的な」と呼び、そのような特定の要素を含む事実を「精神的な」事実と呼びます。

精神的で物理的なものは必ずしも互いに排他的ではないが、重なり合う理由がないことは知っています。

私たちの現在の議論における「精神的な」定義の妥当性に対する疑問は、あまり重要ではありません。私が主張したいのは感覚データが物理的なものである点であり、これが認められた場合、それは現在の調査において感覚データがまた精神的なものであるかどうかに関わらず、私たちにとって無関係です。マッヒ、ジェイムズ、そして「新しい実在論者」たちと異なり、私は精神的で物理的な差異は単なる配置の違いであるとは主張していませんが、現在の論文において述べること自体は彼らの学説と一致し、その立場から導き出すことも可能でした。

感覚データに関する議論では、一般的に混同される2つの質問があります：

1. 感覚的対象が私たちがそれを感じていないときにも存在するのでしょうか？つまり、ある時間点でデータとなった感覚的要素が、その時点を過ぎてデータではなくなるときでも存続するのでしょうか？
2. 感覚データは精神的なものか物理的なものですか？

私は感覚データが物理的なものであると主張しつつ、それらがデータとして存在しなくなった後も常に変化なく存続することはないだろうという可能性を維持します。存在しないことを主張する見解はしばしば誤って、そのような視点から感覚データが精神的なものであると推測されることがありますが、私はこれが現在の問題に対する混乱の一因となったと考えています。もし感覚データがデータとして存在しなくなった後でも存続することは論理的に不可能であるという主張があったとしても、それはそれ自体に精神的な意味を持つことはありません。しかし私が主張するように、その非存続は経験的に確認された因果法則から導かれる可能性のある推測であり、そのためそのような意味を持たないのです。そして私たちは物理の対象として感覚データを扱うことができます。

論理的には、感覚datumは対象であり、その特定のものについて主観が認識しているものです。 datumには主体が部分として含まれるような信念や意志とは異なります。したがって、 datumの存在は主体の存在に論理的に依存していないことがわかります。私が知る限りでは、 $A$ の存在が $B$ の存在に論理的に依存する唯一の方法は、 $B$ が $A$ の部分である場合です。したがって、特定のものがあなたがたが datumとして存在しなくなった後でも persistenceし続ける理由や、他の類似の特定のものが datumとして存在することなく存在する理由はありません。

感覚datumが精神的であるという見解は、その生理的な主観性から一部派生していることは確かですが、それも「感覚」という概念と感覚datumを区別しないことから来ている部分もあります。私は感覚とは、主体が感覚datumに認識する事実を指します。したがって、感覚は主体がその一部分である複合体であり、従って精神的です。一方で、感覚datumは主観に対して対立し、感覚において主体が認識している外部の対象です。感覚datumはしばしば主体の体内に存在しますが、主体の体は主体から独立しており、実際には物質世界の一部であるという点を考慮に入れると、感覚datumが物理的なものと見なされる主要な障害は、感覚datumと感覚を明確に区別し、その主観性が生理的ではなく心理的なものであることが認識されたときに除去されます。

V. "センシビリティ"と"事物"

ただし、「感覚的対象」が物理的世界の最終的な構成要素として認められるなら、共通感覚や物理学の「物質」との到着までには長い困難な旅が必要となる。同じ「物質」の異なる「感覚的対象」を異なる人々が観測する異なる感覚データを組み合わせることが不可能とされる仮定により、「感覚的対象」は単なる主観的な幻影であるべきだと考えられることになった。ある人にはテーブルが長方形の形を呈し、別の人は二つの鋭角と二つの鈍角を持つ形を呈する。ある人には茶色に見える一方で、光を反射する方向からは白く輝いて見える。これらの異なる形状や色は同時に同じ場所に存在することは不可能であり、したがって物理的世界の構成要素ともなるべきことはないと主張されることがある。この議論は私には最近まで不破の論法のように思えた。しかし、Dr. T.P. Nunn の記事「二次的質的特性は観察から独立するか？」において反対意見が巧みに主張されていることが判明した。その仮定された不可能性は「同じ場所」のフレーズから得られた力強さを持っているが、その弱さはこのフレーズに存在する。空間概念はしばしば哲学で（特に反省によってそのような扱いを防ぐべきと自覚している人々でも）、カントが心理的な無知から想像したように、とえられたものとして、単純なものとして、明確なものとして扱われる傾向がある。観測されていない「場所」の曖昧さこそ、実証主義者に困難をもたらし、その対立者に不必要な優位をとえる原因となっている。感覚データには異なる種類の二つの「場所」が含まれており、それはデータが現れる場所と、データが観測される場所である。これらは異なる空間に属するが、私たちは後に見るように、一定の制限のもとではそれらを関連付けることが可能である。同じ物質に対する異なる観察者のそれぞれの「異なる観え方」は観察者が関わる私的な空間でそれぞれ存在する。一方の観察者の私的な世界における任意の場所は、別の観察者の私的世界における任意の場所とは同一ではない。したがって、これらの異なる観え方は同じ場所に組み合わされる問題は生じない；それらすべてが同時に一つの場所に存在することはできないという事実は物理的な現実性を疑う理由にはならない。共通感覚の「物質」は実際にはその全ての観え方と同一視できるかもしれないが、しかし観察者がデータとして経験していない「感覚的対象」も含むべきであるという条件付きでなければならない。


例を挙げてこの点が明確になるでしょう。部屋の中にいる人々が、彼らの言うところの同じテーブルや椅子、壁や絵画を見ているとします。これらの人々の感覚データは完全には一致しませんが、そのデータには十分な類似性があり、それらをいくつかの「物」の観測者たちに対する姿としてグループ化することが可能です。また、部屋の中にある特定の物体が実際の観測者たちに対して示す姿とは別に、他の可能性のある観測者がその物体に対して見ることのできる別の姿があると仮定します。ある人が二人の他の人間に座る位置に座ると、彼には二人の他の人々に対する部屋の姿の中間的な姿が現れるでしょう：しかし、この姿は新たに到着した観測者の感覚器官、神経、脳がないと存在しないとしても、彼が現在占める位置から見れば、部屋の何らかの姿は彼の到着以前に存在していたと考えるのは自然なことです。しかし、この仮定はただ指摘するだけで強調することはありません。

「物」を無防備な偏りで特定の一つの感覚データと同一視することはできませんから、それにはすべての感覚データとは異なる何かが存在し、その下位にあると考えられるようになりました。しかしオカムの剃刀の原理によれば、共感的な思考に至った原始的メタ物理学者たちによって「物」が作られた目的を果たすためには、感覚データのクラスが十分であるとすれば、経済性から物をその感覚データのクラスと同一視することが求められます。これらの感覚データの下位にある実体を否定する必要はありませんが、この不要な存在を主張することを避けることが便利です。ここでの手順は数学の哲学において無用なメタ物理学的モンスターから解放されたものに比喩的に類似しています。

VI. 設計図と推論

「場所」という言葉の曖昧さを分析し、説明する前に、いくつかの方法に関する一般的な注記が必要です。科学的な哲学における最高の原理は次の通りです：

  推論された存在を代替するために、可能な限り論理的構築が使用されるべきです。

数学哲学の領域において構築を推し進める代わりに推論を行ういくつかの例は、この最大公約数の用途を明瞭にするのに役立つでしょう。まず最初に無理数について考えてみましょう。昔は無理数が、有理数列の極限として存在するが、その極限が有理数でないという推測から導かれていましたが、この手続きには、無理数の存在を単なる推量に過ぎないとする欠点がありました。そのため現在の厳格な方法では、このような定義は容認されません。今では無理数を特定の比のクラスとして定義し、比を通じて論理性で構築します。これにより、疑わしい推論から導かれるものではなく、無理数が存在することを確実にします。

また、基数について考えてみましょう。等量の集合には何か共通点があるように見えます。その共通点は、集合から推論して得られる基数とされていますが、基数が集合から推論される限り、その存在は疑問が残ります。これは、任意の仮定的な形而上学的前提を用いて補完する必要があります。特定の集合の基数を、等量のすべての集合のクラスとして定義することで、この形而上学的前提を避けることができます。これにより、算術の形而上学的な不確定要素を取り除きます。

また私は他處で示したように、類似の方法は自己そのものにも適用できます。これらは仮に形而上学的な実在性を持つ必要がないとされ、記号的に構築された架空のものとして扱うことができます。

これらの構築法の進行方法は、これらの場合やすべて類似した場合においても密接に類似している。仮定された推論対象について形式的に扱われる一連の命題がとえられたとき、これらの命題を真にするために仮定された対象に求められる性質を見つける。少々の論理的な巧妙さによって、それらの仮定的対象より少ない仮定的な対象の論理関数を作り出し、その関数が求められる性質を持つことを確認する。この作成された関数を仮定された推論対象に代入することで、命題全体に対する新しい且つ疑いの余地のない解釈を得る。この方法は数学の哲学において豊かであることが知られており、物理学の哲学においても同様に適用可能であることが見出されるだろう。私は物理学へのこの手法の適用に関して独自性を主張することができないため、その提案と刺激をとえてくれた友人で共著者であるドクター・ホワイトヘッドに感謝しなければならない。彼は感覚データと物理学の点、瞬間、粒子の間の領域におけるより数学的な部分への適用に取り組んでいる。

以下の原著英文を忠実に日本語へ翻訳しました：

方法を構築代入法とし、完全に適用すると、物質は感覚データのみで表現される形になります。さらに、他の人の感覚データも含まれるかもしれません。他者の感覚データを知ることは必ずしも推論なしでは不可能ではありませんが、その要素が含まれます。しかし、この理想は現在のところ現実には遠く、長期的な前準備作業が必要であり、それが完全に達成されるのは非常に難しいです。避けられない推論は、いくつかの指導原則に従って処理できます。まず第一に、常に完全に明示的でなければならず、可能な限り一般的な形で述べるべきです。第二に、推論された存在体は、それが可能であれば、とえられたデータから遠く離れたもの（カントの「Ding an sich」のようなもの）ではなく、他の感覚データと類似性を持つべきです。私は許容する推論された存在体は二種類あります：（a）他の人の感覚データ、証言の証拠として最終的には他者の心の類推的議論に基づくもの；（b）知恵が存在しない場所から明らかになる「sensibilia」で、これらを実在と仮定しますが、誰のデータでもありません。これらの二つの推論された存在体のクラスは、おそらく最初のものは挑戦されないでしょう。私はそれを排除できることに最大の喜びを感じますが、ソリシズムの基礎上物理学を確立することができるとすればです。しかし、人間的な感情が論理経済よりも強い人々（恐らく大多数）は、ソリシズムを科学的に満足させるために私の望みとは異なるでしょう。第二種類の推論された存在体は、より深刻な問題を提起します。ある場所で感覚器官や神経構造が存在しない場合でも何かが現れるという主張は、異常であると考えられるかもしれません。私自身もその異常性を感じませんが、これらの仮定的な現象は仮説的枠組みとしてのみ見なすべきであり、物理学の建築物が完成するまで使用されるべきで、可能であれば完成後には取り除くことができるかもしれません。誰のデータでもない「sensibilia」は、説明のために用いる推論的な仮定とし、物理学の最終形態における道徳的部分とは見なされません。

VII. 私的空間と視点の空間

「場所」という言葉における曖昧性について説明しなければならない。これは、同一の感覚データが二人の異なる人々に同時にとえられることがないという事実から始まる。二人の異なる人が見ることのできるものはしばしば非常に類似しており、その類似性は同じ単語を使用してそれらを指すことができ、これなしでは他の人との間での感覚的物体に関する通信は不可能となる。しかし、それでも視点の違いから生じる差異が常に存在するように見える。したがって、各人は自らの感覚データに関していうと、私的な世界に住んでいることになる。この私的な世界にはその独自の空間があるか、あるいはそれよりも多くの空間があるかのように思われる。しかしそれは視覚の空間と触覚の空間を含む他感覚の様々な空間と関連付ける経験によって教えられるものである。しかし、これらの多様な私的な空間については、心理学者にとって興味深いものの、現在の問題に関しては大きな重要性を持つことではない。なぜなら、単に自己中心的な経験が各々の感覚データを一つの私的な空間に統合する能力をもたらすからである。感覚データが存在する場所は私的空間における場所である。したがって、その場所は他の観察者の私的空間におけるどの場所とも異なるものとなる。なぜなら論理経済性の要求に従ってすべての位置が相対的なものであると仮定するのであれば、場所はその中や周囲にあるものの関係によって定義されるため、二つの共有要素を持たない私的世界には同じ場所は存在しないからである。したがって、同一の場所で同一の物体の異なる観察を組み合わせるという問題は生じず、特定の対象が異なる観測者によって異なる形と色に見えるという事実は、すべてこれらの形状と色の物理的な現実性に対する議論にはならない。

さらに、異なるpercipientの私的世界に属する個人的な空間以外には、一つの私的世界が一点としてまたは少なくとも空間的単位として数えられる別の空間がある。この空間は視点の空間と説明されるべきであり、各私的世界は宇宙から特定の視点からの観測結果として現れるものと考えることができる。しかし私は、「視点」の空間という用語を使用することを好む。これは私的世界が誰かがそれを観察する限りのみ実在すると推測させるような印象をとえることを避けるためである。そして同じ理由から、percipientを仮定せずに私的世界を話すときには、「視点」と呼ぶことにする。

異なる視点が一つの空間の中でどのように順序付けられるかを説明しなければならない。これは、同一の物体の異なる視点での観測結果である「sensibilia」が関連付けられることによって達成される。移動や証言により、二つの異なる視点は両方とも同じ「sensibilia」を含まないとしても、非常に類似したものを含むことが発見できる。特定のグループの「sensibilia」が一つの視点の私的世界における空間的順序が、別の視点の私的世界における関連する「sensibilia」の空間的順序と一致したり、非常に類似していることが判明する。このようにして、一つの視点での「sensibile」は別の視点での「sensibile」と関連付けられる。このような関連した「sensibilia」は「同一物体の観測結果」と呼ばれる。

リーブニッツの単体論においては、各単体が宇宙全体を鏡像として反映するため、それぞれの視点には各物体の観測結果である「sensibile」が存在した。しかし私たちの視点のシステムでは、そのような完全性の仮定を行わない。特定の物体はいくつかの視点で観測結果を持つが、おそらくある他の視点ではそうではない。物体をその観測結果のクラスとして定義すると、ある特定の物体θが出現する視点群κの乗算的なクラスの要素であることが判明し、κは互いに排他的な「sensibilia」のクラスの集合である。同様に、視点もその中で出現するもののクラスの乗算的なクラスの要素である。

空間における視点の配置は、特定の物体が異なる視点で見られるときのその物体の異なる観察結果によって行われます。例えば、あるコインがいくつかの異なる視点で見えますが、一部では大きく見え、他の部分では小さく見え、また一部では円形に見え、他には楕円形に見え、その楕円形は歪み具合によって異なります。

これらの円形に見えるすべての視点をまとめ、コインが円形に見える視点を集めた直線上に配置し、コインの観察される大きさの変化に基づいて並べ替えます。また、コインが一定の太さの線条体として見えている視点も平面に配置します（この場合、同じ大きさで見えている多くの異なる視点があり、これらの視点はコインを中心とする円周上に形成されます）。そして、同様に観察される大きさに基づいて並べ替えます。

このような方法により、コインが視覚的に見えるすべての視点を三次元的な空間的順序で配置することができます。経験から知るところでは、コインではなく他のどの物体を選んでも、あるいは同じ物体が異なる視点で見られるときのその違いを利用したどのような方法を選んでも、同じ空間的順序の視点が得られます。この実証的事実こそが、物理学的な一元的な空間を構築することができた理由です。

先ほど説明したように、要素が全体の視点である空間は「視点空間」と呼ばれます。

 VIII. 視点空間における「事物」および「感覚対象」の配置

我々がこれまで構築した世界は六次元のもので、これは三つの視覚的な係列からなり、それぞれは自身が三次元である。現在、各視覚係列内の様々な私的空間との間の関連性を説明しなければならない。この関連性によって物理学の一三次元空間が構築され、その結果、視覚係列と観察者の私的空間の区別が無意識的に曖昧になり、物理学の哲学に大きな影響をとえた。 pennyを取り戻すことにしましょう： pennyが大きく見える視覚係列では、それより小さいように見える視覚係列よりも pennyの方が近いとみなされます。しかし経験によれば、pennyの見え方の大きさは一定の限界を超えることはありません。具体的には、pennyが目の近くにあるとき、もう少し近づいても見えなくなるという点で停止します。触覚ではこの係列を延長し、 pennyが目につくまで続けることができますが、それ以上はできません。もし以前定義した意味での視覚係列の線を沿って移動してきた場合でも、 pennyを想像的に移動させることにより、別の pennyを使って視覚係列の線を延長することができます。これと同じことは pennyで定義された他の任意の視覚係列の線についても成り立ちます。これらのすべての線は特定の場所で交差し、それは特定の視覚係列です。この視覚係列は「pennyがある場所」と定義されます。

二つの場所がとえられた「感覚的」に関連付けられる方法は、現在明らかになった。まず第一に、その「感覚的」が要素である視点の場所がある。これは「感覚的」が見えたという意味での場所である。次に、その「感覚的」が要素であるものの所在がある、すなわち像であり、これもまた「感覚的」が現れたという場所である。ある視点の要素である「感覚的」は、その「感覚的」が像であるものの所在と関連付けられる別の視点に関連付けることができる。心理学者には「から」という場所の方が興味深いため、「感覚的」は主観的なものとして現れるし、観測者が含まれる場所である。物理学者には「に」という場所の方が興味深いため、「感覚的」は物理的な外部のものとして現れる。これらの二つの見解が互いに不一致である理由、その制限と部分的な正当化は、とえられた「感覚的」に関連付けられるこの二つの場所から明らかになる。

物理的なものに視覚空間における位置を割り当てることが可能であることを以前に述べた。これにより、私たちの身体の異なる部分は視覚空間において位置を獲得し、したがって、私達の感覚データが属する視覚が頭の中に存在することについて言及することは意味を持つ（真偽はあまり気にしない）。私達の心は、私達の感覚データが属する視覚と関連付けられているため、この視覚を私達の心が視覚空間における位置として捉えることができる。したがって、上記で定義された意味での頭の中に存在する場合、心が頭の中に存在することについての主張は良い意味を持つ。これにより、特定の物体に対する様々な現象について、それらのうちいくつかは他のものよりも近いと述べることができる；近いものは、その物体に属する視覚がより近い「物体のある位置」にある場合である。したがって、「物体に関するより多くのことは、近くから調査することによって学べる」という主張には真偽のいずれかの意味を見出すことができる。また、「主題と現象がデータとなる別の物体間の介入するもの」という表現にも意味を見出すことができる。

感覚データの主観性の一理由としてしばしば挙げられるのは、ある物体の現象が変化したときにその物体自体が変化していないかと考えるのが困難な場合である。例えば、目を閉じたり、目を歪めたりして物体が二重に見えるときなどである。物体はその現象のクラスとして定義される（上記で採用された定義）ため、当然、その一つの現象が変化するたびに物体自体にも変化がある。しかし、現象がどのように変化するかによって非常に重要な区別がある。私が何かを見た後で目を閉じると、私の目の現象はそのような視覚におけるすべての視覚空間において変化し、一方で多くの物体の現象は変わらない。定義として述べることができるが、何かが変化したときには、その物体に属するどの程度近い現象でも、それよりさらに近い現象に変化がある場合である。

他方で、ある距離を超えないすべての物体の現象が変わらなかった場合、しかり比較的遠くにある物体の現象のみが変化したときには、その変化は別の何かにあると述べることができる。この考慮から、物質についての議論に自然と導かれる。これは私たちの次のトピックである。

IX. 物質の定義

「物理的事物」をその出現のクラスとして定義しましたが、これには物質という概念の定義としては不十分です。私たちは、ある物体の特定の視点での出現が、その物体と視点間の物質によって因果的に影響を受けているという事実を表現したいと考えています。物体と視点との「間に」何かがあるという意味は見つけましたが、物体全体の出現のクラスよりも物質を別のものとして定義したいと思います。これは、物質が出現に及ぼす影響を述べるためです。

通常、私たちは物体について得られる情報が近ければより正確であると仮定します。遠くから見て人間だと気づき、その後ジョーンズだとわかり、さらに彼が笑っていることがわかります。完全な正確性は限界として達成されるべきものであり：私たちはジョーンズの出現が私たちが彼に近づくにつれて極限に向かう場合、その極限をジョーンズの実態と見做すことができます。明らかに物理学的な視点から見て、近くにある物体の出現は遠くにある物体のそれよりも重要であることがわかります。したがって、次のような仮定的定義を立てることができます：

ある物体の「物質」は、その物体からの距離が減少するにつれてその物体の出現の極限です。

この定義には何かが含まれている可能性がありますが、実証的に感覚データからそのような極限を得ることは不可能であるため、不完全です。定義を補足するために構築と定義が必要になるでしょう。しかし、おそらく正しい方向を見つけるヒントを提供しています。

我達現在位置明確，可以概要理解物理学如何从物質逆向推導至感覚数据。ある物体の特定視点での出現は、その物体の構成物質と中間にある物質によって決まります。物体の出現は煙や霧、青い眼鏡、観察者の感覚器官や神経の変化（これらも中間媒体として考慮する必要がある）によって変わります。物体に近づくほど、その出現は中間の物質による影響が少なくなります。物体から離れるにつれて、その出現は元の特性からますます異なっていき、これらの変化の因果法則は中間にある物質を用いて述べられます。非常に小さな距離での出現は、物体以外の原因による影響が少ないため、これらが距離が縮まるにつれて向かう極限値は「実際のもの」であると考えるようになります。これは因果法則を表すためにも重要であり、物質が感覚データよりも「実在的」であるという誤った感覚の源となっています。

例えば、物質の無限分割性について考えてみましょう。ある物体を見つめて近づくと、一つの感覚データが複数に分かれ、それぞれはさらに分割します。したがって、一つの出現は複数のもの代表し、この過程には終わりはありません。したがって、極限において、物体に限りなく近づくと、有限距離では一つの出現に対応する無限の物質単位が現れます。これが無限分割性が生じる方法です。

あるものの因果効果はその物質に完全に内包されています。これはある種の経験的な事実ですが、「因果効果」を正確に定義するのは難しいため、それを正確に述べるのは困難です。


事物の質料について経験的に知ることは、正確ではない。非常に小さな距離から物の表象を把握することは不可能であり、その表象の限界も正確には推測できないからである。しかし、観察できる表象を通じて、それらは大まかに推定される。そして結果として、これらの表象は物理学によって、私たちの直近の周囲にある質料の関数として表現されることがわかる。例えば、遠くの物体の視覚的な表象は、目につく光波の関数である。これは思考を混乱させるが、実際的な困難にはならない。

一つの表象だけでは、その他の同時的な表象を決定することはできないが、それによって一定の程度の決定がなされる。物の隠れた構造について可能な限りの決定は、精密な動力学的推論を通じてのみ達成される。

X. 時間

全体的な時間は、全体的な空間と同様に構築であると考えられる。物理学自体も、相対性に関する議論を通じてこの事実を認識している。

二人称の経験に属する二つの視点間には、直接的な時間関係、すなわち前後関係がある。これは歴史を異なる経験によって分割する方法と似ているが、経験や精神的な要素を導入しないで分割できる。特定の「感覚的」ものよりも以前または以降、あるいは同時に存在することを定義した「生略記」を用いることで、一連の視点を得ることができる。これらはすべて一人の人間の経験の一部となり得るが、必ずしもそうなる必要はない。この方法によって、世界の歴史は互いに排他的な生略記の数々に分割される。

我々は現在、異なる伝記における時間の関連付けを行う必要がある。一般的には、とえられた（瞬間的）なものの二つの異なる視点での出現が属する異なる伝記において同時であるとみなすことが自然だが、これは便利ではない。例えば、AがBに叫び、BはAの叫びを聞いた瞬間に返事をするとする。この場合、Aが自分の叫びを Heard したときからBの叫びを Heard するまでの間には時間的な差がある。よって、AとBの同じ叫びを Heard した時間を完全に同時としてしまうと、特定の事件と完全に同時であるが互いに異なるイベントが生じる可能性がある。

これを避けるため、「音速」を仮定する。つまり、BがAの叫びを Heard した時間は、Aが自分の叫びを Heard した時間からBがBの叫びを Heard した時間までの半分と仮定される。このようにして関連付けが行われる。

音について述べたことは当然ながら光にも適用される。一般的な原則は、特定の瞬間においてあるものの異なる視点での出現を一緒にグループ化するとき、それらすべてがその瞬間に存在するとみなすのではなく、これらは事物から外側に広がり、その性質により異なる速度で拡散することである。一つの伝記における時間と別の伝記における時間を直接的に関連付ける手段がないため、特定の瞬間においてあるものの出現をグループ化する時間的な分類は一部仮定的である。その動機は主に、同じイベントと完全に同時であるイベントが互いに完全に同時に存在することというような教条の検証を確保し、因果法則の表現における便利さを確保することにある。

第十一節 物や物質の持続性

物理学の世界と感覚の世界を結びつけるためには、以下の三つの主要な問題が生じる：

1. 一つの空間の構築
2. 一つの時間の構築
3. 持続的なものや物質の構築

我々は既に最初と第二の問題を考慮した。残る第三の問題について検討する必要がある。

異なる視点での関連した現象が一瞬の物理的な「もの」を形成する方法について見た。現在は、異なる時間における現象が同じ「もの」として組み合わされる方法と、その結果として得られる恒久的な物理的な「物質」について考える必要がある。

当然、物理学で技術的に基礎となる恒常な素材の仮定は、形而上学的には正当化されない：多くの人々によって同時に見られた一つの「もの」と同様に、異なる時間に同じまたは異なる人々によって見られる一つの「もの」も構築である。実際には、特定の「感覚的現象」の特定のグループだけである。

私たちは、「もの」の瞬間的な状態が異なる視点からの「感覚的現象」の集まりであることを見た。これらは一貫した時間内にすべて同時に存在するわけではないが、「ものがある場所」から速度によって広がる。この状態で「もの」が存在している時間は、これらの現象が発生する時間の下限となる。

今度は、異なる時間における別のセットの現象が同じ「もの」として属する理由について考える必要がある。

この目的のために、少なくとも最初は単一の人生経験に限定することもできる。もしとえられた人生経験内で二つの「感覚的現象」が一つの「もの」の出現であると常に言えるなら、異なる人生経験内の「感覚的現象」を同じ瞬間的な状態の出現として結びつける方法について見た通り、必要なすべてを得ることができる。

まず最初に注意すべきは、常識における「ものの同一性」が物理における「物質の同一性」と常に一致しない場合があることである。人間の体は常識では一つの恒久的な「もの」だが、物理学ではその物質は常に変化している。大まかには、常識的概念は通常の感覚的データでの出現の連続性に基づいており、物理的概念は非常に小さな距離からの出現の連続性に基づいている。おそらく、常識的概念は完全な精密性を有する能力がないと考えられる。したがって、物質の恒久性についての概念に焦点を当てるべきである。


二つの異なる時間における同じ物質の出現の最初の特性は、\textit{連続性}です。二つの出現は、観察によってつながる一連の中間的な存在で結びついています。もし時間と空間が密接な係列を形成するなら、これらの中間的存在もまた密接な係列を形成しなければなりません。葉の色は夏と秋では異なりますが、私たちはその変化が徐々に起こると信じ、二つの特定の時間における色が異なる場合でも、中間的な時間において、それらの時間における色の中間的な色があると考えます。

しかし、連続性に関して重要な考慮事項は二つあります。

まず第一に、これは主に仮定です。私たちは何かを継続的に観察することはほとんどありませんし、観察していない間にそれが観察された状態の間の中間的な状態を通じて移行すると仮定するのは仮説的です。連続性が直接確認できるのは不連続な観察中だけですが、非常に速い運動の場合、例えば爆発では、連続性は直接確認する能力を失います。したがって、私たちは感覚データから仮定的な「センシビラ」の補完が存在し、それが連続性を保つために必要であるという仮説を推測することができます。そしてそのためにはそのような補完がある可能性があります。しかし、すでにこのような仮定的な「センシビラ」を使用していることを考慮に入れると、必要な「センシビラ」を含むこの点は通過させることにします。

第二に、連続性は物質同一性の決定的な基準ではありません。多くの場合、岩や山、テーブル、椅子など、出現がゆっくりと変化する場合には連続性が十分ですが、均一な液体の一部のような他の場合では全く機能しません。海の一滴から別の時間別の一滴へは感覚的に連続的な段階で移動できます。海流の効果から海水の動きを推測しますが、直接の感覚的観察と連続性の仮定だけではその動きを推測することはできません。

以上の文章を忠実に日本語へ翻訳しました：

加えて必要とされる特性は、動力学の法則に従うことです。一般的な感覚から恒常的なものとして認識されるものを出発点とし、時々合理的と思われる修正のみを行ってもたらされる「sensibilia」の集まりが、特定の単純な法則、すなわち動力学の法則に従うことが見つかります。異なる時間における「sensibilia」を同一の物質として考慮することで、運動という概念を定義できます。これは一定期間中の運動中に存在する何かが持続しているという仮定や構築が必要とします。すべての「sensibilia」とその出現時間がとえられた期間において観察される運動は、異なる時間における「sensibilia」を同一の物質として組み合わせる方法により異なります。したがって、世界全体の歴史が特定の点でとえられている場合でも、運動がどのようなものであるかという問いは一定の程度において任意性があります。経験から、動力学の法則を満たすように運動を決定することができること、そしてこの決定は概ね一般的な感覚的な意見と一致することを示します。したがって、このような決定が採用され、異なる時間における出現が同一の物質に属すると見做されるかどうかを実際的または理論的に判断する基準となります。すべての物質が全ての時間において持続されることについては、定義によってその保証ができると考えられます。

以下の原著英文を忠実に日本語へ翻訳しました：

物理学の実験的成功によって結論を推奨するには、その結論が物理学者によって何が証明されているかを考慮する必要があります。証明されるのは、感覚データを超えて仮説が検証不能であるとしても、それらは感覚データと矛盾していない点です。むしろ、十分な「感覚的実在」の集合がとえられたとき、すべての感覚データを計算的に予測できる理想的なものとなっています。物理学では、感覚データを係列に集めることを見積もりました。各係列は一つの「もの」に属すると考えられ、物理法則に関して一般には一つのものに属さない係列とは異なる行動模式を持ちます。二つの表象が同じものに属するかどうかは明確であるためには、表象をグループ化する方法は唯一あるべきです。この点を証明することは非常に困難ですが、現在の目的ではこの点を通過し、一つしかないという仮定を立てることができます。したがって、以下の定義を下記のように設定することができます：「物理学的なものとは、その内容が物理法則に従う係列の感覚的実在である」。このような係列が存在することは実験的事実であり、物理学の検証可能性を構成します。

XII. 幻想、幻視、夢

次に、当システムにおいて、通常の物理学との関連性が欠けたような感覚データはどのように位置づけるべきか問うます。このような感覚データは種類が異なり、異なる扱いが必要です。しかし、これらすべては「虚構」であると呼ばれるものであり、したがって、議論に着手する前に現実と虚構に関する論理的な注記が必要です。

Mr. A. Wolfは言います：

「透明な活動の係としての心像の概念は、夢や幻視の可能性を説明できないため、私はその妥当性が低いと考えます。裸露で透明な活動が存在していないものに向かってどのように指向されるのか、存在されていないものを理解する方法が実現不可能であるように思われます。」

この声明は、おそらく多くの人が賛成するものでしょう。しかし、二つの反論が存在します。まず第一に、いかなる「透明」でない活動も、何も目指すことができないと考えるのは難しいです：関係の項は単なる非実体であってはなりません。第二に、夢象が「存在しない」と「とえられていない」という主張には理由がなく、私はそのような理由を提供することは不可能だと確信しています。まず二番目の点から始めましょう。

（1）夢象が「とえられていない」という信念は、目覚めた生活において感覚的データと対応する「もの」の区別ができていないことから来ていると考えます。夢の中で、夢者が想像しているような対応する「もの」はありません；したがって、もし目覚めた生活で「もの」がとえられているとする（例えばミンゲンが主張するように） なら、夢と目覚めた生活との間にはとえられた性において差異があるはずです。しかし、我々が主張しているように、「もの」は決してとえられるものではなく、その「もの」を構成する感覚的データのうちの一つだけがとえられるなら、夢で我々が理解するものは目覚めた生活で理解するものと同じくらいとえられていると言えるでしょう。

同じような議論は、夢象が「存在しない」という点にも適用されます。夢象は夢者の視点による私的な空間の中で位置付けられています；問題は他の私的な空間との関連性にあり、従って視覚的空間との関連性にあります。しかし、「存在する」というデータとしての意味においては、夢象は目覚めた生活の感覚的データと同じくらい「存在している」のです。

（2）「幻覚」や「虚構」の概念とその補対的な概念である「実在」の概念は、深い論理的な混乱を含む形で一般的に使用されています。対を成す言葉、「実在」と「虚構」、「存在」と「非存在」、「有効」と「無効」などはすべて、「真偽」から派生したものです。「真偽」は、除いて派生的な意味では適用されない限り、主張のみに適用されます。したがって、上記の対を成す言葉が適切に適用される場面では、我々は主張や文脈なしで意味を持たない不完全な表現と対処する必要があります。つまり、「説明」にはこれらの言葉が適用されますが、「名詞」には適用されません：言い換えれば、データに基づいて定義された存在や非存在についてのみ適用されるものであり、データに基づかないデータそのものには全く適用されないのです。


任意のデータ $x$ に対して、「存在」や「非存在」という言葉を使用することは意味がありません。たとえば、我々は「もちろん $x$ は存在する。それ以外ではデータとして存在しないことができないからだ」と言うかもしれません。しかし、そのような主張は実質的に意味がないのです。一方で、「私の現在の感覚データは存在する」や「$x$ は私の現在の感覚データである」という命題は真実であり、重要です。これらの命題から「$x$ は存在する」という結論を導く推論は、未経験の論理者にとっては説得力のあるものですが、その結論自体は単に意味がないだけでなく、完全に無意味なものです。「私の現在の感覚データは存在する」と言うことは、「ある対象が「私の現在の感覚データ」を描写している」という意味です。しかし、$x$ は（仮定の場合）名前であり描写ではないため、「ある対象が $x$ を描写している」とは言えません。ワイトヘッド博士と私は、他處で（同引用）記述し、シンボルの助けを借りてこの点を完全に説明しました；シンボルなしでは理解するのが難しいため、ここではその証明を繰り返すことはありませんが、現在の問題への適用について進めていきます。

「存在」は描写のみに適用されるという事実は、文法的に適切な名詞を使用することで隠されています。これらの名詞は実質的に描写へと変換されます。たとえば、「ホメロスは存在するか」という問いは正当です；しかし、「ホメロス」は「ホメロス詩集の作者」を意味し、これは描写です。同様に、「神は存在するか」という問いも正当であり、「神」は「最高の存在」や「最実在の対象」などと解釈できます。もし「神」が名詞であれば、神はデータであるため、その存在については議論の余地はありません。カントは不明確ながら感じていた存在と他の述語との区別は、描写理論によって明瞭になり、「存在」という概念自体が形而上学の基本的な概念から完全に除外されることを見ることができます。

存在について述べたことは、実在についても同様に適用されます。実在は実存と同義であるかもしれません。幻想や幻覚、夢における直感的対象については、「存在するか」「実在するか」を尋ねることは意味がありません。そこにはあるのですし、それ以上議論することはありません。しかし、これらの対象から推測される「事物」や他の「感覚的対象」の存在性や実在性について尋ねることは正当です。これらの「虚構的な事物」や他の「感覚的対象」が虚構であることに加えて、それらがデータでないことに注意を向けられなかったことが、夢の対象が虚構であるという見解につながりました。

これらを詳細に検討するために、現実主義に対する既存の主張を適用することができます。ただし、以下は主に他人が以前述べたことの繰り返しです。

1. まず最初に、異なる観察者の前で同一物体が異なる形状や色を示すという正常な視覚の多様性があります。これは、同じ物体が異なる観察者に対して異なる形と色を示す場合について言及します。ロックの水は熱くも冷たくもあるように見えるものに該当します。私たちは異なる視点を持つシステムによってこれらの状況を十分に説明できますし、これらは現実主義に対する証拠にはなりません。

2. また、異なる感覚間の相関関係が不通常である場合があります。水中で曲がった棒はその一例です。人々は棒が見え上手く直っていると主張しますが、これは触ると直っているということを意味します。棒が触れると曲がっていると感じるという誤った推論があるだけです。棒が写真に写るときも同じように見えますし、「グロスタス卿によれば、『写真は嘘をつかない』」と述べるように、写真には虚構はありません。重複視もここに該当しますが、この場合の不通常な相関は生理学的なものであり、写真では作用しない可能性があります。重複視したときに「事物」が二つになっているかどうか尋ねるのは誤りです。「事物」は一連の「感覚的対象」で構成され、観察者がデータとする視覚的「感覚的対象」だけが二つに分かれます。この現象は純粋な生理学的な説明で十分であり、私たちが通常注目する物体から得る通常の視覚的感覚データよりも説明が必要ないかもしれません。

(3) 今私たちは、夢のような例に到達しました。これらの夢は、夢中で何か疑念を抱かせるものがないかもしれませんが、以前のデータや後のデータと不一致であるという理由から非難されます。もちろん、しばしば夢の中の物体が通常の行動を示さない場合があります：重い物体が飛ぶ、固体が溶ける、赤ちゃんが豚に変化したり、もっと大きな変化を経験するなど。しかしこれらはすべて不思議な出来事は夢の中で起こる必要はありませんし、そのような出来事が夢の物体が「現実的でない」と呼ばれる理由ではありません。夢者が目覚めると、過去と未来との連続性の欠如からこれらの物体を非難します；そして他の私的な世界とは関連がないという欠如から他人はこれらを非難します。後者の根拠を省略すると、私たちがそれらを非難する理由は、夢の中から推論される「もの」が物理学の法則に従って目覚めからの感覚データから推論される「もの」と組み合わないことです。これは、夢の中から推論される「もの」も非難するために使用できるかもしれません。夢データは確かに「もの」の出現ですが、それらは夢者が想像するような「もの」ではありません。私は心理学者による夢に関する心理学的な理論を攻撃することはありません。しかし、物理的原因が非常に明白な場合もあることは確かです。たとえば、ドアが鳴る音が海戦の夢や戦艦や海や煙の像を生むかもしれません。その全体の夢はドアが鳴る現象の出現ですが、睡眠中の身体（特に脳）の特殊な状態により、この出現はドアが鳴ったときに予想されるものとは異なり、そのため夢者は誤った信念を持つことになります。しかし彼の感覚データはまだ物理的であり、完成された物理学が含み計算するようなものです。


（4）最後の幻覚の種類は、一人の人間の経験の中で発見できず、他の人々との体験との不一致を通じてしか発見できないものです。夢が十分に目覚めた生活と整合性を保つ場合、この分類に属する可能性がありますが、主な例は狂気を引き起こす再発性の感覚幻覚です。このような状況で患者が他人が「狂っている」と呼ぶようになるのは、その人が自身の経験の中で幻覚的な感覚データが他の視点での「感覚的対象」や「物質」に通常の種類の関連性を持つことを示すものが存在しないからです。もちろん彼は証言を通じてこれを学ぶことができますが、おそらく彼は証言が偽であると信じ、自己欺瞞だと考えることを簡単にするでしょう。私には見当たらない限り、患者はこのような状況で自己の狂気と友人の虚偽という二つの同等に満足する仮説間で決定するための理論的な基準はありません。

上記の例から見て取れるように、私たちが欺瞞的とみなす異常な感覚データは、その性質自体において他のどのデータよりも優位性を持たず、他の「感覚的対象」や「物質」との関連性または因果的なつながりにおいてのみ差異があります。通常の関連性と接続性は無意識的な期待の一部となり、心理学家以外にはデータの一部であるようにさえ思われることで、誤ってこれらのデータが虚構であると考えるようになります。しかし実際には、これらは誤った推論の原因に過ぎません。異常な関連性や接続性が発生することにより、感覚から物事を推測する難しさと、物理学を感覚データで表現する困難さが増します。しかしながら、その異常性は物理的または生理的に説明可能であり、従って哲学的な反論ではなく単なる複雑化要因に過ぎません。

私は結論として、感覚データが物理的世界の実質的な一部であるという見解に対する有効な反論は存在しないと主張し、この見解は物理学の経験的検証可能性を説明する唯一のものだと考えます。本稿では、概略的な前駆的な素描のみを提示しました。特に、時間が物理的世界の構築において果たす役割は、上記の説明から見てもより根本的であると考えられます。さらに詳細化すれば、未感知された「感覚的存在」の役割はおそらく無限に小さくできるでしょう。おそらく、「物質」の歴史を用いて瞬時の現象から導かれる推論を補完することで実現されるでしょう。

脚注：

 \textit{Proc. Arist. Soc.}, 1909-1910, pp. 191-218.

 この主題に関しては、Mr. A.A. Robb の \textit{A Theory of Time and Space} (Camb. Univ. Press) を比較してください。彼の理論が最も興味深く新鮮な部分を現在の目的のために省略した上で、彼の見解をここに提案しました。Robb 氏は、同じタイトルを持つパンフレットでその理論の概要を提示しています (Heffer and Sons, Cambridge, 1913)。

 "Natural Realism and Present Tendencies in Philosophy," \textit{Proc. Arist. Soc.}, 1908-1909, p. 165.

 \textit{Die Erfahrungsgrundlagen unseres Wissens}, p. 28.

 参照：\textit{Principia Mathematica}, Vol. I, * 14, and Introduction, Chap. III。存在の定義については、* 14. 02 を参照してください。

 参照：Edwin B. Holt, \textit{The Place of Illusory Experience in a Realistic World.} "The New Realism," p. 303, これはこの点と二重視覚についても言及しています。


\chapter{第IX章：原因の概念について --- {ON THE NOTION OF CAUSE}}

IX

因果の概念について


以下の論文では、まず「原因」という言葉が誤解を招く関連性と深く結びついており、哲学的な用語から完全に排除すべきであることを主張したいと思います。次に、科学において哲学が仮定する「因果法則」の代わりに使用されている原理（もし存在すれば）を探求したいと思います。さらに、因果に関する誤った概念と密接に関連していると私は考えられる、特に目的論と決定論における混乱を明らかにしたいと思います。

すべての哲学者は、どの学派でも、因果関係が科学の基本的な公理や仮定の一つであると考えています。しかし、奇妙なことに、重力天文学のような高度な科学では、「原因」という言葉が全く使われていません。ジェームズ・ウォード博士は『自然主義と無知論』の中で、この事実を物理学に対する批判の根拠として挙げています。彼は、世界について最終的な真実を明らかにしたいと考える人々にとって、原因の発見がその仕事であるべきだと考えているようです。しかし、物理学はそれらの原因を探求すらしていません。私には、哲学がこのような立法的機能を持つべきではない理由があり、物理学が原因を求めなくなったのは、実際にはそんなものは存在しないからだと考えます。因果法則は、哲学者たちの間で通用している多くのものと同じく、時代遅れの遺物であり、王政のように誤って害がないと信じられているために生き残っているだけだと私は信じています。

哲学が通常「原因」と理解することを知るためには、バーレンドの辞書を参照しました。私の期待を超える結果を得ることができました。そこで以下の三つの互換性のない定義を見つけました：

  "因果律. (1) タイムシリーズにおける事件間の必然的な関連..."

  "原因（概念）. 何かが別の過程によって起こる結果としてのプロセスを思い浮かべたり、認識したりする内容が含まれているもの..."

  "原因と結果. (1) 原因と結果...は区別できる2つの要素、状態、または現実の側面を示す相対的な用語であり、最初が存在しなくなるたびに第二が即座に存在し始め、第二が存在するたびに第一が即座に存在しなくなるような関係にあるものである。"

これら三つの定義を逐一考慮するものとする。第一の定義は明らかに「必要」という概念がなければ理解不能である。この項において、ボールドウィンの辞書には次のように記されている：—

「NECESSARY. これは真であり、かつどのような状況下でも真であるものを指す。したがって、「必然」の概念には単なる物理的な強制以上の要素が含まれており、その発生には一般的な法則がある」

因果関係の概念は必要性と密接に関連しているため、上記定義をlingerすることなく考察することは逸脱とは言えないとする。可能であれば、この定義から何か意味を見出すことを試みる価値がある。現状では、非常に明確な意味を持つものではない。

最初に注意すべき点は、「どのような状況下でも真である」という表現が意味を持つ場合、その主語は述語関数であり、単なる命題ではないということである。 命題は単純に真か偽のいずれかであり、それ以上のことはない：状況についての議論はあり得ない。「チャールズ1世の頭が切り落とされた」は夏でも冬でも同じように真であり、日曜日でも月曜日でも変わらない。したがって、「どのような状況下でも何かが真である」という表現が価値がある場合、その「何か」は述語関数である、すなわち変数を含む表現であり、変数に値が割り当てられると命題となるものでなければならない。言及された「状況」は、変数が持つ異なる値すべてである。したがって、「必要」という概念が「どのような状況下でも真である」という意味を持つ場合、「xが人間であれば、xは死ぬ」は必要であり、任意の可能なxの値に対して真であるからである。従って、我々は以下の定義に至る：—

「NECESSARYは述語関数に対する述語であり、その引数または引数すべてについて常に真であることを意味する」

不幸なことに、Baldwinの辞書における定義によれば、「必要である」とは「すべての場合に真」であり、「真」であることの両方が必要とされています。これら二つは不一致です。命題のみが「真」であり、 propositional function のみが「すべての場合に真」であることになります。したがって、その定義自体は意味をなしません。おそらく意図された定義は次のようであるべきです：「命題がすべての場合に真であるとき、それは propositional function の値であり、それがすべての場合に真であるための引数または引数全てがとえられる」。しかし、この定義を採用すると、同一の命題がその論理項の一つを選んで propositional function の引数とするか否かによって必要であるか偶然的であるかが変わります。例えば、「ソクラテスが人間であれば、ソクラテスは死人である」はソクラテスを引数と選んだ場合に必要であり、ソクラテスや「人間」を選んで propositional function の引数とする場合はそうではありません。また、「ソクラテスが人間であれば、プラトゥノは死人である」はソクラテスまたは「人間」を選んで引数とする場合に必要ですが、プラトゥノまたは「死人」を選んで引数とする場合はそうではありません。しかし、この困難を克服するためには、どの構成要素が引数と見なされるべきかを明確に指定することで、次のような定義に至ります：

「あるとえられた構成要素に関して命題が必要であるとは、その構成要素を任意の方法で変更しても命題が意味を持つ限り真であることを指す。」

この定義を引用された因果律の定義に適用します。明らかに、引数は以前のイベントが発生する時間です。したがって、因果関係の一例は次のようになります：「イベント [Math: e\textit{{1}] が時刻 [Math: t}{1}] に発生すると、その後にはイベント [Math: e\textit{{2}] が発生します」。この命題は [Math: t}{1}] を変更しても真であることを意図しており、すなわち [Math: t\textit{{1}] のどの方法でも変更されても真であることを指します。因果関係は普遍的な法則として次のようになります：「任意のイベント [Math: t}{1}] がとえられたとき、その後にはイベント [Math: e\textit{{2}] が発生するようなイベント [Math: e}{2}] がある」。しかし、これを見事な定義にするためには、[Math: e\_{2}] がどれだけ後ろに発生するかを指定する必要があります。したがって、原則は次のようになります：

「任意のイベント [Math: e\textit{{1}] に対して、あるイベント [Math: e}{2}] と時間間隔 [tau] が存在し、[Math: e\textit{{1}] が発生すると、その後 [tau] の間隔で [Math: e}{2}] が発生する」。

私はまだこの法則が真偽について考慮する時期ではない。現在は、因果律そのものとは何かを明らかにすることに焦点を当てるべきである。したがって、上の定義の中で他の定義に移る。

第二の定義については長く議論する必要はない。まず第一に、これは心理学的なものであり、因果関係を考える際に私たちに関心のあるのは、「過程の考えや感覚」ではなく、その過程自体であるからである。第二に、これは円環論理を導入しており、別の過程「によって結果として」起こることと述べることで、定義すべき因果の概念をすでに含んでいるからである。

第三の定義は最も正確であり、明瞭さにおいて欠けるところがありません。しかし、この定義によって主張される原因と結果の時間的連続性には大きな困難があります。二つの瞬間は連続していませんから、時間係列は密接であるため、原因または結果あるいは両方が、定義が正しい場合、有限な時間を過ごす必要があります。実際、定義の文法によれば、双方とも有限な時間内に存在すると推測されます。

しかし、この問題は二つの難点を生じさせます：まず、原因が自己内部での変化を含む過程である場合、因果関係はその初期部分と後期部分との間で必要となります（因果律が普遍的であるなら）。さらに、初期部分は結果とは連続しておらず、したがって定義によれば結果に影響をとえることはありません。従って、原因の初期部分を変更しても結果には影響がないため、真の原因として定義されたものは達成されません。なぜなら、定義は複数の原因を排除しているからです。

一方で、原因が自己内部での変化を伴わない純粋な静態である場合、まず第一に自然の中にそのような原因を見つけることはできませんし、第二に、原因が平穏に存在した後突然結果として爆発するということが不思議であり、論理的に可能であっても受け入れられません。なぜなら、原因はいつでも同様の状態で存在していたり、結果を生じさせずに変化することもできていたからです。

したがって、因果関係が時間上連続しているという見解は致命的な矛盾に直面します：原因と結果が時間的間隔 tau で分かれている必要があります。これは先述の第一定義の解釈において仮定されました。

因果律に関する同様の声明は、バロードの最初の定義から引き出されたものと同じ本質的な内容を他の哲学者も述べています。例えばジョン・スチュアート・ミルは次のように言います：


「因果法則、その認識は推論科学の主な柱であり、それは観察によって得られる恒常的な順序の不変性が自然の事象とその前に行われた他の事象との間で見つかるという、身近な真実です。」

そして、哲学者たちによって述べられた法則は価値がないと正しく認識したベルソンですが、それでもこの法則が科学で使用されていると仮定し続けます。彼は次のように言います：

「つまり、この法則 [因果法則] は、現象はその条件によって決定されるという意味であり、言い換えれば同じ原因が同じ結果を生むということです。」

また、以下にも述べています：

「物理現象を観察し、これらの現象は法則に従います。これは以下の意味を持ちます：（1）以前観察された現象 a, b, c, d は同じ形で再び発生する可能性があります；（2）条件 a, b, c, d の後に現れた特定の現象 P が、これらの条件のみによって現れ、同じ条件が再度存在した時点で欠かさず再発生します。」

ベルソンによる科学への大部分の攻撃は、この原理が使用されていると仮定するところに基礎を置いています。実際には、そのような原理を使用していませんが、哲学者たちは—ベルソンも含め—しばしば相互に科学に関する見解を得てしまい、直接科学から得ません。この原理自体について、異なる学派の哲学者間では相当な合意があります。しかし、すぐに問題となるいくつかの困難が生じます。原因と結果の間に複数の原因については現在考慮しませんが、他のより重大的な問題を考慮する必要があります。上記の法則の述べ方から二つの問題は以下の通りです：

（1）「事象」とは何を指すのか？

（2）原因と結果の間にはどの程度の時間差があることができるのか？

(1) 法において「事件」とは、再発する可能性が高いものであることが明らかで、そうでなければ法は単純化されてしまう。従って、「事件」は特定のものではなく、多くの例を有する普遍的な概念である。また、「事件」は宇宙全体を超えるものではないことも推論できる。再発する確率が非常に低いからである。「事件」とは何を指すかと言えば、火をつけることや自動機にコインを入れることのようなものだ。このような事件が再発するためには、その定義は狭くならず、火をどのように打つか、コインの温度はどの程度かなど具体的な条件を述べる必要はない。もしこれらの考慮事項が重要であるとすれば、「事件」は最多一度しか起こらないことになり、法は情報提供の役割を失ってしまう。従って、「事件」は多くの時間的な個別の発生例を含む広範な普遍的概念で定義されるべきである。

(2) 次に時間を考慮するが、哲学的には原因と結果は時係列上連続していると考えられるかもしれない。しかし、既に述べた理由によりこれは不可能であり、したがって無限小の時間間隔がないため、原因と結果には有限な時間間隔 [τ] があるべきである。しかしこれは即座に不可克服の困難を引き起こす。どの程度短い時間間隔 [τ] を設定しても、その間に予想される結果を防ぐような何かが起こる可能性がある。コインを入れた後、チケットを取り出す前に地震があり、機械が傾き、計算が狂う。予想される効果を得るために確実であるためには、環境に何が影響しないかを知らなければならない。しかし、これにより仮定された原因は単独で効果を保証するのに十分ではないことが意味される。そして環境を含めると、再発の確率は減少し、最終的には全体の環境を含めた場合、再発の確率はほとんど nil となるまで減ってしまう。

これらの困難にもかかわらず、日常生活中で比較的頼りになる一連の規則性は確かに存在する。これらは因果律という仮定的な法則を提示した。それらが機能しない場合、より良い表現が見つかったと信じられ、それが常に機能しなければならないと考えられる。私はこのような序列が実際に一度も失敗しないものがあることを否定していない。ある一定の質量以上で動いている石が、一定の厚さ未満のガラスに接触すると必ず割れるという法則は存在するかもしれない。また、これらの規則性を観察することは、科学の初期段階において有用であると認めている。空中に浮かぶ物体が通常落下することという観察は重力法則への一歩となった。しかし私は、このような不変の一連の規則性が存在し、科学がそれを探求することを目指しているとは否定する。

これらの規則性は、「事件」の定義における一定の曖昧さに依存していることを忘れてはならない。物体が落下することは量的で質的な漠然とした主張であり、科学はその落ちる速度を知りたいと考えている。これは物体の形状と空気の密度によります。実際には、真空中で落下する場合に近い一貫性があることは認められる。ガリレオが観察した限りでは、その一貫性は完全だった。しかし後になって、すらすらとは異なる角度や高度にも影響があることが明らかになった。理論的には、太陽と月の位置も影響をとえるだろう。つまり、科学の進歩とともに、最初に観察された粗雑な一連の規則性から遠ざかり、前因後果の分化が深まり、関連する前因の範囲が広がっていく。

したがって、「同じ原因、同じ結果」という原理は哲学者が科学にとって重要だと考えるべきではない。予測可能な結果を得るためには十分に詳細な前因が必要であり、その複雑さは再現性を希薄にする。したがって、この原理に基づくとすれば、科学は全く成果を出せないだろう。

これらの考慮の重要性は、科学的手法についてより正確な説明を導くことに一部あり、また人間の意志と類似性が原因の概念に誤りを生む源となる類推を排除することにも一部ある。後者の点については、いくつかの例によって理解しやすくなる。

（1）「因果は多少なりとも類似しているべきである。」この原理は偶然主義の哲学において重要な位置を占めており、現在でも完全には消えていない。例えば、心が以前に何も意識的なものがない宇宙から発展したとは考えにくいという信念があるが、その理由の一つは物質が心とあまり似ていないため、それを引き起こすことができないとされるからである。または、特に、我々の性質の中でも「優れた部分」は、常に何か少なくとも同程度に優れたものが存在してそれが原因となった場合にのみ説明できると考えられている。これらの見解は、因果関係の過度に単純化された法則を前提としているように思われる。科学的な意味での「原因」と「結果」について正しく定義すれば、通常は非常に異なるものであることが示されるからである。実際には、「原因」は全体の宇宙の二つの状態であり、「結果」は特定の出来事である。

（2）「因果は意志と類似しているべきである。なぜなら、原因と結果間に理解できる関係性がなければならないからである。」この原理は、私は哲学者が明示的にそれを否定するときでも無意識に持っていると考えられることが多い。我々が以前に考慮したように、心が純粋な物質的世界から生じたとは考えにくいという見解においてその作用があると思われる。しかし、「理解できる」という言葉の意味は私には不明である。それは「想像に馴染みやすい」という意味だと考える人が多いようだ。意志の行為とその実現との間の関係性は、他のどの意味でも最も「理解できない」ものである。しかし明らかに、原因と結果間に求められるような種類の関係性は、因果律が仮定する「事件」間だけが持つことができるだろう。物理学のような科学において因果律を置き換える法則は、任意の二つのイベント間でどのような関係性を探しても場所を提供しないからである。

(3) 「原因は何かしらの意味で結果を強制するが、その逆は成り立たない。」この信念は決定論への嫌悪感に大きく影響をとえているように思われるが、実際には二番目の原理と結びついており、その原理を放棄すれば崩れてしまう。以下のように「強制」と定義できる。「ある状況がAが何かをしたいと思うのを防いだり、何かをしないために促すとき、その状況はAを強制するという」この定義には、「原因」という言葉の意味を見つけることが前提となっている。現在私が明確にするのは、強制は非常に複雑な概念であり、望みが叶わないことを含む点である。人が自分の望むとおり行動している限り、強制は存在しないし、以前の出来事を使ってどれだけ彼の願いを予測しても強制は存在しない。そして望みが関としていない場面では、強制について議論する余地がない。したがって、一般的には原因が結果を強制すると考えるのは誤解に過ぎない。

同じ原理の曖昧な表現では、「決定」という言葉に「強制」が置き換わる。「原因は何かしらの意味で結果を決定するが、その逆は成り立たない」と言われる。この「決定」という用語の意味は明確ではない。私が知る限り、唯一の正確な意味は関数や一対多の関係である。複数の原因を認めつつ効果は認めない場合、つまりとえられた原因が特定の結果を生むことを示す一方で、とえられた結果が多くの可能性の中から一つを選んだことを示す場合、私たちは原因が結果を決定し、その逆は成り立たないと述べることができる。しかし複数の原因を認めることは、効果を曖昧かつ狭い視点でとらえ、原因を正確かつ広い視点でとらえることから生じる。多くの前因が「原因」に該当する場合、それはその人の死が曖昧かつ狭い視点で捉えられているためである。しかし、「原因」として毒薬の摂取、「効果」として五分後の世界全体の状態をとらえると、結果の多様性ではなく原因の多様性が生じる。したがって、「原因」や「効果」の間の想定される対称性の欠如は虚構である。

(4) 「原因は存在しなくなった後には作用しない。なぜなら、存在しなくなったものは何もないからである。」これは一般的な格言であり、広くは明示されない偏見でもある。私はこの格言が過去の影響がある現在に及ぼしていることから、ベルソンの「durée」が魅力的である理由の一部と捉えている。過去が現在に影響を及ぼす限り、過去は何かしらの形で存在していると考える。この格言の誤りは、「原因が作用する」という仮定にある。「意志が作用する」とは、意志によって起こることが実際に起こったときにその意志が作用すると考えられるが、何もが作用するには意志しかない。意志が「作用する」という信念は、意識的にも無意識的にも意志に類似させることから生じる。すでに述べたように、原因があるとしても、その効果と有限な時間間隔を保たなければならないので、結果として存在しなくなった後にその効果を引き起こす。

上記の意志が「作用する」という定義に対して反論があり得る。「意志はただそれが望むものを引き起こしたときにのみ作用する。それだけが通常の意味での意志が「作用する」ことを示すものである」という主張だが、これは我々が攻撃している因果律の見解を含んでいるため、定義としては適さない。意志が作用するとは、類似の状況で類似の意志が通常は望むものを引き起こす法則があるときにその意志が作用すると述べることができる。しかし、これは曖昧な概念であり、まだ考慮していないアイデアを導入する。最も重要なのは、「作用する」という一般的な概念が原因を否定した場合に開かれる定義としては利用できないということである。

(5) 「原因は存在しない場所で作用しない。」これは広く普及している格言であり、ニュートンに対して抗議され、遠隔作用に対する偏見の源となっている。哲学では、この格言は瞬間的な作用を否定し、それから単一論やライプニッツの単子論へと導き出される。時間的連続性に関する類似の格言と同じく、この格言も「原因が作用する」という仮定に基づいている、すなわち、それが意志に不透明な形で類似しているという仮定である。また、時間的連続性の場合と同様に、この格言から導かれる推論は全く根拠がない。

私は今や、どのような法則や法則群が因果律の仮定を置き換えることができるかについて問うことに戻る。

まず、伝統的な法則が考慮するような順列の範囲を超えていない限り、多くの場合で観察され、かつ未だに失敗したことがない順列があれば、その順列が将来的の事例でも成り立つ可能性が高いという帰納的確率を認めることができます。石が窓を破壊するという順列であれば、今後も同じように破壊すると予想できます。もちろん、この帰納的原理は正当に疑問視されるべきものであるかもしれませんが、現在の議論ではその真実性を不変として扱います。つまり、どの頻繁に観察される順列についても、前件が原因で後件が結果と述べることができます。

しかし、このような特別な順列は伝統的な因果関係とは非常に異なる点があります。まず第一に、これまで未だに観察されていない事例においての順列はただの可能性であり、一方で因果関係は必要性であると想定されていました。これは単に我々が真の因果関係を発見したことを確信していないという意味ではなく、現在の私たちの感覚における因果関係の場合でも、観察に基づいてある事象が起こると別の事象も起こる可能性が高いというだけです。したがって、夜が昼の原因であると述べる場合でも、実際にはBがAを伴わないケースがあるとしても、打火石を打つことがその燃え上がる原因となるでしょう。打火石が湿っているために点火しない場合もありますが、それでも打火石を打つことは燃え上がる原因となります。

第二に、この意味での因果関係はすべての事象に対して仮定されるわけではありません。私たちはそれらの順列を見つける限りのみ、因果関係を信じるだけであり、常にそのような順列が存在することを前提にはしません。

第三に、十分に頻繁な順列であれば現在の私たちの感覚における因果関係となります。例えば、夜が昼の原因であると述べることも拒否されません。我々がこの表現を避ける理由は、その順列が失敗する可能性があるという想像が容易だからであり、因果関係は時間的な間隔を持つからです。そのような順列は他の状況によって中断される可能性があるため、どんな順列でもそのような中断が起こることで失敗する可能性があります。ミルはこの夜と昼の例について述べています。

彼等は「原因」という言葉を使用するためには、前件が常に後件を伴ってきただけでなく、現状の事物の構造が存続する限り、その関係が常に続くという信念を持つべきである必要がある。

この意味では、ミルが考えていたような因果法則を見つけることを諦める羽目になる。我々が観察したどのような因果的な連鎖も、いつでも新たな法律によって否定される可能性があり、それらの法則は進歩した科学が追求する種類のものではない。

第四に、確率的な順序法則は日常や科学の初期段階では有用だが、科学が成功すると別の法則によって置き換える傾向がある。重力法則は、どの科学でも同様な状況を示す例である。互いに重力をもつ物体の運動においては、原因と結果という言葉を使う余地がない；ただ単に公式があるだけである。特定の微分方程式が存在し、これらはシステム内の任意の粒子について各瞬間に常に成り立つものであり、とえられた構造と速度、または二つの時間点における構造を知っている限り、他の任意の早期または後の時間点における構造を理論的に計算することが可能となる。つまり、任意の瞬間の構造はその瞬間と二つの指定された瞬間における構造の関数である。この主張は物理学全体に適用され、重力だけの特別な場合に限られるものではない。しかし、このようなシステムにおいては「原因」という言葉や「結果」という言葉を適切に使うことはできない。

疑問なき事、古い「因果法則」が長く哲学の書籍に支配的影響をとえ続けた理由は、ほとんどの哲学者が関数の概念に馴染みがないためであり、そのため彼らは過度に簡略化した表現を求めているからである。科学の法則における恒常性とは、「同じ」原因が「同じ」結果を繰り返すというような単純な一致ではなく、「関係」の一致性が求められる。すなわち、「関係」の一致性もまた単純すぎる表現であり、正確な表現は「微分方程式の一致性」である。この事実を非数学的な言葉で正確に述べることは不可能であり、最も近い表現としては以下のようなものがある：「宇宙の任意の瞬間における状態と、その瞬間に任意の部分が変化している速度の変化率との間には恒常的な関係があり、この関係は一対多であり、つまり宇宙の状態がとえられたとき、その速度の変化率は確定的である」というものだ。科学の実践において何かを実際に発見できる「因果法則」であれば、上記の命題は哲学の書籍に見られるあらゆる「因果法則」よりも名にふさわしい。

以上の原理について、いくつかの観察を行うべきである—

(1) 上記の原理は、先験的または自明なものではなく、「思考の必然性」でもないことは明白である。また、科学の前提でもなく、そのような意味での前提ではない：これは、自己確認的な法則が多数から導き出される経験的一般化である。

(2) 法則は過去と未来を区別しない：未来は過去を「決定」し、過去も未来を「決定」するという意味において同じように作用する。ここでの「決定」という言葉は純粋に論理的な意味を持つ：ある数の変数が別の変数を「決定」するとは、その他の変数がそれらの変数の関数である場合を指す。

(3) 法則は、ある十分に小さな体の内での事象の進行が、宇宙内のどの二つの状態においても、その小さな体からかなり離れた部分に関してのみ異なる場合でも、ほぼ同じであるという条件下で、経験的に検証可能ではありません。例えば、太陽係における惑星の運動は、固定星がどのように分布されているかに関係なく、すべての固定星が太陽よりも惑星に比べて非常に遠方に存在する限り、ほぼ同じであることが予想されます。万有引力が距離に直接比例した場合、最も遠く離れた星が惑星の運動に最大の影響をとえるなら、世界は現在と同様に規則的で数学的な法則に従うかもしれませんが、その事実を発見することは決して不可能です。

(4) 古来的な「因果律」は科学では仮定されていませんが、「自然の均一性」と呼ばれる何かが仮定され、あるいは帰納的な根拠に基づいて受け入れられています。自然の均一性は、「同じ原因、同じ結果」という平凡な原理を主張しませんが、法則の永久性という原理を主張します。つまり、例えば構成関数として加速を示す法則が観測可能な過去全体で成り立つことが確認された場合、それは将来も同様に成り立つと期待されます。または、それが自身で成り立たない場合でも、過去に関して仮定された法則と一致する別の法則が将来に適用されることを期待します。この原理の根拠は単に、非常に多くの事例でその真実性が確認されているという帰納的な根拠であり、従って原理自体は確実なものではなく、その確率は正確には評価できない程度にしか保証されません。

自然界の均一性は、科学の実践において仮定されるものの、その一般性からして、すべての科学的推論が誤りとなるような大前提として扱われることは適切ではない。自然法則全てが永久に存続すると仮定することは、特定の法則が永久に存続すると仮定するよりも確率が低く、また特定の法則が特定の日付まで有効だと仮定することもその日付以降の有効性を保証するよりも確率が低い。具体的な科学的状況では必要なことを仮定し、それ以上のことは仮定しない。1915年の航海暦を作成する際には、重力法則がその年終了まで真実であると仮定するが、1916年については次の航海暦の巻が出るまで仮定しない。この手続きは、自然の均一性が先験的に知られているものではなく、経験的な一般化であるという事実に基づく。すなわち、「すべての人間は死ぬ」というようなものと同じようにである。このような場合、とえられた具体的な例から新しい例へ直接論じる方が良い。結論はどちらの場合も確率的だが、前者の方法ではより高い確率を持つ。

すべての科学において二種類の法則を区別しなければならない：一つは経験的に検証可能でおそらく近似的なもの；もう一つは検証不可能で正確である可能性があるもの。重力法則は、太陽係への適用に関しては、太陽係外の物質がそのような目的では無視できると仮定した場合に経験的に検証可能である。この仮定をもとに、我々はこれが大体真実であると考えるが、普遍的な重力法則が正確であると信じているものについては経験的に検証することはできない。この点は、「相対的に隔離されたシステム」と呼ぶべきものに関連して非常に重要である。

以下のように定義できる：

特定の期間中に相対的に隔離されたシステムとは、その期間内で誤差を一定範囲内に保って同じ行動を示すと予想されるシステムであり、宇宙の残り部分がどのような構成であれ同様の行動を示す。

一係が「実質的に隔離されている」と呼ばれる期間がある場合、その期間中は、他の宇宙の状態が指定された誤差よりも大きい影響を及ぼす可能性こそ存在するとしても、そうした状態が実際に起こらない理由があるからである。

厳密には、我々はシステムが相対的に隔離されている点を明確に指定すべきである。例えば、地球は落下物体に関しては相対的に隔離されているが、潮汐に関してはそうではない；経済現象に関しては実質的に隔離されているものの、ジェボンズの太陽黒子説が正しい場合、商業危機に関してはすら実質的に隔離されていなかった。

我々は注意しなければならないのは、事前に一係が隔離されていることを証明することはできないということである。これは、このシステムだけに約束された均一性を述べられる観察結果から推論されるからである。もし宇宙全体の完全な法則が分かっていたら、その法則から一係の隔離状態を導くことができたであろう；例えば万有引力の法則がある場合、太陽係のこの点での実質的な隔離は、その近傍にほとんど物質がないという事実によって推測される。

しかし、隔離されたシステムが科学法則の発見の可能性を提供する重要な役割を持っているだけであり、科学の完成した構造において理論的には重要ではないということに注意しなければならない。

「事件Aが「原因」として別の事件Bを引き起こす」という哲学的な基本的なケースは、実質的に隔離されたシステムの最も単純化された例である。ある期間中、特定の条件下でAが発生すると常にBが続くという一般科学法則がある場合、その場合、AとBはその期間を通じて実質的に隔離された一係となる。

しかし、これが起こることは幸運であり、それは特別な状況によるものであり、他の宇宙が異なる状態であっても同じ法則を適用していたとしても真でなかったであろう。

因果関係が果たした重要な役割は、過去から将来を推論する可能性、またはより一般的には、特定の指定された時間からの任意の時間のイベントを推論することである。このような推論が可能な一係は「決定論的な」一係と呼ばれることができる。以下のように決定論的な一係を定義できる：

任意のそのような推論が可能な一係を「決定論的な」一係と呼ぶことができる。

システムが「決定論的」であるとは、とえられたデータ
[$e_1, e_2, \dots, e_n$、時刻 $t_1, t_2, \dots, t_n$] について、このシステムに関するものとして、任意の時刻 $t$ におけるシステムの状態を $E_t$ とすると、関数的な関係

\[ E_t = f(e_1, t_1, e_2, t_2, \dots, e_n, t_n, t) \quad \text{(A)} \]

が存在することを指す。この式において $t$ はとえられた期間内の任意の時刻であるとしても、その期間外ではこの式は成り立たない可能性がある。もし宇宙全体がこのようなシステムであれば、決定論は宇宙に真実であり、そうでなければそうではない。このようなシステムの一部であるシステムを「決定された」と呼び、それ以外のものは「恣意的な」と呼ぶ。

イベント [$e_1, e_2, \dots, e_n$] をこのシステムの「決定因子」と呼ぶ。注意すべきは、一つのセットの決定因子を持つシステムが一般に多くの決定因子を持つことである。たとえば、惑星の運動の場合、任意の二つのとえられた時間における太陽係の構成はそのような決定因子となる。

また、心理物理学的並行性の仮説から別の例を挙げることができる。目的のために、ある脳状態に対応する一定の心状態が常に存在し、逆もまた然りであると仮定しよう。つまり、それらは一対一の関係にあるのであり、一方は他方の関数であり、その逆もまた然りである。さらに、実際的に確実なことは、ある特定の脳が二度と完全に同じ状態になることは稀であるため、とえられた宇宙全体の一定の状態に対応する一定の脳状態があることである。したがって、特定の人間の心状態と全宇宙（物質的および心理的な）状態との間に一対一の関係がある。これにより、もし心理物理学的並行性が真実であれば、とえられた n つの宇宙全体の状態が宇宙全体の決定因子であるなら、特定の人間の n つの心状態は全宇宙（物質的および心理的な）の決定因子となることが推論できる。

上の図は、心と物質の関係について思索した人々に見舞われた定められた混乱と密接に関連しています。しばしば考えられるのは、脳の状態がとえられたときの心の状態が確定的であるとともに、物質的世界が決定論的なシステムである場合、心は何かしらの意味で物質に「支配」されていると考えられ、逆に物質は心に支配されていないとされています。しかし、心の状態がとえられたときに脳の状態も確定的である場合、物質を心に「支配」されるものとして考慮するのと同じくらい、心を物質に「支配」されるものとして考慮することも正確です。理論的には、物質を一度も言及せずに心の歴史を計算し、最後には物質が対応する歴史を通じて進んだことを導くことができます。実際、脳と心との関係が1対多でない限り、一方的に心に依存することになります。逆に、ベルソンが仮定するように、関係が多対1である場合、物質は心に一方的に依存します。しかし、この依存性は常に論理的なものであり、私たちを望まないことを強制されるという人々の本能的な想像とは異なります。

また、機械論と目的論の例を挙げてみよう。システムは、特定の物質の位置が特定の時間にとえられた場合など、純粋に材料的な決定因子を持つとき、「機械的」と定義されることがある。私たちは知る世界における心と物質の世界が機械的システムであるかどうかは開かれた問題であり、仮説として、それが機械的システムであるとする。この仮説——私は主張する——無論、宇宙が或いは或いでは「目的論的」なシステムであるかどうかという問いには何らの光を投げかけない。正確に言えば、「目的論的」なシステムとは何かを定義するのは難しいが、我々が採用する特定の定義によって議論は大きく影響されない。大まかに言うと、目的論的システムは目的が達成されるシステムであり、つまり、より深いものやより高貴なものやより根本的なものやより普遍的なものなどと呼ばれる特定の欲望が実現されることを伴うシステムである。現在の事実——それが事実であるとする——が宇宙が機械的であるという事実は、上記の意味で宇宙が目的論的かどうかという問いには何らの影響もとえない。すべての願いが達成される機械的なシステムやすべての願いが妨げられる機械的なシステムも存在する可能性がある。したがって、実際の世界が或いはどの程度まで目的論的であるかという問いを決定するのは、それが機械的であることを証明することではなく、その世界が機械的でないことを望む理由は存在しない。

これら全ての問題において、推論できるものと実際に決定されたものが混同される非常に大きな困難があります。一時的に未来が「決定されている」という意味を考えてみましょう。その意味には科学的法則によらず独立して決定されるという重要な意味があります。つまり、未来はただ起こるだけではなく、それがどのようなものになるかは既に決まっていると解釈できるのです。

過去はすでに起こった事実によって決定されているとみなしますが、記憶が前向きではなく後ろ向きであるため、未来も同様にその事実によって決定されているべきだと考えます。「しかし」と言われ、「あなたは過去を変えることはできませんが、未来は一定程度変えることができます」という見解は、私があらゆる誤りを取り除こうとしたものと同じような因果関係に関する誤りに基づいているように思います。

既知の過去を他の何かに変えることは不可能です—これは矛盾法則の単なる応用です。既知の過去があると分かっているのであれば、それを異なるものに望むのは無意味です。同様に、未来も既知である場合（例えば、予想される日食の場合）を異なるものに望むのも無意味です。

しかし、「しかし」という反論があり、「私たちの願いはしばしば存在しない場合と比べて未来を異なるものにすることができますし、過去にはそのような影響はありません」と言われます。これはまた単なる同義語です。効果が原因から続くものであると定義すると、明らかに過去に影響をとえることはできません。しかし、現在の願いが異なる場合でも、過去も異なる可能性があるということを意味するわけではありません。

現在の私たちの願いは過去によって制約されているため、過去が異なる場合でなければ現在の願いも異なることができません。したがって、現在の願いが異なる場合、過去も異なるはずです。もちろん、過去は既に決まっているので、現在の願いも同じように既定のままです—これはまた矛盾法則の応用です。

事実は（1）願望は一般的に知識不足によるものであり、したがって未来に関する願望の方が過去に関する願望よりも頻繁である；（2）未来に関連する願望とその現実化は「実質的に独立したシステム」を形成することが多い、つまり多くの未来に関する願望が現実化される—のような点に過ぎません。しかし、過去は記憶によって知ることができるが未来はそうではないという偶然の事実からくる主な違いがあるように思われます。

Although "決定"という意味において、未来はその単なる事実であると予定されているからこそ決定的だとされる場合でも、それでも一部の決定論反対派、特にM・ベルジャンソーや pragmatists を説得するのに十分なのは明白です。しかし、人々が未来を決定的と言及する際には、過去に機能として示され、少なくとも理論的には計算できる形で未来を展開することを指していることが多いとは限りません。

そのような公式がいかなる複雑さの如何に関わらず認められる場合でも、特定の測定可能な量によってとえられた一定の瞬間における状態が関数であるシステムはすべて決定論的システムとみなされるべきかのように思われます。これを例に考えてみましょう。単一の物質粒子について、時間 $t$ での座標は [Math: x\textit{{t}, y}{t}, z\textit{{t}] です。粒子がいかなる方法で動くにせよ、理論的には関数 [Math: f}{1}, f\textit{{2}, f}{3}] が存在し、

[Math: x\textit{{t} = f}{1}(t), y\textit{{t} = f}{2}(t), z\textit{{t} = f}{3}(t).]

と表すことができます。これにより、理論的には時間 $t$ における物質宇宙の全体的な状態は $t$ の関数として展開可能であることが推測されます。したがって、我々の宇宙は上記で定義された意味での決定論的となります。しかし、これが真実であれば、その宇宙について決定論的であることを述べるだけでは情報が伝えられません。確かに、関わる公式は厳密に無限の複雑さを持つ可能性があり、実際に書き下したり理解することは不可能かもしれません。しかし、知識から見れば、それ自体において、上記の考慮が妥当であれば、物質宇宙は決定論的でなければならず、法則に従うべきであると推測されるべきです。

しかし、これには明確に意図されたものとは異なっている。この見解と意図された見解の違いは次のように見ることができる。ある公式がこれまでの事実に適合している場合、例えば重力法則の場合、その他の無数の公式がある。これらの公式は過去において実証的に区別できないが、将来にはそれらからますます離れる。したがって、恒久的な法則があると仮定しても、逆平方律が将来にまで適用されることを理由として推測することはできない；それはこれまで区別できなかった他の法則である可能性がある。過去の事実が一つの法則を遵守する場合でも、それらは他の法則も遵守しており、これらは過去には区別できなかったが将来には離れるものである。したがって、常に存在し、これまでに破られなかった法則が初めて破られている瞬間がある。実際には科学は、事実に適合する最も単純な公式を選択している。しかし、これは明らかに方法論的な戒律であり、自然の法則ではない。最も単純な公式が時間とともに適用できなくなると、まだ適用できる最も単純な公式が選択され、科学には公理が否認されたという感覚がない。したがって、多くの科学分野においてこれまでに非常に単純な法則が見つかったという事実は、あらかじめとえられた理由ではなく、かつ過去の法則が継続すると推論するための証拠としても使用できない。これはなぜなら、常に以前に真であった法則が否認されているからである；ただし、進んだ科学ではそれよりも単純な法則が残っている。さらに、先進的な科学の状態から他の科学の将来の状態を推論することは誤りであり、それはおそらくそのような単純で容易に確認できる法則を遵守してきた分野は進んだ科学であり、他の分野ではそうではないためである。


我々が考慮している困難は、主として時間が公式に入れられることのない原理によって解決されるか、あるいは完全に解決されるかもしれません。すべての機械的法則は、加速が構造と時間の両方ではなく、単に構造の関数であることを示しています；そしてこの時間の無関心性の原理は、科学的な法則にも拡張できます。実際、自然の均一性を「科学法則が時間という引数を含まない」という意味と解釈することもできるでしょう（ただし、そのような形でとえられた場合を除き）。この考慮事項が我々の困難を完全に克服するかどうかは分かりませんが、それでもそれは困難を大幅に軽減します。

自由意志に関する問題への適用例として、これは述べたことを説明するのに役立ちます。

（1）意志に対する決定論は、我々の意志がいくつかの決定論的システムの一部であるという教義です。つまり、上記で定義した「決定」の意味において意志が決まっているということです。この教義が真偽如何は事実上の問題であり、以前の議論が正しいと仮定すれば、どちら側にも先験的な考慮事項はありません。一方では、因果関係という先驗的なカテゴリーは存在せず、観察された均一性のみがあります。実際、意志に関する観察された均一性があるため、意志が決まっているという経験的証拠があります。しかし、その証拠が圧倒的であると主張することは非常に冒険的であり、いくつかの意志も他のこと同様に完全には決まらない可能性もあります。

（2）しかし、他方では、決定論に対する自由意志の主観的な感覚は問題に関係ありません。そのような見解が関係があるという視点は、原因が効果を強制するか、または自然が法を人間の命令のように遵守させると信じる信念に基づいています。これらは単なる人間中心主義的迷信であり、原因と意志の類似性、そして自然法則と人間の命令との類似性によるものです。我々の内省的な意識にある自由とは異なる意志が強制されないという感覚は、その意志を我々自身が選んだものであることを意味します。伝統的な因果関係理論の欠点の一つは、決定論と我々が内省的に認識する自由との人工的な対立を作り出したことです。

(3) 一般に、意志が決定されるかどうかという問題以外に、その意志が機械的に決定されるかどうかという問題がある。すなわち、それらが先に定義された機械係の一部であるかどうかの問題だ。これは、ある特定の材料データをとえると、すべての意志がこれらのデータの関数となるような法則があるかどうかの問題だ。またここでも、実証的な証拠は存在するが、それが全ての意志について結論的ではない。しかし重要なのは、意志が機械係の一部である場合でも、これによって物質が心志を凌駕することを意味しないことだ。同じシステムが物質的決定要因を受け入れる可能性があるとしても、それはまた精神的な決定要因も受け入れられる可能性がある。したがって、機械係は意志のセットと同様に、事実のセットによって決定される可能性もある。従って、意志が機械的に決定されるという見解を嫌う理由は誤りであるように思われる。

(4) 決定性としばしば関連付けられる「必然性」の概念は、決定論から正当に導き出されるべき混乱した概念だ。必要性について話すときによく混同される三つの意味がある：—

([α]) 行動が必然的であるとは、その行動を他の方法で行いたいとしても、それが行われることをいう。決定論はこの意味での必然性を含意しない。

([β]) 命題関数が必然的であるとは、すべての値が真であるときをいう。これは現在の議論には関連がない。

([γ]) とえられた要素に関して命題が必然的であるとは、その要素を引数とし、必要な命題関数の値として真であるときをいう。すなわち、その要素を変えても真であるときをいう。決定論的なシステムにおいては、意志とその決定要因との間の接続が必然的である場合がある。これは、決定因子が発生する時間という要素を変化させたとしても一定の時間間隔を保つ限りだ。しかしこの意味での必然性は純粋に論理的なものであり、感情的な重要性はない。

我達現在まとめ一下因果律に関する議論。哲学的な表現では、因果律は誤りであり、科学では使用されていないことが分かった。次に、科学的法則の性質を考慮し、一つの事象Aが常に別の事象Bを伴うというよりも、特定の時間における一定の事象（我々はこれを決定因子と呼んだ）と他の事象（過去、現在、または将来の時間）との間の関数的な関係を述べることが科学的法則であることがわかった。先験的なカテゴリーが関わっていないことも確認した：科学的法則の存在は純粋に経験的事実であり、必ずしも普遍的ではなく、 trivial で科学的に無用な形を除いては。一つの決定因子のセットを持つシステムは、別の全く異なる種類の決定因子を持つ可能性があることも見つかった。例えば、機械的に決定されたシステムは、目的論的にまたは意志的に決定される可能性もある。最後に、自由意志の問題を考えた：ここでは、意志が決定されていると仮定する理由は強いが結論的ではないことがわかった。意志が機械的に決定されているとしても、反省によって示される自由を否定する理由にはならないし、機械的事象が意志によって決定されていないということも示されない。自由意志対決定論の問題は、我々が正しい場合、主に虚構的なものだが、一部はまだ決定的に解決できない状態にある。

FOOTNOTES:

 命題関数は変数や未確定な要素を含む表現で、変数に具体的な値が割り当てられると命題となる。例としては、「AはAである」、「$x$は数である」がある。変数は関数の引数と呼ばれる。

 \textit{論理学}, 3巻, 第5章, §2。

 \textit{時間と自由意志}, p.199。

 \textit{時間と自由意志}, p.202。

 同上，§6


\chapter{第X章：直接体験による知識と記述による知識 --- {KNOWLEDGE BY ACQUAINTANCE AND KNOWLEDGE BY DESCRIPTION}}

知識の認識と説明による知識


次の論文の目的は、「この～～」についての命題を知りながら、その「この～～」が誰か何であるかを知らない場合における我々の知識について考慮することです。例えば、私は最多得票者となる候補者が選出されることを知っていますが、最多得票者となる候補者が誰であるかは分かりません。私が考慮したい問題は、このような状況での主語が単なる説明によってとえられた場合における我々の知識は何であるかという点です。私は以前において、純粋に論理的な観点からこの問題を検討しましたが、ここでは知識理論との関連も考慮し、論理的な議論と並行してこの問題を検討したいと思います。上記の論理的議論により、本稿では論理的部分はできるだけ簡潔にします。

為了明瞭地闡明「認識」と「描写」の対比を理解するために、まず「認識」という言葉の意味について説明します。私はある対象に対して認識しているとは、その対象に対する直接的な認知関係を持っていることを指し、つまりその対象自体に直接的に気づいている状態を指します。ここで私が言う認知関係は、判断を構成するものではなく、描写を構成するものです。実際、私は認識という関係は、私をSと呼び、対象をOと呼ぶときの主語と客体との関係の逆であると考えています。つまり、SがOに対して認識しているとは、OがSに描写されることと同じことです。しかし、「認識」という言葉の付属的な意味や自然な展開は、「描写」という言葉よりも異なります。

まず、一般的な認知の用語と同様に、「認識」について話すときには、対象が実際に私の心に現れていないとしても、一度は私の心にあったり、機会があれば再び現れる場合でも、その対象に対して認識していると言える自然な状況があります。これは私が他のことを考えているときでも2+2=4を知っているという意味と同じです。

また、「認識」は「描写」と比較して、我々が関心を持つ事実の関係性を強調する意図を持っています。私は、描写について話すときに、対象に過度に焦点を当てるために、主語を見落とす危険があると考えています。その結果、主語がないという見方につながり、これにより物質主義に至ります。あるいは、描写されるものは主語の一部であるという見方につながり、これがイデアル主義へと進み、ソリシズムへと進むためには最も激しい扭曲が必要となります。

私は、「認識」という言葉を使用することで、主語と対象の二元性を保持したいと考えています。この二元性は私にとって認知に関する基本的な事実であると考えています。したがって、「認識」を選んだのは、その言葉が認識する主体の必要性を強調しているからです。

当われが認識する対象の種類を問うとき、最初で最も明白な例は \textit{感覚データ} である。私は色を見たり音を聞いたりするとき、直接的にその色や音と対面している。これらの状況において私が認識する感覚データは、一般的にそしておそらく常に複雑なものである。特に視覚の場合ではそれが明らかである。もちろん、私が意味するのは物理的な物体が複雑であるだけでなく、直接的に認識される感覚的対象自体が複雑で空間的な部分を持つことである。複雑なものを認識することができるのにその構成要素を認識しないことは可能なのかという問題は容易な問いではないが、全体としてそれは不可能ではない理由がないように思われる。この問題は自己意識に関連して鋭く現れ、現在私たちはそれを簡潔に考慮しなければならない。

反省において、私たちは自己を対象とする複雑な状態や関係に気づくような直感的な認識を持つように思われる。太陽を見ると、私はしばしばそのことに気づきながら太陽自体も意識する。食事の望みを持つと、そのことに気づきながら自身の望みも意識する。しかし、自己だけを対象とする意識状態を見つけ出すのは難しいし、おそらく不可能である。自己意識の性質について議論するのは大きすぎるし、私たちの主題とはあまり関連がないので、ここでは詳しくは述べない。もし私たちは自己と直接的な認識を持っていないと仮定すれば、単純な事実を説明することは難しいかもしれないが、不可能ではない。複合体「自己-認識-A」に気づき、かつ「私はAを認識している」という命題も知っているという事実は明らかである。この複合体は分析され、もし「私」が直接的な認識の対象でないものだとすると、「私」は説明によって知られている何かと考える必要がある。自己と直接的な認識を持っていないという見方を維持したいのであれば、次のように議論できる：私たちは認識そのものを認識しており、それは関係であることを知っている。また、私たちは自身の複合体に気づき、そこでは認識が関連する関係であることがわかる。したがって、この複合体には認識される対象という構成要素があり、つまり主項と述語項があるべきである。その主項を「私」と定義する。したがって、「私」は「私が意識する意識の主項」を意味する。「私」の定義としてはこれほど適切なものはないように思われるが、それでも必要であれば、自己と直接的な認識を持っていると仮定し、「私」はその特定の対象の正当な名前であると考えるか、あるいは自己意識の別の分析を見つけるべきである。したがって、自己意識は複合体を知ることなくその構成要素を知ることができるかどうかという問いに光を当てるものではないと見做すことができる。しかし、この問いは現在の目的には重要でないため、さらに議論するつもりはない。

これまで考慮した意識はすべて、特定の存在に対する意識であり、大まかには感覚データと呼ばれるべきもの全般に該当する可能性があります。知識論的な観点から見れば、内省的知識は視覚や聴覚から得られる知識と同じレベルにあるのです。しかし、上記のような対象に対する意識だけでなく（これを「個体の意識」と呼びましょう）、それとは異なる意味での「共相の意識」も存在します。共相の意識は「概念化」と呼ばれ、共相自体は「概念」と呼ばれます。特定の黄色を認識するだけでなく、十分な数の黄色を見た後で十分な知性があれば、黄色という共相を認識することもあります。この共相は、「黄色は青と異なる」「黄色は緑よりも青に似ていない」といった判断における主語となります。「これは黄色だ」において「これ」は感覚データです。また、共相の関係も意識の対象となります。上と下、前と後、類似性、欲望、そして意識そのものなど、これらすべてが認識できる対象であると考えられます。

関于関係，有人可能会主張我達从未直接意識到普遍的関係本身，而是只意識到包含該関係的複合体。例如，可以説我達並不直接知道“之前”これ種関係，尽管我達可以理解諸如“これ个在那个之前”的命題，並且可以直接意識到諸如“これ个在那个之前”これ様的複合体。然而，これ一観点難以とこれ様一个事実相調和：即我達経常知道以関係為謂語的命題，或者関係的論元不是明確給定的対象而是“任何事物”。例如，我達知道如果一个事物在另一个事物之前，而另一个事物又在第三个事物之前，則第一个事物也在第三个事物之前；在これ里涉及的事物並不是明確的事物，而是“任何事物”。除非我達熟悉“之前”これ一関係本身，而不是僅僅了解某个具体対象相対于另一給定対象処于之前的関係，否則很難理解これ様的事実。更直接地説：从複合体的意識中得出的判断如“これ个在那个之前”，構成了一種分析；如果我達不了解所用術語的意義，就不会理解これ種分析。因此，我達必須假設我達熟悉“之前”的意義，而不僅僅是其実例。

由此看来，在我達意識到的対象中有至少両類，即个别事物和普遍概念。我将所有存在物以及一个或多个構成部分為存在的複合体包括在个别事物中，如“これ个在那个之前”、“これ个在那个之上”、“これ个黄色”。而将所有没有个别事物作為其構成部分的対象包括在普遍概念中。因此，“普遍-个别”的析取涵盖了所有対象。我達也可以称之為“抽象-具体”的析取。これと“観念-感知”的対立並不完全対応，因為被記住或想象的事物属于个别事物的范畴，但很難被称為感知。（相反，我達熟悉的普遍概念可以と観念相等同。）

由此可見，在我達直接意識到的対象中不包括物理対象（相対于感覚数据而言），也不包括他人的意識。これ些事物是通過我所謂的“描述性知識”為我達所知的，我達現在必須考慮これ一点。

以下が翻訳です：

「説明」とは、「～～である」や「～～の」のような文句を指します。形式が「～～である」という文句を「曖昧な説明」と呼び、形式が「～～の」（単数形）を「確定的な説明」と呼びます。したがって、「男である」とは曖昧な説明であり、「鉄の仮面を持つ男」とは確定的な説明です。曖昧な説明に関連する問題はいくつかありますが、それらは私が議論したい事柄に直接関わらないため、それらを避けて進みます。私は議論したいのは、ある確定的な説明に対応する物体が存在することは知っているものの、その物体自体とは接していなければ、私たちのその物体に関する知識の性質についてです。これは確定的な説明に関わる唯一の問題です。したがって、以下では「説明」という言葉を用いる際は、「確定的な説明」を指すものとします。つまり、確定的な説明とは形式が「～～の」（単数形）の文句を指します。

ある物体が「説明によって知られている」とは、その物体が「～～である」という状況で、一つだけ特定の性質を持つ物体があることを我々が知っているとき、つまり、一つだけの物体が特定の性質を持っているということが分かっているときに、そう言います。そして通常、我々が同一の物体について接して認識しているとは限らないと仮定します。鉄の仮面を持つ男は存在したことは知られていますし、彼に関する多くの命題も知られていますが、彼が誰であるかは分かりません。最多得票数を得た候補者が選ばれることが分かっていますが、実際には最多得票数を得る候補者と接して認識している可能性がありますが、どの候補者がその人であるかは分かりません。つまり、「Aは最多得票数を得る候補者である」という形の命題を一つも知らないということです。「a」が我々が接して認識している何かであるとき、「～～の」形式の命題「aは～～である」というものを指します。

ある物体が存在することは分かっている場合、その物体自体とは接していないとしても、それが実際の物体である可能性がある場合でも、我々が「aは～～である」という形の命題を一つも知らないとき、「単なる説明的な知識」があると呼びます。

当われが「何かが存在する」と言うとき、それはただ一つの対象がその何かであることを意味します。「$a$は何かである」という命題は、$a$が何かに属しており、他には何も属していないことを意味します。「ジョセフ・ラマー卿は連合候補である」というのは、「ジョセフ・ラマー卿は連合候補であり、他の誰もそうではない」ということです。「連合候補が存在する」は「誰かは連合候補であり、他の誰もそうではない」ということを意味します。したがって、何かを知っている対象について我々がその何かであることを認識しているときには、その何かが存在することを知りますが、それ以外の対象を知らなくてもその何かが存在することが分かり、実際にはその何かでないどんな対象も知らない場合でも同様です。

普通の言葉すら、適切に使用される限りでは通常は実質的な説明であるものです。つまり、正規名詞を使用する人の心の中に含まれる考えを明示的に表現するためには、通常はその正規名詞を説明に置き換える必要があります。さらに、同じ言葉でも異なる人間や異なる時間では必要な説明が異なります。唯一一定なものは、名詞が適切に使用される限りにおいて適用される対象です。しかし、この対象が一定である限り、名詞を含む命題の真偽はその特定の説明によって通常は変わらないものです。

以下が翻訳です：

いくつかの例を挙げましょう。ビスマルクについて述べられた何かという仮定を考えます。自己に対する直接的な認識があるとするなら、ビスマルク自身は自分の名前を使って、彼が認識していた特定の人を指すことができたかもしれません。この場合、彼が自己について判断を行ったとしても、彼自身がその判断の構成要素となる可能性があります。ここでは、通常希望されるように直接的な用法が現れ、単に特定の対象を表すために使用されます。しかし、ビスマルクを知る人が彼について判断を行う場合、状況は異なります。この人はビスマルクと関連付けられた感覚データによって認識した身体や心を知っていたのです。つまり、彼らはこれらの感覚データに関連付けられた身体と心を通じて認識されたものであり、より一層の物理的な身体や心はその通りでした。それは、説明によって認識されたものです。もちろん、友人が彼を思い出す際にどの特徴が浮かぶかは非常に偶然によるもので、その説明は友人の頭の中に偶然存在します。重要な点は、彼自身の対象について直接的な認識をしていないにもかかわらず、様々な説明が全て同じ対象を指していることを知っていることです。

当我們、ビスマルクを知らない者が彼について判断するとき、私たちの頭の中で浮かぶ説明はおそらく、歴史的な知識の不確かな塊となるでしょう——多くの場合、彼を特定するために必要な情報よりも遥かに多いです。しかし、例として説明のためには、ビスマルクを「ドイツ帝国初代首相」と考えるとしてみましょう。ここでのすべての単語は抽象的で、「ドイツ」のみが具体的な意味を持ちます。「ドイツ」もまた、人々によって異なる意味を持つでしょう。一部の人々にとってはドイツへの旅行、別の人は地図上のドイツの外見など、様々な思い出を呼び起こします。

しかし、適用可能な説明を得るためには、私たちが知っている特定の要素を参照する必要があります。過去、現在、未来（確定的な日付とは異なる）、ここと彼方、他者が教えてくれたことなどについて言及する際には、そのような参照が含まれます。したがって、ある特定のものに適用される説明を得るためには、私たちが知っている特定の要素を参照せざるを得ないようです。つまり、私たちが描写されたものの知識が、ただその描写から導かれる論理的なものだけであるなら、適用可能な説明は私たちが知っている特定の要素を参照する必要があります。

例えば、「最も長寿の人」は誰かに適用される説明ですが、この人についての判断をするためには、その説明から得られる知識を超えた情報が必要ありません。しかし、「ドイツ帝国初代首相は機敏な外交官だった」と述べるとき、私たちが知っている特定の要素（通常は耳にした証言や読んだもの）に基づいて真実性を確認できます。

心理学的に考えると、他者への情報を伝えることなく、実際のビスマルクに関する事実が判断の重要性をとえますが、私たちは本当に持っている思考には、私たちが知っている一つまたは複数の特定要素が含まれており、それ以外は概念のみで構成されています。場所の名前——ロンドン、イギリス、ヨーロッパ、地球、太陽係——もまた、使用されるたびに、私たが知っている一つまたは複数の特定要素から始まる説明を含みます。

私は、メタ物理学で考慮される宇宙についてさえも、このような特定要素との関連性があると疑っています。論理においては、存在するものだけでなく、存在する可能性のあるものや可能なものについても考慮する必要があるため、実際の特定要素への参照が含まれません。

当たるところについて述べるとき、しばしば私たちは説明を含む形式ではなく、実際に描写されたものについて断言しようとする意図があるように思われる。つまり、ビスマルクについて何か述べるとき、もし可能であれば、彼自身が持つ判断を下すことを望む。それは、彼自身の構成要素である彼自身による判断であるからだ。しかし、実際のビスマルクは私たちは知らないため、この目標には必ずしも成功しない。しかし、私たちはBという名前のオブジェクトが存在し、それがビスマルクであり、かつ彼は機知に富んだ外交官だったことを知っている。したがって、私たちは「Bは機知に富んだ外交官だった」という命題を描写することができる。ここでBはビスマルクという名前のオブジェクトである。

異なる説明を使用してもコミュニケーションが可能である理由は、実際のビスマルクに関する真実な命題が存在し、説明が正しくいかに変化しようとも、その命題自体は同じであることを私たちは知っているからだ。この真実な命題を描写し、それが私たちにとって興味深いものだが、その命題自体にはまだなじみがないため、その命題そのものを知らないが、それは真実であることは知っている。

個別の事柄に対する認識からの移行には Various Stages がある：私たちは彼を知った人々にとってはビスマルク、歴史を通じて彼を知る人々にとってはビスマルク、鉄の仮面を持つ男、人類中最長寿の人々などがある。これらは逐次的に個別の事柄に対する認識から遠ざかっているし、普遍的な領域においても似たような階層が存在する。多くの普遍的命題と同様に、多くの個別的事柄も説明によってしか私たちは知らない。しかし、個別的事柄の場合と同じく、説明によって知られているものは最終的には認識によって知られているものに帰着される。

述語を含む命題の分析における根本的な認識論的原理は次の通りである：\textit{我々が理解できるすべての命題は、我々が認識している全ての要素から成り立っているべきである.}既に述べたことから明らかになるように、この原理を推奨する理由と、最初に見えたようにその原理に反する命題に対処する方法について説明しよう。まず、その原理が真実であると考えられる理由から始めよう。


主な理由は、判断や仮定を行うことが可能であると信じるのは、それが何についての判断や仮定であるかを知らなければ行えないという点にあると思われるからです。たとえば、ジュリウス・カエサルについての判断を行ったとしても、実際のジュリウス・カエサルそのものがその判断の要素ではないことが明白です。しかし、さらに進む前に、私が「このものやそのものは判断や理解した命題の要素である」という言葉をどのように解釈しているか説明しておくことが良いかもしれません。

まず最初に、判断について：判断は、ある心が何らかの存在に対して関係を持つ出来事としてとらえます。これは、判断されるものが構成する存在たち、すなわち、判断されるものからなる複合体の一部です。たとえば、AがBを愛しているという判断がある場合、その判断は特定の瞬間に私、A、愛、そしてBとの四項関係である「判断」という名前の関係によって存在することになります。つまり、私は判断を行うときには、自分自身、A、愛、そしてBという一定の複合体があり、その関係性が「判断」であることになります。

この見解に対する私の理由は他處で述べられていますので、ここでは繰り返しません。この見解に基づいて、判断の要素は、それが判断である複合体の要素です。したがって、上記の場合では、要素は自分自身、A、愛、そしてB、そして「判断」です。しかし自分自身と「判断」という要素はすべての私の判断で共有されるため、特定の判断の区別的な要素はA、愛、そしてBとなります。

次に、「命題を理解する」という意味について説明します。「AがBを愛している」という命題を理解するとは、私とAと愛との間に関係がある別の関係があり、それが「AがBを愛している」ということを仮定することである。AがBを愛しているという仮定ができるとき、「AがBを愛している」命題を理解できるということです。したがって、十分な知識がない場合でも、ある命題を理解することはしばしば起こります。

仮定も判断同様に多項関係であり、心はその一つの要素です。その他の要素は「仮定された命題の要素」と呼ばれます。したがって、私は述べた原理は次のように再述できます：「仮定や判断の関係が発生するとき、その関係によって仮定や判断を行う心と関連している項は、その心が認識している項である」これはつまり、何かについての判断や仮定を立てるにはそれが何であるかを知らなければならないということです。私はこの原理の真実性が理解した瞬間に明らかだと考えていますので、以下ではこの原理を前提とし、説明する判断に含む記述を分析する指針として使用します。

以下が翻訳です：

ジュリウス・シラスに戻ります。彼自身は、私が下すどんな判断にも含まれていないことは認められるでしょう。しかし、ここで議論すべき点は、判断が「アイデア」と呼ばれる何かから成り立っているという見解です。そしてそれは私の判断に含まれている「ジュリウス・シラスのアイデア」であると主張するのです。この見解が説得力を持つ理由は、「アイデア」に関する誤った理論を形成していないかにあると考えます。

私の「アイデア」としてジュリウス・シラスを意味するのは、私が知っていることの集合です。例えば、彼はガルフを征服し、3月15日に暗殺され、学校の生徒たちにとって厄介な人物であるということです。ここで私は、そして主張するところでは、ジュリウス・シラスについて判断するとき私が頭に持っている実際の内容を明らかにするために、その名前を適切な名詞から説明に置き換えるべきだと認めます。しばしば私の考えを表現するためには「名は『ジュリウス・シラス』である男」という説明が役立つでしょう。彼について私が忘れてしまった他のことに関わらず、少なくともその名前については忘れていないことが明らかです。

しかし、アイデアの理論が提案された方法で見当違いであると私は信じています。この理論によれば、ある外側のものに対する「アイデア」と呼ばれる精神的な存在があると考えられ、そして判断は精神的な出来事であるため、その成分は判断する人の心の成分であるべきだと主張します。しかし、この視点ではアイデアが私と外部の物事との間の障壁となり、我々は知識において実際のものではなく、それらのアイデアにしか到達しないと考えられます。この視点における心、アイデア、対象の関係は完全に不明瞭であり、私は観察によって発見できる何らかの理由も存在せず、アイデアが心と対象の間に入ることを正当化するものはありません。

私はこの視点が関係性への嫌悪感から生じていると考えています。そして、我々は何か「内側」にあるものを持っている必要があると感じ、それが対象を知る状態であると呼びたいと考えます。しかし、このような見解は即座に無限の退化を招きます。アイデアが対象に対するアイデアを持っていることによって説明されなければならず、そのアイデアもまた対象に対するアイデアを持っていることになります。つまり、無限に続く悪循環です。

したがって、我々は対象と知覚するときに関連性があるだけであり、心の内側に何か「アイデア」と呼ばれる成分が必要であるという主張者は誤りだと信じます。これは大きな問題であり、十分に議論すると私たちの主題から遠ざかる可能性があります。したがって、私は上記の指摘と、判断する際には対象そのものではなく、仮定される純粋な精神的存在ではなく、実際の対象が判断という複合体の成分であることを確認します。

したがって、私が「ジュリウス・セシリア」をジュリウス・セシリアの説明に置き換える必要があると言っているのは、「ジュリウス・セシリア」という概念ではなく、説明自体であることを意味します。たとえば、私たちの説明が「名前が『ジュリウス・セシリア』だった男」であるとする場合、判断が「ジュリウス・セシリアは暗殺された」とするなら、「名前が『ジュリウス・セシリア』だった男は暗殺された」となります。ここでは、「ジュリウス・セシリア」は私たちが知っている音や形であり、その判断の他の構成要素（「は」を除く tense を含む）は私たちが知っている概念です。したがって、この判断はすべて私たちは知っている構成要素に削減されますが、ジュリウス・セシリア自体はこの判断の構成要素から消えてしまいます。しかし、これは短時間後に説明する前提条件を含みます：すなわち、「名前が『ジュリウス・セシリア』だった男」全体が判断の構成要素であるべきでないということです。つまり、この文は全体として判断に入るべき意味を持たないようにすべきです。したがって、この判断の適切な分析では、この文を部分的な判断の一部として扱うのではなく、分解する必要があります。「名前が『ジュリウス・セシリア』だった男は暗殺された」という判断は、「ただ一つの男だけが『ジュリウス・セシリア』と呼ばれており、その男は暗殺された」と解釈されるかもしれません。ここでは明らかに、「名前が『ジュリウス・セシリア』だった男」に対応する構成要素はありません。したがって、この文を判断の構成要素として解釈すべき理由はありませんし、すべての判断の構成要素について知るためにはこの文を分解しなければならないことがわかりました。この結論は認識論に関する考察から達成されましたが、論理的な考慮事項も同様に私たちを導きます。これらを簡潔にレビューする必要があります。


一般的には、意味と表象という二つの側面を区別することが多い。「ワーヴリーの著者」という表現の場合、その意味は著作者とワーヴリーとの関係を持つ一定の複合体である一方で、表象はスコットである。同様に、「羽根のない二足歩行者」も複合的な意味を持ち、その構成要素には二足歩行と羽根の無さが含まれるが、その表象は人間というクラスである。したがって、「スコットはワーヴリーの著者である」とか「人間は羽根のない二足歩行者に等しい」と言ったときには、表象の同一性を主張しているのであり、この主張は意味の多様性から価値がある。私は、意味と表象の対立という虽然可以真実地解釈，但如果将其視為根本性的則是誤導的これ一観点持有异議。我認為，除了専有名詞以外，表象並不是命題的一个構成要素；専有名詞是指那些不賦予対象某種性質而僅僅命名它的詞語，例如“我”和“このもの”。

表象不是命題構成要素的理由之一是，即使我達不了解表象也能知道命題。人達在不知道「ワーヴリーの著者は小説家である」という命題表示的是スコット的情况下也知道これ个命題。これ一理由已経得到了充分的强調。

第二の理由は、「この何か」に関する命題が、「この何か」に意味合がある場合でも成立することである。例えば「金色の山は存在しない」とか「円形の四角形は自相矛盾する」というような例を挙げることができる。もし意味と対象との二重性を保存するのであれば、ミンローと同じく、「この何か」のような対象が存在すると主張せざるを得ない。これらの対象は実在しないとしてもである。すなわち、存在する円形の四角形は存在しているが、実在しないという。ミンローはこれを見做さないが、私はそれが矛盾ではないとは見ることができない。実際、私には明らかに思われるが、「円形の四角形のような対象は存在しない」という判断は、そのような対象が存在することを前提していない。これが認められれば、論理形式による類推により、すべての「この何か」に関する命題が実際にその「この何か」を構成要素として含むわけではないという結論に至る。

ミンズ・ジョーンズは、「現在のフランス王」という対象について矛盾する性質をもつ予言を認めることは困難がないと主張しています。彼女は、この対象が自己矛盾しているわけではないが、単に存在しないだけであるという理由で、円形の正方形とは異なりその対象自体が矛盾しているわけではないと主張することができます。しかし、これは問題の本質には触れていません。このような主張に対する実際の反論は、「排中律」を伝統的な形式「AはBでありかつBではない」というものではなく、「命題は真であるか偽であるかのいずれでもない」という形式で述べるべきであるということです。伝統的な形式は、予言が主題と述語に対する属性をとえるような予言に対してのみ適用されます。命題についての法則を主題や述語について述べる代わりに述べると、すぐに明らかになるように、「現在のフランス王」に関する予言や「円形の正方形」に関する予言は例外ではなく、他の予言と同じく真であるか偽であるかのいずれでもないという能力が欠如しています。ミンズ・ジョーンズは、「スコットは『ウェーバリー』の著者である」という主張は、「スコット」や「『ウェーバリー』の著者」という用語の表意性間の同一性を主張していると論じています。しかし、この主張の異なる意味に対する選択にはいくつかの困難があります。まず第一に、注意すべきは、「『ウェーバリー』の著者」は「スコット」のような単なる名前ではないということです。「スコット」は特定の人を指すために慣習的に使用される音や形状であり、その人に関する情報をとえず、意味と呼ばれるものとは区別された表意性を持っていません。（上記で考慮したように、通常、名詞も説明の代りに立っていることが多いという事実を無視します。）しかし、「『ウェーバリー』の著者」は単なる「スコット」の名前であるだけでなく、単語「the」「author」「of」「Waverley」がそれぞれ何を指すかによって、その要素が単なる慣習に属しています。これらの単語がどのような意味を持つかをとえられた場合、「『ウェーバリー』の著者」という用語は任意ではなくなります。スコットは「ウェーバリー」の著者であるという主張を述べるとき、我々は「スコットはシーウォルターである」と言った場合とは異なり、これらが一人の人間の二つの名前であると述べていないということに注意しなければなりません。人間の名前はその人が呼ばれるものであり、スコットが「ウェーバリー」の著者として呼ばれることは多量であっても、それが彼を著者にするわけではありません。彼が実際に『ウェーバリー』を書くことが必要であり、これは名前の関係とは無関係な事実でした。

もし名詞の指称が同一であると主張するのであれば、\textit{指称}とは単に名詞が命名されるものの関係を意味してはなりません。実際には、「スコット」の意味は「ワーヴリーの作者」という言葉の指称に近いと言えるでしょう。「スコット」に対する関係は、「スコット」がスコットを意味することと同じで、「著者」に対する関係は、「著者」がその概念を意味することと同じです。したがって、「ワーヴリーの作者」という言葉で区別する場合、「スコット」には意味があるものの指称はありません。また、「スコットはワーヴリーの作者である」と言うとき、「ワーヴリーの作者」の意味は主張に関連しています。もし指称だけが関連するとすれば、同じ指称を持つ他の表現でも同様の命題になります。したがって、「スコットはマリオンの作者である」という命題は「スコットはワーヴリーの作者である」という命題と同じです。しかし、これは明らかにそうではなく、最初の命題からスコットがマリオンを書いたことを学び、二番目の命題からスコットがワーヴリーを書いたことを学ぶことになりますが、最初のものはワーヴリーについて何も教えてくれず、二番目のものはマリオンについて何も教えてくれません。したがって、「ワーヴリーの作者」という言葉の意味は「スコットはワーヴリーの作者である」という主張に関連していることは明らかです。

Thus 我々は「ウェーバリーの著者」が単なる名前ではなく、その意味が現れる命題において関連があることを合意した。
すなわちミス・ジョーンズが言うように、「スコットはウェーバリーの著者である」という主張が指称の同一性を主張している場合、我々は「ウェーバリーの著者」の指称が「ウェーバリーの著者」という言葉で意味されるものを指す指称であると見なさなければならない。
そこで「ウェーバリーの著者」という言葉の意味をMと呼ぶことにしよう。つまりMは「ウェーバリーの著者」の意味である。
すると我々は「スコットはウェーバリーの著者である」という主張が「スコットはMの指称である」という意味であるとするべきだ。
しかしここで我々は同じ形式の命題で自己説明をしているだけであり、実質的な解釈に進展していない。
「Mの指称」もまた「ウェーバリーの著者」同様に意味と指称を持つことになる。理論を検討している場合、その意味をM'と呼ぶことにしよう。すると我々の主張は「スコットはM'の指称である」となり、これにより無限の帰納的解釈が生じる。
したがって、我々の命題を指称の同一性を主張するものとして見なす試みは失敗し、別の分析を見つけることが必要となる。この分析が完了したら、「指称の同一性」という表現を再解釈できるようになるだろう。これは基本として取り扱われる限り不明瞭であり続けるからだ。


まず注意すべき点は、"ワーヴリーの作者"に関する任意の命題において、スコットが明示的に言及されていない限り、その命題における表象自体すなわちスコットは登場せず、代わりに表象概念が変数で表現されます。例えば「ワーヴリーの作者はマリオンの作者だった」と述べた場合、我々は必ずしも両者がスコットであることを意味しているわけではありません——我々はスコットという人物を忘れているかもしれません。我々は何か一人の男がワーヴリーとマリオンの作者であったということを言っているに過ぎません。つまり、ある人がワーヴリーとマリオンを書いたし、他には誰もそれらを書かなかったということです。したがって、その同一性は変数の同一性、すなわち不確定な主語（「某人」）の同一性に相当します。これが我々がスコットの名前を知らずとも"ワーヴリーの作者"に関する命題を理解できる理由です。例えば、「ワーヴリーの作者は詩人だった」と述べた場合、それは「一人だけがワーヴリーを書いたし、彼は詩人だった」を意味します。「ワーヴリーの作者はスコットだった」と述べた場合、それは「一人だけがワーヴリーを書いたし、彼はスコットだった」という意味です。ここでは、同一性は変数、すなわち不確定な主語（「彼」）とスコットとの間にあることになります。「ワーヴリーの作者」は分析され、命題の構成要素として現れなくなりました。


分析する必要がある「Waverleyの作者」という表現について、その理由は次のように述べることができます。つまり、「Waverleyの作者はMarmionの作者である」と言うとき、「は」が同一性を表していることは明らかです。また、「共通の意味域」すなわちスコットはこの命題には含まれていないことも確認しました。さらに、「Waverleyの作者」と「Marmionの作者」の意味（もし存在する場合）は異なることがわかりました。

言葉が命題の文言に含まれる任意の意味において、その意味が命題の構成要素であるとすると、「スコット」は実際のスコットという人を指し、「作者」は特定の普遍的概念を指す同一の意味で用いられます。したがって、上記の命題における「Waverleyの作者」という表現が下部複合語として存在する場合、その「意味」は「Marmionの作者」の「意味」に等しいものであるべきです。これは明らかに成り立たないため、唯一の解決策は「Waverleyの作者」という表現自体には意味がないが、その部分を構成する表現には意味があるということです。

つまり、上記の命題の適切な分析では、「Waverleyの作者」が消える必要があります。これは上の命題を「誰かがWaverleyを書いたし、それ以外の人間は書かなかったし、その同じ人がMarmionも書いたし、それ以外の人間は書かなかった」という意味に分解することで実現されます。これにより、「x」がWaverleyとMarmionを書いたし、それ以外の人は書かなかったという命題関数が真であることが可能となり、ある値の「x」が真をもたらす一方で他の値はそうならないことがわかります。

したがって、当たる判断の真正な主語は命題関数であり、未確定の構成要素を含む複合体で、その構成要素が決定されると命題となります。

我達現在定義述語の表意。もし「$a$ はこのようにと呼ばれるものである」という命題が真であることを知っている場合、すなわち $a$ はこのようにと呼ばれるものであり他のものは何もないと分かっている場合、$a$ を「このようにと呼ばれるもの」の述語の表意と呼びます。我々が自然に「このようにと呼ばれるもの」について述べる多くの命題において、$a$ に「このようにと呼ばれるもの」と置き換えると、その真偽は変わらないか、または同じです。また、「このようにと呼ばれるもの」を他の表意と同じである述語で置き換えると、その真偽も変わりません。したがって、実用的な観点から見れば、我々は表意よりも説明に興味を持ちます。なぜなら、説明の中に現れる多くの命題の真偽を決定するのは表意だからです。さらに、以前述べた説明と認識の関係においても、しばしば表意に到達したいという欲求があり、それが阻害されるのは認識不足によるものです。そのような場合、説明は表意に近づくための手段に過ぎません。したがって、表意が命題の中で部分的な構成要素であると誤解する自然な傾向がありますが、論理的根拠や認識論的根拠から見ればこれは誤りです。実際の対象（それが存在すれば）は、説明の中に現れる命題において明示的に述べられなければ、その表意は命題の構成要素ではありません。そのため、このような命題を理解するには説明の構成要素に対する認識が必要ですが、表意については必要ありません。分析を適用した最初の結果として、「the so-and-so」が文法的な主語である命題では、変数を主語とします。すなわち、「何かだけがこのようにと呼ばれるものであり、その何かはこのようなものである」という形の命題を得ます。さらに「the so-and-so」についての命題の分析は、変数の性質、つまり $some$, $any$, および $all$ の意味に関する問題に合流します。これは難しい問題で、現在そのことについて何も述べません。

概要を述べると、私たちは主観的知識と説明的知識という二種類の対象に関する知識を区別しました。これらの中で、主観的知識だけが対象そのものを心に引き寄せます。私たちは感覚データや多くの共相、そしておそらく自分自身について主観的知識を持ちますが、物理的な物体や他の意識については持っていないのです。ある性質やいくつかの性質を持つ何かを知るとき、それが私たちが認識している性質やいくつかの性質を持つ一つの対象であると知っている場合、私たちは説明的知識を持っていると言います。つまり、その性質やいくつかの性質は一つの対象に属し、それ以上にはないとき、私たちはその一つの対象について説明的に認識していると言えます。物理的な物体や他の意識に関する私たちの知識はすべて説明的知識であり、これらの説明には感覚データが含まれることが多いです。私たちが理解できる全ての命題は、主観的知識を持っているもので構成されています。主観的知識を持っていない要素は私たちは理解できません。判断は「アイデア」と呼ばれる心的な要素から成り立つわけではなく、ある心と特定の対象、個体や共相が含まれる出来事から成り立っています（少なくとも一つは共相である）。正しく分析された判断の構成要素となるすべての対象は、その判断の構成要素である心が認識している対象です。この結論は、命題の中で現れる説明的表現を分析することを強制し、これらの表現で示される対象がそのような表現が含まれる判断の構成要素ではない（これらの対象が明示的に述べられている場合を除く）と言わざるを得ません。これはまた、当該の判断は同一性の表出と意味の多様性との間の主張ではなく説明されるべきであるという見解（論理的な理由でも推奨されている）を導きます。つまり、スコット自身が私たちの判断の構成要素ではないし、その判断を同一性の表出と意味の多様性との間の主張によって説明することはできないのです。このような判断は説明的表現を分解し、変数を導入し、命題関数を最終的な主題とする分析が必要です。実際、「〜である」というような判断は「xが〜であり、それ以外は何でもないし、xは〜である」が真になる可能性があるということを意味します。このような判断の分析には多くの新しい問題が生じますが、これらの問題についての議論は現在の論文では行われません。


脚注：

 過後参照。

 \textit{Philosophical Essays}「真理の性質」。ウィトゲンシュタイン氏により説得されたが、この理論はやや単純すぎると思われ、しかし要求される修正は上記の議論[1917年]に影響をとえない。

 Meinong参照、\textit{Ueber Annahmen}。以前はミンオング氏の見解とは異なり、仮定する関係は単なる提示関係であると考えていたが、現在ではその見解は誤りであり、ミンオング氏は正しいと認識している。しかし現在の私の見解は、判断や仮定において客観性が一つではなく、判断または仮定の各成分は心に対する多項関係にあるという理論に基づいている。

 この見解は最近ジェーンズ博士E.E.C.によって主張されている。「新しい論理法則とその意味」、$Mind$, 1911年1月号。

 私は「私」という言葉を厳密な意味での名詞から除外し、「これ」という用語だけを残すべきであると考える[1917年]。

 Meinong, \textit{Ueber Annahmen}, 2版、ライプツィグ、1910年、第141ページ。

 $Mind$, 1910年7月号、第380ページ。

 $Mind$, 1910年7月号、第379ページ。

 私が主張する理論は、\textit{Principia Mathematica}、巻I、導入、第III章で完全に述べられており、その論理的根拠も示されている。また、$Mind$, 1905年10月号では簡略化された形で述べられている。

 この用語は、判断に含まれる心理学的な何かを指すために使用し、それが何かであるという問題に対する偏見を抱くつもりはないことを意図していない。

索引

アッキレスと亀、80番以降、89番以降

認識の関係、209番以降

アレクサンドル、125番

アメリカ実在論者たち、134番

アリストテレス、42番、76番、97番


 Bacon、41番

ベルゲンソン、14番以降、22番、105番、128番、185番以降、203番

ベーコン、97番、132番

ブレイク、1番

ボサンケット、99番

ブラッド、89番注


微分積分学、82番


カントロジー、64番、81番以降、85番、91番

キャリー、50番、82番

原因の概念、135番注、180番以降

キリスト教と断念、51番

 Chuang Tz[)u]、106番

恒久的なものや物質の構築、169番以降

論理的構築、155番以降


ダーウィン、15番、23番、43番

デデキント、64番、81番以降、85番

 Descartes、97番、126番

説明、175番、214番以降


教育、37番以降

ユークリッド、62番、92番、94番

進化論、23番以降、28番


ファノ、93番

ファラデイ、34番

自由意志、205番以降

フレーゲ、78番注


ガリレオ、42番

グラスゴー卿、177番注

善と悪、26番以降


ヘーゲル、8番、10番、18番、85番、97番、105番以降

ハイネ、113番

ヘラクレイト、1番以降、10番

ヘルツ、34番

ホルト、177番注

 ヒューム、1番、97番


無限の数学的なもの、84番以降


ジェイムズ、ウィリアム、100番

ジェーンズ博士E.E.C.、224番注、225番

判断、219番以降


カント、85番、96番、97番、99番、118番以降

知識の直接的把握、209頁から；
  説明による知識、214頁から


ラプラス、23頁


レーニッツ、76頁、79頁、82頁から、97頁、126頁、144頁、160頁


ロック、97頁


論理、その法則、68頁から


マカーソルとテイラーの定理、95頁


マラス、43頁


数学、58頁から；
  メタフィジクスとの関係、74頁から；
  論理との関係、75頁から；
  無限小との関係、82頁から


物質の性質、125頁から；
  物質の定義、164頁から


マクスウェル、34頁


意味と指称、223頁から


マイエンロング、174頁、220番脚注、225頁


 militarism、50頁


ミールズ、185頁、193頁から

神秘主義と論理、1頁から


必然性の概念、207頁から

ニーチェ、22頁、50頁

ヌーン、125頁、137番脚注、153頁


パルメニデス、7頁から、18頁、21頁

個別的事物の認識、210頁から

ペアノ、78頁から、93頁から

視点、139頁から；
  視点の空間、158頁から

哲学と論理、111頁

物理学、感覚データと、145頁から

ピアース、76番脚注


プラトン、1頁から、10頁、30頁、60頁、97頁

実用主義、22頁、105頁


現実主義と分析的手法、120頁から

理由と直感、12頁から

関係の論理、76頁

ロブ、167番脚注


サンテヤーナ、20頁

感覚データ、147頁、210頁から；
  物理学との関係、145頁から

感覚的存在体、148頁から

空間、138頁から；
  私的空間、158頁から；
  論理的な問題、114頁から；
  物理学の問題、115頁から；
  知識論的な問題、118頁から

決定論的なシステム、199頁；
  実質的に隔離されたもの、198頁；
  相対的に隔離されたもの、197頁；
  機械的もの、201頁


時間、10頁、21頁から、141頁から、167頁から

エリトリアのゼノンのパラドックス、90頁から


一元性と多元性、18頁から

普遍的存在の認識、212頁から


ウォード、180頁


ヴェイエルシュトラス、80頁、82頁、95頁

ホイットヘッド、117頁、157頁、175頁

ウルフ、173頁


エリトリアのゼノン、64頁、80頁、84頁、89頁から


[図版]

ジョージ アレン \& ウィン Ltd

ロンドン：ミュージアム通り40号、ワーグ1

オークランド：ウィンダム通り24号

ボーディー：グラハム通り15号、バールドレジデンス、ボーディー1

カルカッタ：チットラランジャン通り17号、カルカッタ13

ケープタウン：ロン通り109号

カラチ：メサーソンズ・エSTATE、ウッド通り2号、カラチ2

ニューデリー：アージムリゲート拡張13-14号、ニューデリー1

サンパウロ：ジュール9日通り1138番地－51室

シンガポール、東南アジアおよび遠東：プリンセップ通り36c

シドニー、ニュー・サウス・ウェールズ：ブラドバリーハウス、ヨーク通り55号

トロント：ウェストランド通り91号

本文中の印刷ミスは次の通り訂正しました：

111ページ：charateristicをcharacteristicに
71ページ：inexpugnableをinexpungableに



\backmatter

\chapter*{解題：ラッセルの論理原子論と現代科学・数学の基礎}
\addcontentsline{toc}{chapter}{解題：ラッセルの論理原子論と現代科学・数学の基礎}

バートランド・ラッセルの哲学体係の核心は、「不完全で主観的な人間の知覚や言語の曖昧さを、厳密な論理分析と数学的構成によって脱神話化すること」にあります。

ホワイトヘッドとの記念碑的共著『プリンキピア・マテマティカ（数学原理）』において、算術と解析の全体係を純粋論理学から演繹しようとしたラッセルは、本書『神秘主義と論理』に収められた諸論考において、その数理的厳密さを現実世界の実在性や物理学の認識論へと果敢に拡張しました。

とりわけ「センス・データと物理学の関係」における「推論された実体ではなく、論理的構成物としての物理的客体」という公理的定式化（オッカムの剃刀の論理学的極限）は、現代の場の量子論、情報幾何学、認知科学における実在論論争の先駆をなすものです。また、「記述の理論（Theory of Descriptions）」は、存在しない対象（例：「現在のフランス国王」）をめぐる数千年来の形而上学の混乱を、命題関数の量化記号によって一挙に氷解させました。

ラッセルが示した「感情や願望による甘美な慰めを退け、事実と論理の冷厳な光を直視する勇気」は、現代を生きる知的生活者にとっても、永遠に色褪せない理性の羅針盤であり続けることでしょう。

\end{document}
